\documentclass[10pt]{article}
\usepackage[T1,T2A]{fontenc}
\usepackage[utf8]{inputenc}
\usepackage[english,ukrainian]{babel}
\usepackage{geometry}
\usepackage{longtable}
\usepackage{array}
\usepackage{multirow}
\usepackage{hyperref}
\usepackage{mathptmx}
\renewcommand{\familydefault}{\rmdefault}

\geometry{a4paper,left=2cm,right=2cm,top=2cm,bottom=2cm}

\begin{document}

\begin{flushright}
ЗАТВЕРДЖЕНО\\
Наказом Державна судова адміністрація України\\
від \_\_ \_\_\_\_\_\_\_\_\_ 2026 р. № \_\_\_\\
(в редакції наказу Державна судова адміністрація України\\
від \_\_ \_\_\_\_\_\_\_\_\_ 2026 р. № \_\_)
\end{flushright}

\vspace{2cm}

\begin{center}
\textbf{\Large Єдина судова інформаційно-телекомунікаційна система (ЄСІКС)}\\
\textbf{\Large «Спеціалізовані Суди»}\\
\vspace{1cm}
\textbf{\Large ЦІЛЬОВИЙ ПРОФІЛЬ БЕЗПЕКИ}\\
\vspace{0.5cm}
\large UA.46207985.СЗІ.ЦПБ-41
\end{center}

\newpage

\footnotesize
\begin{longtable}{|>{\centering\arraybackslash}p{0.5cm}|>{\raggedright\arraybackslash}p{4cm}|>{\raggedright\arraybackslash}p{5cm}|>{\raggedright\arraybackslash}p{2cm}|>{\raggedright\arraybackslash}p{3cm}|}
\hline
\textbf{№} & \textbf{Вимога з безпеки інформації} & \textbf{Вимога ГПБ} & \multicolumn{2}{c|}{\textbf{ЦПБ}} \\
\cline{4-5}
& & & \textbf{Захід захисту} & \textbf{Налаштований зміст заходу захисту}\\
\hline
\endfirsthead

\hline
\textbf{№} & \textbf{Вимога з безпеки інформації} & \textbf{Вимога ГПБ} & \multicolumn{2}{c|}{\textbf{ЦПБ}} \\
\cline{4-5}
& & & \textbf{Захід захисту} & \textbf{Налаштований зміст заходу захисту}\\
\hline
\endhead

\hline
\endfoot

\hline
\endlastfoot

\multicolumn{5}{|l|}{\textbf{1. УПРАВЛІННЯ ДОСТУПОМ (AC)}} \\ \hline
1 & ПОЛІТИКА ТА ПРОЦЕДУРИ УПРАВЛІННЯ ДОСТУПОМ (AC-1) & a. Розробити, задокументувати та поширити [Призначення: серед визначеного організацією персоналу або ролей]:\newline 1. 2. [Вибір (один або декілька): Рівень організації; Рівень місії/бізнес-процесу; рівень системи] політики контролю доступу, яка:\newline (a) містить мету, сферу застосування, ролі, відповідальність, зобов'язання керівництва, координацію між організаційними підрозділами та систему контролю відповідності (compliances);\newline (b) відповідає чинному законодавству, нормативним документам, директивам, нормам, політикам, стандартам і керівним документам. Процедури, що сприяють реалізації політики управління доступом і відповідних заходів управління доступом.\newline b. Призначити на посаду [Призначення: визначену організацією посадову особу] для управління, документування і розповсюдження політики та процедур контролю доступом.\newline c. Переглянути та оновити:\newline 1. поточну політику управління доступом [Призначення: з визначеною організацією частотою] та [Призначення: події, визначені організацією];\newline 2. поточні процедури управління доступом [Призначення: з визначеною організацією частотою] та [Завдання: події, визначені організацією]. & AC-1 & За політику та процедури управління доступом в системі ЄСІКС відповідає підсистема IAM/PKI (Identity and Access Management / Public Key Infrastructure), яка керує обліковими записами, ролями та доступом користувачів згідно з вимогами та політиками ЄСІКС.\vspace{1mm}\newline\noindent\rule{\linewidth}{0.4pt}\vspace{1mm}
Параметр: Визначено персонал або ролі, на які поширюється політика контролю доступу\newline Тип: list\newline Значення: \textbf{admin, security\_\allowbreak{}officer}\vspace{1mm}\newline\noindent\rule{\linewidth}{0.4pt}\vspace{1mm}
Параметр: Визначено персонал або ролі, на які поширюються процедури контролю доступу\newline Тип: list\newline Значення: \textbf{admin, security\_\allowbreak{}officer}\vspace{1mm}\newline\noindent\rule{\linewidth}{0.4pt}\vspace{1mm}
Параметр: Визначено посадову особу, яка керуватиме політикою та процедурами контролю доступу\newline Тип: list\newline Значення: \textbf{admin, security\_\allowbreak{}officer}\vspace{1mm}\newline\noindent\rule{\linewidth}{0.4pt}\vspace{1mm}
Параметр: Визначено частоту, з якою переглядається та оновлюється поточна політика контролю доступу\newline Тип: list\newline Значення: \textbf{default\_\allowbreak{}deny\_\allowbreak{}rule, abac\_\allowbreak{}rule\_\allowbreak{}1} \\ 
 &  &  &  & Параметр: Визначено події, які вимагають перегляду та оновлення поточної політики контролю доступу\newline Тип: list\newline Значення: \textbf{default\_\allowbreak{}deny\_\allowbreak{}rule, abac\_\allowbreak{}rule\_\allowbreak{}1}\vspace{1mm}\newline\noindent\rule{\linewidth}{0.4pt}\vspace{1mm}
Параметр: Визначено частоту, з якою переглядаються та оновлюються поточні процедури контролю доступу\newline Тип: integer\newline Значення: \textbf{30} \\ \hline
2 & УПРАВЛІННЯ ОБЛІКОВИМИ ЗАПИСАМИ (AC-2) & a. Визначити та задокументувати типи облікових записів системи, дозволених для використання в ІС для підтримки цілей, завдань, функцій і процесів організації.\newline b. Призначити менеджерів облікових записів для управління системними обліковими записами.\newline c. Створити умови для групового та рольового членства.\newline d. Визначити авторизованих користувачів інформаційної системи, членство в групі та ролі, а також дозволи доступу (наприклад, привілеї) та інші атрибути (за потреби) для кожного облікового запису.\newline e. Вимагати схвалення [Призначення: визначеною організацією відповідальною особою або роллю] запитів на створення облікових записів системи.\newline f. Створювати, активувати, змінювати, деактивувати та видаляти системні облікові записи відповідно до [Призначення: визначених організацією політики, процедур та умов].\newline g. Впровадити моніторинг використання облікових записів системи.\newline h. Повідомляти адміністраторів облікових записів у межах [Призначення: визначеного організацією часового періоду для кожної ситуації]:\newline 1. коли облікові записи більше не потрібні;\newline 2. коли користувачі звільнені чи переведені;\newline 3. коли використовуються індивідуальні системи або наявні зміни, які потребують нових знань.\newline i. Авторизувати доступ до системи на основі:\newline 1. Дійсної авторизації доступу.\newline 2. Передбачуваного використання системи.\newline 3. Інших атрибутів, що вимагаються організацією.\newline j. Проводити перегляд облікових записів на відповідність вимогам управління обліковими записами з [Призначення: визначеною організацією частотою].\newline k. Впровадити процес повторного випуску облікових даних спільного/групового облікового запису (якщо він буде розгорнутий), коли особи виходять з групи.\newline l. Узгодити процеси управління обліковими записами з процесами звільнення та переводу (передачі повноважень) персоналу. & AC-2 & Політики управління доступом, привілеями та віддаленим доступом реалізуються централізовано через підсистеми IAM (Identity and Access Management) та ABAC (Attribute-Based Access Control) ЄСІКС з дотриманням принципу найменших привілеїв.\vspace{1mm}\newline\noindent\rule{\linewidth}{0.4pt}\vspace{1mm}
Параметр: account\_\allowbreak{}management\newline Тип: boolean\newline Значення: \textbf{true} \\ \hline
3 & УПРАВЛІННЯ ОБЛІКОВИМИ ЗАПИСАМИ - ДЕАКТИВАЦІЯ ОБЛІКОВИХ ЗАПИСІВ (AC-2(3)) &  & AC-2(3) & Політики управління доступом, привілеями та віддаленим доступом реалізуються централізовано через підсистеми IAM (Identity and Access Management) та ABAC (Attribute-Based Access Control) ЄСІКС з дотриманням принципу найменших привілеїв.\vspace{1mm}\newline\noindent\rule{\linewidth}{0.4pt}\vspace{1mm}
Параметр: Облікові записи деактивуються протягом часового періоду, коли облікові записи більше не пов'язані з користувачем або фізичною особою\newline Тип: list\newline Значення: \textbf{admin, security\_\allowbreak{}officer}\vspace{1mm}\newline\noindent\rule{\linewidth}{0.4pt}\vspace{1mm}
Параметр: Облікові записи відключено протягом часового періоду, коли облікові записи порушують політику організації; облікові записи деактивуються протягом , якщо вони були неактивними протягом \newline Тип: list\newline Значення: \textbf{default\_\allowbreak{}deny\_\allowbreak{}rule, abac\_\allowbreak{}rule\_\allowbreak{}1}\vspace{1mm}\newline\noindent\rule{\linewidth}{0.4pt}\vspace{1mm}
Параметр: Визначено період часу, протягом якого необхідно деактивувати облікові записи\newline Тип: integer\newline Значення: \textbf{30}\vspace{1mm}\newline\noindent\rule{\linewidth}{0.4pt}\vspace{1mm}
Параметр: Визначено період часу неактивності, після закінчення якого облікові записи будуть деактивовані; AC-02(03)(a) облікові записи деактивуються протягом часового періоду, коли термін дії облікових записів минув\newline Тип: list\newline Значення: \textbf{login, logout, failed\_\allowbreak{}attempt} \\ \hline
4 & УПРАВЛІННЯ ОБЛІКОВИМИ ЗАПИСАМИ - ВИХІД ІЗ СИСТЕМИ ЗА ВІДСУТНОСТІ АКТИВНОСТІ (AC-2(5)) &  & AC-2(5) & Політики управління доступом, привілеями та віддаленим доступом реалізуються централізовано через підсистеми IAM (Identity and Access Management) та ABAC (Attribute-Based Access Control) ЄСІКС з дотриманням принципу найменших привілеїв.\vspace{1mm}\newline\noindent\rule{\linewidth}{0.4pt}\vspace{1mm}
Параметр: Користувачі повинні виходити з системи, коли період очікуваної бездіяльності або опис часу, коли потрібно вийти з системи\newline Тип: integer\newline Значення: \textbf{30}\vspace{1mm}\newline\noindent\rule{\linewidth}{0.4pt}\vspace{1mm}
Параметр: Визначено часовий період очікуваної бездіяльності або опис, коли потрібно вийти з системи\newline Тип: integer\newline Значення: \textbf{30} \\ \hline
5 & УПРАВЛІННЯ ОБЛІКОВИМИ ЗАПИСАМИ - ДЕАКТИВАЦІЯ ОБЛІКОВИХ ЗАПИСІВ ОСІБ З ВИСОКИМ РІВНЕМ РИЗИКУ (AC-2(13)) &  & AC-2(13) & Політики управління доступом, привілеями та віддаленим доступом реалізуються централізовано через підсистеми IAM (Identity and Access Management) та ABAC (Attribute-Based Access Control) ЄСІКС з дотриманням принципу найменших привілеїв.\vspace{1mm}\newline\noindent\rule{\linewidth}{0.4pt}\vspace{1mm}
Параметр: Облікові записи користувачів деактивуються протягом періоду часу з моменту виявлення значних ризиків\newline Тип: integer\newline Значення: \textbf{30}\vspace{1mm}\newline\noindent\rule{\linewidth}{0.4pt}\vspace{1mm}
Параметр: Визначено період часу, протягом якого необхідно деактивувати облікові записи фізичних осіб, які становлять значний ризик\newline Тип: integer\newline Значення: \textbf{30} \\ \hline
6 & ЗАБЕЗПЕЧЕННЯ ДОСТУПУ (AC-3) & Застосовувати затверджені повноваження для логічного доступу до інформації та ресурсів системи відповідно до чинної політики (правил) управління доступом. & AC-3 & Усі запити до ресурсів системи примусово проходять через Policy Decision Point (PDP) модуля `abac`. Прямий доступ до даних в обхід механізмів авторизації заборонено архітектурно.\vspace{1mm}\newline\noindent\rule{\linewidth}{0.4pt}\vspace{1mm}
Параметр: access\_\allowbreak{}enforcement\newline Тип: boolean\newline Значення: \textbf{true} \\ \hline
7 & УПРАВЛІННЯ ІНФОРМАЦІЙНИМИ ПОТОКАМИ (AC-4) & Застосувати затверджені повноваження для управління потоком інформації всередині системи та між пов'язаними системами на основі [Призначення: визначеними організацією політиками управління інформаційним потоком]. & AC-4 & Політики управління доступом, привілеями та віддаленим доступом реалізуються централізовано через підсистеми IAM (Identity and Access Management) та ABAC (Attribute-Based Access Control) ЄСІКС з дотриманням принципу найменших привілеїв.\vspace{1mm}\newline\noindent\rule{\linewidth}{0.4pt}\vspace{1mm}
Параметр: Затверджені повноваження застосовуються для контролю потоку інформації всередині системи та між підключеними системами на основі політики управління інформаційними потоками\newline Тип: list\newline Значення: \textbf{default\_\allowbreak{}deny\_\allowbreak{}rule, abac\_\allowbreak{}rule\_\allowbreak{}1}\vspace{1mm}\newline\noindent\rule{\linewidth}{0.4pt}\vspace{1mm}
Параметр: Визначено політики управління інформаційними потоками всередині системи та між підключеними системами\newline Тип: list\newline Значення: \textbf{default\_\allowbreak{}deny\_\allowbreak{}rule, abac\_\allowbreak{}rule\_\allowbreak{}1} \\ \hline
8 & РОЗМЕЖУВАННЯ ОБОВ'ЯЗКІВ (AC-5) & a. Розмежувати і документувати [Призначення: визначені організацією обов'язки окремих осіб].\newline b. Установити правила авторизації доступу для підтримки розмежування обов'язків. & AC-5 & Політики управління доступом, привілеями та віддаленим доступом реалізуються централізовано через підсистеми IAM (Identity and Access Management) та ABAC (Attribute-Based Access Control) ЄСІКС з дотриманням принципу найменших привілеїв. \\ \hline
9 & МІНІМІЗАЦІЯ ПОВНОВАЖЕНЬ (AC-6) & Впровадити принцип мінімізації повноважень, який дозволяє користувачам (або процесам, що діють від імені користувачів) здійснювати лише такі авторизовані звернення, які необхідні для виконання визначених завдань відповідно до цілей (призначення, місії) організації та функцій. & AC-6 & Профілі безпеки і ролі конфігуруються таким чином, що користувачам надаються лише ті привілеї, які необхідні для виконання їхніх посадових обов'язків. Механізм дозволів `abac` працює за принципом 'Default-Deny'.\vspace{1mm}\newline\noindent\rule{\linewidth}{0.4pt}\vspace{1mm}
Параметр: least\_\allowbreak{}privilege\newline Тип: boolean\newline Значення: \textbf{true} \\ \hline
10 & МІНІМІЗАЦІЯ ПОВНОВАЖЕНЬ - АВТОРИЗОВАНИЙ ДОСТУП ДО ФУНКЦІЙ БЕЗПЕКИ (AC-6(1)) &  & AC-6(1) & Політики управління доступом, привілеями та віддаленим доступом реалізуються централізовано через підсистеми IAM (Identity and Access Management) та ABAC (Attribute-Based Access Control) ЄСІКС з дотриманням принципу найменших привілеїв.\vspace{1mm}\newline\noindent\rule{\linewidth}{0.4pt}\vspace{1mm}
Параметр: Визначені особи або ролі з авторизованим доступом до функцій безпеки та інформації, що має відношення до безпеки\newline Тип: list\newline Значення: \textbf{admin, security\_\allowbreak{}officer} \\ \hline
11 & МІНІМІЗАЦІЯ ПОВНОВАЖЕНЬ - НЕПРИВІЛЕЙОВАНИЙ ДОСТУП ДО НЕЗАХИЩЕНИХ ФУНКЦІЙ (AC-6(2)) &  & AC-6(2) & Політики управління доступом, привілеями та віддаленим доступом реалізуються централізовано через підсистеми IAM (Identity and Access Management) та ABAC (Attribute-Based Access Control) ЄСІКС з дотриманням принципу найменших привілеїв.\vspace{1mm}\newline\noindent\rule{\linewidth}{0.4pt}\vspace{1mm}
Параметр: Користувачі облікових записів (або ролей) системи з доступом до функцій безпеки або інформації, що стосується безпеки, повинні використовувати непривілейовані облікові записи або ролі під час доступу до незахищених функцій\newline Тип: list\newline Значення: \textbf{admin, security\_\allowbreak{}officer} \\ \hline
12 & МІНІМІЗАЦІЯ ПОВНОВАЖЕНЬ - ПРИВІЛЕЙОВАНІ ОБЛІКОВІ ЗАПИСИ (AC-6(5)) &  & AC-6(5) & Політики управління доступом, привілеями та віддаленим доступом реалізуються централізовано через підсистеми IAM (Identity and Access Management) та ABAC (Attribute-Based Access Control) ЄСІКС з дотриманням принципу найменших привілеїв.\vspace{1mm}\newline\noindent\rule{\linewidth}{0.4pt}\vspace{1mm}
Параметр: Привілейовані облікові записи в системі обмежено персонал або ролі\newline Тип: list\newline Значення: \textbf{admin, security\_\allowbreak{}officer}\vspace{1mm}\newline\noindent\rule{\linewidth}{0.4pt}\vspace{1mm}
Параметр: Визначено персонал або ролі, яким мають бути обмежені привілейовані облікові записи в системі\newline Тип: list\newline Значення: \textbf{admin, security\_\allowbreak{}officer} \\ \hline
13 & МІНІМІЗАЦІЯ ПОВНОВАЖЕНЬ - ПЕРЕГЛЯД ПОВНОВАЖЕНЬ КОРИСТУВАЧА (AC-6(7)) &  & AC-6(7) & Політики управління доступом, привілеями та віддаленим доступом реалізуються централізовано через підсистеми IAM (Identity and Access Management) та ABAC (Attribute-Based Access Control) ЄСІКС з дотриманням принципу найменших привілеїв.\vspace{1mm}\newline\noindent\rule{\linewidth}{0.4pt}\vspace{1mm}
Параметр: Повноваження, призначені ролям і класам, переглядаються з частотою для перевірки необхідності таких повноважень\newline Тип: integer\newline Значення: \textbf{30}\vspace{1mm}\newline\noindent\rule{\linewidth}{0.4pt}\vspace{1mm}
Параметр: Визначено частоту перегляду повноважень, призначених ролям або класам користувачів\newline Тип: integer\newline Значення: \textbf{30}\vspace{1mm}\newline\noindent\rule{\linewidth}{0.4pt}\vspace{1mm}
Параметр: Визначено ролі або класи користувачів, яким призначено повноваження \newline Тип: list\newline Значення: \textbf{admin, security\_\allowbreak{}officer} \\ \hline
14 & МІНІМІЗАЦІЯ ПОВНОВАЖЕНЬ - АУДИТ ВИКОРИСТАННЯ ПРИВІЛЕЙОВАНИХ ФУНКЦІЙ (AC-6(9)) &  & AC-6(9) & Політики управління доступом, привілеями та віддаленим доступом реалізуються централізовано через підсистеми IAM (Identity and Access Management) та ABAC (Attribute-Based Access Control) ЄСІКС з дотриманням принципу найменших привілеїв. \\ \hline
15 & МІНІМІЗАЦІЯ ПОВНОВАЖЕНЬ - ЗАБОРОНА НЕПРИВІЛЕЙОВАНИМ КОРИСТУВАЧАМ ВИКОНУВАТИ ПРИВІЛЕЙОВАНІ ФУНКЦІЇ (AC-6(10)) &  & AC-6(10) & Політики управління доступом, привілеями та віддаленим доступом реалізуються централізовано через підсистеми IAM (Identity and Access Management) та ABAC (Attribute-Based Access Control) ЄСІКС з дотриманням принципу найменших привілеїв. \\ \hline
16 & НЕВДАЛІ СПРОБИ ВХОДУ В СИСТЕМУ (AC-7) & a. Встановити обмеження на [Призначення: визначену організацією кількість] послідовних неуспішних спроб входу користувача в систему впродовж [Призначення: визначеного організацією часового періоду].\newline b. Автоматично виконати [Вибір (один або декілька): блокування облікового запису/вузла на [Призначення: визначений організацією часовий період]; блокування облікового запису/вузла, доки він не буде розблокований адміністратором; затримання наступної команди входу в систему за [Надання: визначеним організацією алгоритмом затримки]; виконати [Призначення: визначені організацією дії]], коли перевищено максимальну кількість невдалих спроб входу в систему. & AC-7 & Політики управління доступом, привілеями та віддаленим доступом реалізуються централізовано через підсистеми IAM (Identity and Access Management) та ABAC (Attribute-Based Access Control) ЄСІКС з дотриманням принципу найменших привілеїв.\vspace{1mm}\newline\noindent\rule{\linewidth}{0.4pt}\vspace{1mm}
Параметр: Застосовано обмеження на кількість послідовних неуспішних спроб входу користувача протягом періоду часу\newline Тип: integer\newline Значення: \textbf{30}\vspace{1mm}\newline\noindent\rule{\linewidth}{0.4pt}\vspace{1mm}
Параметр: Визначається кількість послідовних неуспішних спроб входу користувача, дозволених протягом певного періоду часу\newline Тип: integer\newline Значення: \textbf{30}\vspace{1mm}\newline\noindent\rule{\linewidth}{0.4pt}\vspace{1mm}
Параметр: Визначено період часу, яким обмежується кількість послідовних неуспішних спроб входу користувача\newline Тип: integer\newline Значення: \textbf{30}\vspace{1mm}\newline\noindent\rule{\linewidth}{0.4pt}\vspace{1mm}
Параметр: Вибрано одне або декілька з наступних ЗНАЧЕНЬ ПАРАМЕТРІВ: \{заблокувати обліковий запис або вузол на період часу; заблокувати обліковий запис або вузол до зняття адміністратором; затримати наступний запит на вхід за алгоритмом затримки; повідомити системного адміністратора; виконати іншу дію\}\newline Тип: string\newline Значення: \textbf{AES-256-GCM} \\ 
 &  &  &  & Параметр: Період часу, на який буде заблоковано обліковий запис або вузол (якщо вибрано)\newline Тип: integer\newline Значення: \textbf{30}\vspace{1mm}\newline\noindent\rule{\linewidth}{0.4pt}\vspace{1mm}
Параметр: Визначено алгоритм затримки наступного запиту на вхід (якщо вибрано)\newline Тип: string\newline Значення: \textbf{AES-256-GCM}\vspace{1mm}\newline\noindent\rule{\linewidth}{0.4pt}\vspace{1mm}
Параметр: Інша дія, яка буде виконана після перевищення максимальної кількості невдалих спроб (якщо вибрано)\newline Тип: integer\newline Значення: \textbf{3} \\ \hline
17 & НЕВДАЛІ СПРОБИ ВХОДУ В СИСТЕМУ - ОЧИЩЕННЯ АБО СТИРАННЯ МОБІЛЬНОГО ПРИСТРОЮ (AC-7(2)) &  & AC-7(2) & Політики управління доступом, привілеями та віддаленим доступом реалізуються централізовано через підсистеми IAM (Identity and Access Management) та ABAC (Attribute-Based Access Control) ЄСІКС з дотриманням принципу найменших привілеїв.\vspace{1mm}\newline\noindent\rule{\linewidth}{0.4pt}\vspace{1mm}
Параметр: Інформація очищується або стирається з мобільних пристроїв на основі вимог або методів очищення або стирання після <AC-07(02)\_\allowbreak{}ODP[03 кількість> послідовних, невдалих спроб входу на пристрій\newline Тип: integer\newline Значення: \textbf{3}\vspace{1mm}\newline\noindent\rule{\linewidth}{0.4pt}\vspace{1mm}
Параметр: Визначається кількість послідовних невдалих спроб входу в систему до того, як інформація буде очищена або стерта з мобільних пристроїв\newline Тип: integer\newline Значення: \textbf{3} \\ \hline
18 & ПОПЕРЕДЖЕННЯ ПРО ВИКОРИСТАННЯ СИСТЕМИ (AC-8) & a.\newline b. Демонструвати користувачам [Призначення: визначене організацією сповіщення або банер про використання системи] перед тим, як надавати доступ до системи, що забезпечує безпеку та приватність відповідно до чинних законів, нормативних документів, наказів, директив, політик, правил, стандартів і керівних принципів, які зазначають, що:\newline 1. користувачі здійснюють доступ до урядової системи;\newline 2. використання системи може контролюватися, реєструватися та підлягати аудиту;\newline 3. несанкціоноване використання системи забороняється та приводить до кримінальної та цивільної відповідальності;\newline 4. використання системи означає згоду на моніторинг і запис дій користувача. Зберігати сповіщення або банер на екрані, доки користувачі не визнають умови використання та не приймуть явних дій для входу в систему або подальшого доступу до системи.\newline c. Для загальнодоступних систем:\newline 1. демонструвати інформацію про умови використання системи [Призначення: визначені організацією умови], перш ніж надавати подальший доступ до загальнодоступної системи;\newline 2. демонструвати посилання, якщо такі є, на моніторинг, запис або аудит, які узгоджуються з акомодацією приватності для таких систем, які зазвичай забороняють такі дії;\newline 3. мати опис авторизованого використання системи. & AC-8 & Політики управління доступом, привілеями та віддаленим доступом реалізуються централізовано через підсистеми IAM (Identity and Access Management) та ABAC (Attribute-Based Access Control) ЄСІКС з дотриманням принципу найменших привілеїв.\vspace{1mm}\newline\noindent\rule{\linewidth}{0.4pt}\vspace{1mm}
Параметр: Визначається сповіщення або банер про використання системи, який система буде показувати користувачам перед наданням доступу до системи; визначено умови використання системи, які система відображатиме перед наданням подальшого доступу; повідомлення про використання системи відображається користувачам перед наданням доступу до системи, яке містить повідомлення про конфіденційність і безпеку, що відповідають чинним законам, нормативним документам, наказам, директивам, політикам, правилам, стандартам і керівним принципам; у повідомленні про використання системи зазначено, що користувачі отримують доступ до урядової системи; у повідомленні про використання системи зазначено, що використання системи може контролюватися, реєструватися та підлягати аудиту; у повідомленні про використання системи зазначено, що несанкціоноване використання системи заборонено і тягне за собою кримінальну та цивільну відповідальність; у повідомленні про використання системи зазначено, що використання системи означає згоду на моніторинг і запис дій; сповіщення або банер залишається на екрані до тих пір, поки користувачі не визнають умови використання та не приймуть явні дії для входу в систему або подальшого доступу до неї; для загальнодоступних систем інформація про використання системи умови відображається перед наданням подальшого доступу до загальнодоступної системи; для загальнодоступних систем відображаються будь-які посилання на моніторинг, запис або аудит, які узгоджуються з положеннями про конфіденційність таких систем, що зазвичай забороняють ці види діяльності; для загальнодоступних систем додається опис авторизованого використання системи\newline Тип: list\newline Значення: \textbf{default\_\allowbreak{}deny\_\allowbreak{}rule, abac\_\allowbreak{}rule\_\allowbreak{}1} \\ \hline
19 & БЛОКУВАННЯ ПРИСТРОЮ (AC-11) & a. Заборонити подальший доступ до системи шляхом ініціювання блокування пристрою після [Призначення: визначеного організацією періоду] бездіяльності або після отримання запиту від користувача.\newline b. Зберігати блокування пристрою, поки користувач не відновить доступ, використовуючи встановлені процедури ідентифікації та автентифікації. & AC-11 & Сеанс користувача автоматично блокується системою (`synrc/n2o` session expiry) після визначеного періоду бездіяльності, вимагаючи повторної автентифікації.\vspace{1mm}\newline\noindent\rule{\linewidth}{0.4pt}\vspace{1mm}
Параметр: session\_\allowbreak{}lock\_\allowbreak{}timeout\newline Тип: integer\newline Значення: \textbf{900} \\ \hline
20 & БЛОКУВАННЯ ПРИСТРОЮ - ПРИХОВАНІ ДИСПЛЕЇ (AC-11(1)) &  & AC-11(1) & Політики управління доступом, привілеями та віддаленим доступом реалізуються централізовано через підсистеми IAM (Identity and Access Management) та ABAC (Attribute-Based Access Control) ЄСІКС з дотриманням принципу найменших привілеїв. \\ \hline
21 & Припинення сеансу (AC-12) & Сеанс користувача має завершуватися автоматично після [Призначення: визначених організацією умов або тригерних подій, що вимагають припинення сеансу]. & AC-12 & Користувач може ініціювати завершення сеансу (logout). Крім того, система примусово завершує сеанси при досягненні максимального часу життя токена або зміни контексту безпеки.\vspace{1mm}\newline\noindent\rule{\linewidth}{0.4pt}\vspace{1mm}
Параметр: session\_\allowbreak{}termination\newline Тип: boolean\newline Значення: \textbf{true} \\ \hline
22 & АВТОМАТИЗОВАНЕ МАРКУВАННЯ (AC-15) &  & AC-15 & Політики управління доступом, привілеями та віддаленим доступом реалізуються централізовано через підсистеми IAM (Identity and Access Management) та ABAC (Attribute-Based Access Control) ЄСІКС з дотриманням принципу найменших привілеїв. \\ \hline
23 & АТРИБУТИ БЕЗПЕКИ ТА ПРИВАТНОСТІ (AC-16) & a. Визначити засоби для асоціювання (пов'язання) [Призначення: визначених організацією типів атрибутів безпеки та приватності], що приймають [Призначення: визначені організацією значення атрибутів безпеки та приватності] з інформацією, яка зберігається, обробляється та/або передається.\newline b. Пов'язані атрибути безпеки та приватності мають створюватися і зберігатися разом з інформацією.\newline c. Встановити дозволені [Призначення: визначені організацією атрибути безпеки] для [Призначення: систем, визначених організацією].\newline d. Визначити дозволені [Призначення: визначені організацією значення або діапазони] для кожного з встановлених атрибутів безпеки та приватності. & AC-16 & Політики управління доступом, привілеями та віддаленим доступом реалізуються централізовано через підсистеми IAM (Identity and Access Management) та ABAC (Attribute-Based Access Control) ЄСІКС з дотриманням принципу найменших привілеїв.\vspace{1mm}\newline\noindent\rule{\linewidth}{0.4pt}\vspace{1mm}
Параметр: Визначено типи атрибутів безпеки, які мають бути пов'язані зі значеннями атрибутів безпеки для інформації, що зберігається, обробляється та/або передається; визначено значення атрибутів конфіденційності для типів атрибутів конфіденційності; визначено системи, для яких мають бути встановлені дозволені атрибути безпеки; визначено системи, для яких мають бути встановлені дозволені атрибути конфіденційності; визначено атрибути безпеки, визначені як частина AC-16(a), які дозволені для систем; визначено атрибути конфіденційності, визначені як частина AC-16(a), які дозволені для систем; визначено значення атрибутів або діапазони для встановлених атрибутів; визначено частоту, з якою слід переглядати атрибути безпеки на предмет відповідності; визначено частоту, з якою слід переглядати атрибути конфіденційності на предмет відповідності; надано засоби для асоціювання типів атрибутів безпеки зі значеннями атрибутів безпеки для інформації, що зберігається, обробляється та/або передається надано засоби для асоціювання типів атрибутів конфіденційності зі значеннями атрибутів конфіденційності для інформації, що зберігається, обробляється та/або передається виникають зв'язки з атрибутами; зв'язки атрибутів зберігаються разом з інформацією; на основі атрибутів, визначених у AC-16\_\allowbreak{}ODP[01] для систем, встановлюються наступні дозволені атрибути безпеки: атрибути безпеки; на основі атрибутів, визначених у AC-16\_\allowbreak{}ODP[02] для систем, встановлюються наступні дозволені атрибути приватності: атрибути конфіденційності ; AC-16(d)\newline Тип: list\newline Значення: \textbf{default\_\allowbreak{}deny\_\allowbreak{}rule, abac\_\allowbreak{}rule\_\allowbreak{}1}\vspace{1mm}\newline\noindent\rule{\linewidth}{0.4pt}\vspace{1mm}
Параметр: Визначено наступні допустимі значення або діапазони атрибутів для кожного з встановлених атрибутів: значення або діапазони атрибутів; проводиться аудит змін до атрибутів; атрибути безпеки перевіряються на відповідність частота; атрибути конфіденційності перевіряються на відповідність частота\newline Тип: list\newline Значення: \textbf{default\_\allowbreak{}deny\_\allowbreak{}rule, abac\_\allowbreak{}rule\_\allowbreak{}1} \\ \hline
24 & ВІДДАЛЕНИЙ ДОСТУП (AC-17) & a. Встановити та задокументувати обмеження на використання, вимоги до конфігурації/підключення та рекомендації щодо здійснення кожного типу віддаленого доступу.\newline b. Авторизувати віддалений доступ до системи, перш ніж будуть дозволені такі підключення. & AC-17 & Віддалений доступ до системи керується додатковими політиками (наприклад, перевірка IP-адреси, VPN, mTLS), які інтегруються з `synrc/ca` та `abac` для забезпечення безпеки доступу за межами захищеного периметра.\vspace{1mm}\newline\noindent\rule{\linewidth}{0.4pt}\vspace{1mm}
Параметр: remote\_\allowbreak{}access\_\allowbreak{}control\newline Тип: boolean\newline Значення: \textbf{true} \\ \hline
25 & ВІДДАЛЕНИЙ ДОСТУП - КЕРОВАНІ ТОЧКИ КОНТРОЛЮ ДОСТУПУ (AC-17(3)) &  & AC-17(3) & Політики управління доступом, привілеями та віддаленим доступом реалізуються централізовано через підсистеми IAM (Identity and Access Management) та ABAC (Attribute-Based Access Control) ЄСІКС з дотриманням принципу найменших привілеїв. \\ \hline
26 & ВІДДАЛЕНИЙ ДОСТУП - ПРИВІЛЕЙОВАНІ КОМАНДИ ТА ДОСТУП (AC-17(4)) &  & AC-17(4) & Політики управління доступом, привілеями та віддаленим доступом реалізуються централізовано через підсистеми IAM (Identity and Access Management) та ABAC (Attribute-Based Access Control) ЄСІКС з дотриманням принципу найменших привілеїв. \\ \hline
27 & БЕЗДРОТОВИЙ ДОСТУП (AC-18) & a. Установити обмеження на використання, вимоги до конфігурації/підключення та рекомендації щодо здійснення бездротового доступу.\newline b. Авторизувати бездротовий доступ до системи, перш ніж будуть дозволені такі підключення. & AC-18 &  \\ \hline
28 & БЕЗДРОТОВИЙ ДОСТУП - АВТЕНТИФІКАЦІЯ ТА ШИФРУВАННЯ (AC-18(1)) &  & AC-18(1) & Політики управління доступом, привілеями та віддаленим доступом реалізуються централізовано через підсистеми IAM (Identity and Access Management) та ABAC (Attribute-Based Access Control) ЄСІКС з дотриманням принципу найменших привілеїв.\vspace{1mm}\newline\noindent\rule{\linewidth}{0.4pt}\vspace{1mm}
Параметр: Бездротовий доступ до системи захищений за допомогою шифрування\newline Тип: string\newline Значення: \textbf{AES-256-GCM} \\ \hline
29 & БЕЗДРОТОВИЙ ДОСТУП - ВІДКЛЮЧЕННЯ БЕЗДРОТОВОЇ МЕРЕЖІ (AC-18(3)) &  & AC-18(3) & Політики управління доступом, привілеями та віддаленим доступом реалізуються централізовано через підсистеми IAM (Identity and Access Management) та ABAC (Attribute-Based Access Control) ЄСІКС з дотриманням принципу найменших привілеїв. \\ \hline
30 & КОНТРОЛЬ ДОСТУПУ ДЛЯ МОБІЛЬНИХ ПРИСТРОЇВ (AC-19) & a. Встановити обмеження на використання, вимоги до конфігурації, вимоги до підключення і рекомендації щодо впровадження мобільних пристроїв, контрольованих організацією.\newline b. Авторизувати підключення мобільних пристроїв до систем, які експлуатуються організацією. & AC-19 & Політики управління доступом, привілеями та віддаленим доступом реалізуються централізовано через підсистеми IAM (Identity and Access Management) та ABAC (Attribute-Based Access Control) ЄСІКС з дотриманням принципу найменших привілеїв. \\ \hline
31 & КОНТРОЛЬ ДОСТУПУ ДЛЯ МОБІЛЬНИХ ПРИСТРОЇВ - ВИКОРИСТАННЯ ПИСЬМОВИХ ТА ПОРТАТИВНИЙ ПРИСТРОЇВ ДЛЯ ЗБЕРІГАННЯ ДАНИХ (AC-19(1)) &  & AC-19(1) & Політики управління доступом, привілеями та віддаленим доступом реалізуються централізовано через підсистеми IAM (Identity and Access Management) та ABAC (Attribute-Based Access Control) ЄСІКС з дотриманням принципу найменших привілеїв. \\ \hline
32 & КОНТРОЛЬ ДОСТУПУ ДЛЯ МОБІЛЬНИХ ПРИСТРОЇВ - ВИКОРИСТАННЯ ПЕРСОНАЛЬНИХ ПОРТАТИВНИХ ПРИСТРОЇВ ЗБЕРІГАННЯ ДАНИХ (AC-19(2)) &  & AC-19(2) & Політики управління доступом, привілеями та віддаленим доступом реалізуються централізовано через підсистеми IAM (Identity and Access Management) та ABAC (Attribute-Based Access Control) ЄСІКС з дотриманням принципу найменших привілеїв.\vspace{1mm}\newline\noindent\rule{\linewidth}{0.4pt}\vspace{1mm}
Параметр: КОНТРОЛЬ ДОСТУПУ ДЛЯ МОБІЛЬНИХ ПРИСТРОЇВ - ВИКОРИСТАННЯ ПЕРСОНАЛЬНИХ ПОРТАТИВНИХ ПРИСТРОЇВ ЗБЕРІГАННЯ ДАНИХ [Вилучено: Включено в MP-07]. AC-19(03) КОНТРОЛЬ ДОСТУПУ ДЛЯ МОБІЛЬНИХ ПРИСТРОЇВ - ВИКОРИСТАННЯ ПОРТАТИВНИХ ПРИСТРОЇВ ЗБЕРІГАННЯ ДАНИХ З НЕІДЕНТИФІКОВАНИМ ВЛАСНИКОМ [Вилучено: Включено в MP-07]\newline Тип: list\newline Значення: \textbf{admin, security\_\allowbreak{}officer} \\ \hline
33 & КОНТРОЛЬ ДОСТУПУ ДЛЯ МОБІЛЬНИХ ПРИСТРОЇВ - ОБМЕЖЕННЯ ДЛЯ ЗАСЕКРЕЧЕНОЇ ІНФОРМАЦІЇ (AC-19(4)) &  & AC-19(4) & Політики управління доступом, привілеями та віддаленим доступом реалізуються централізовано через підсистеми IAM (Identity and Access Management) та ABAC (Attribute-Based Access Control) ЄСІКС з дотриманням принципу найменших привілеїв.\vspace{1mm}\newline\noindent\rule{\linewidth}{0.4pt}\vspace{1mm}
Параметр: Використання незахищених мобільних пристроїв на об'єктах, що містять системи, які обробляють, зберігають або передають секретну інформацію, заборонено, за винятком випадків, коли на це є спеціальний дозвіл уповноваженої посадової особи; AC-19(04)(b)(01) заборона підключення незахищених мобільних пристроїв до систем з обмеженим доступом застосовується до осіб, яким уповноважена посадова особа дозволила використовувати незахищені мобільні пристрої на об'єктах, що містять системи, які обробляють, зберігають або передають інформацію з обмеженим доступом; AC-19(04)(b)(02) дозвіл уповноваженої посадової особи на підключення незахищених мобільних пристроїв до незахищених систем вимагається від осіб, яким дозволено використовувати незахищені мобільні пристрої на об'єктах, що містять системи, які обробляють, зберігають або передають інформацію з обмеженим доступом; AC-19(04)(b)(03) заборона використання внутрішніх або зовнішніх модемів чи бездротових інтерфейсів у складі незахищених мобільних пристроїв поширюється на осіб, яким уповноваженою посадовою особою дозволено використовувати незахищені мобільні пристрої під час виконання службових обов'язків, а також на осіб, які не мають права на використання таких пристроїв AC-19(04)(b)(04)[01] вибірковий огляд та перевірка незахищених мобільних пристроїв та інформації, що зберігається на них, посадовими особами є обов'язковими; AC-19(04)(b)(04)[02] дотримання політики обробки інцидентів застосовується у разі виявлення секретної інформації під час випадкового огляду та перевірки незахищених мобільних пристроїв\newline Тип: list\newline Значення: \textbf{admin, security\_\allowbreak{}officer}\vspace{1mm}\newline\noindent\rule{\linewidth}{0.4pt}\vspace{1mm}
Параметр: Підключення захищенихмобільних пристроїв до засекречених систем обмежено відповідно до політики безпеки\newline Тип: list\newline Значення: \textbf{default\_\allowbreak{}deny\_\allowbreak{}rule, abac\_\allowbreak{}rule\_\allowbreak{}1}\vspace{1mm}\newline\noindent\rule{\linewidth}{0.4pt}\vspace{1mm}
Параметр: Визначено посадових осіб із захисту інформації, відповідальних за огляд та перевірку захищених мобільних пристроїв та інформації, що зберігається на цих пристроях\newline Тип: list\newline Значення: \textbf{admin, security\_\allowbreak{}officer}\vspace{1mm}\newline\noindent\rule{\linewidth}{0.4pt}\vspace{1mm}
Параметр: Визначено політики безпеки, що обмежують підключення захищених мобільних пристроїв до засекречених систем\newline Тип: list\newline Значення: \textbf{default\_\allowbreak{}deny\_\allowbreak{}rule, abac\_\allowbreak{}rule\_\allowbreak{}1} \\ \hline
34 & КОНТРОЛЬ ДОСТУПУ ДЛЯ МОБІЛЬНИХ ПРИСТРОЇВ - ПОВНЕ ШИФРУВАННЯ ПРИСТРОЇВ ТА СХОВИЩ ІНФОРМАЦІЇ (AC-19(5)) &  & AC-19(5) & Політики управління доступом, привілеями та віддаленим доступом реалізуються централізовано через підсистеми IAM (Identity and Access Management) та ABAC (Attribute-Based Access Control) ЄСІКС з дотриманням принципу найменших привілеїв. \\ \hline
35 & ВИКОРИСТАННЯ ЗОВНІШНІХ СИСТЕМ (AC-20) & a. [Вибір (один або кілька): Встановіть [Призначення: умови, визначені організацією]; Визначте [Призначення: визначені організацією засоби контролю, які, як стверджується, будуть реалізовані на зовнішніх системах]], узгоджені з довірчими відносинами, встановленими з іншими організаціями, які володіють, експлуатують та/або обслуговують зовнішні системи, дозволяючи уповноваженим особам:\newline 1. доступ до системи із зовнішніх систем;\newline 2. обробляти, зберігати або передавати керовану організацією інформацію за допомогою зовнішніх систем;\newline b. Заборонити використання [Призначення: організаційно-визначені типи зовнішніх систем]. & AC-20 & Політики управління доступом, привілеями та віддаленим доступом реалізуються централізовано через підсистеми IAM (Identity and Access Management) та ABAC (Attribute-Based Access Control) ЄСІКС з дотриманням принципу найменших привілеїв.\vspace{1mm}\newline\noindent\rule{\linewidth}{0.4pt}\vspace{1mm}
Параметр: Вибрано одне або більше з наступних ЗНАЧЕНЬ ПАРАМЕТРІВ: \{встановити умови та положення; визначити заходи захисту\}\newline Тип: list\newline Значення: \textbf{}\vspace{1mm}\newline\noindent\rule{\linewidth}{0.4pt}\vspace{1mm}
Параметр: Визначено умови та положення, що відповідають довірчим відносинам, встановленим з іншими організаціями, які володіють, експлуатують та/або обслуговують зовнішні системи (якщо вибрано)\newline Тип: list\newline Значення: \textbf{} \\ \hline
36 & ВИКОРИСТАННЯ ЗОВНІШНІХ СИСТЕМ - ОБМЕЖЕННЯ НА АВТОРИЗОВАНЕ ВИКОРИСТАННЯ (AC-20(1)) &  & AC-20(1) & Політики управління доступом, привілеями та віддаленим доступом реалізуються централізовано через підсистеми IAM (Identity and Access Management) та ABAC (Attribute-Based Access Control) ЄСІКС з дотриманням принципу найменших привілеїв.\vspace{1mm}\newline\noindent\rule{\linewidth}{0.4pt}\vspace{1mm}
Параметр: Авторизовані особи мають право використовувати зовнішню систему для доступу до системи або для обробки, зберігання чи передачі інформації, що контролюється організацією, лише після перевірки виконання заходів безпеки та конфіденційності, зазначених у політиці безпеки та конфіденційності організації, а також планах безпеки та конфіденційності (якщо такі застосовуються)\newline Тип: list\newline Значення: \textbf{admin, security\_\allowbreak{}officer}\vspace{1mm}\newline\noindent\rule{\linewidth}{0.4pt}\vspace{1mm}
Параметр: Авторизовані особи мають право використовувати зовнішню систему для доступу до системи або для обробки, зберігання чи передачі інформації, що контролюється організацією, лише після збереження погоджених угод про підключення або обробку системи з структурою організації, на якій розміщена зовнішня система\newline Тип: list\newline Значення: \textbf{admin, security\_\allowbreak{}officer} \\ \hline
37 & ВИКОРИСТАННЯ ЗОВНІШНІХ СИСТЕМ - ПЕРЕНОСНІ ПРИСТРОЇ ЗБЕРІГАННЯ ДАНИХ (AC-20(2)) &  & AC-20(2) & Політики управління доступом, привілеями та віддаленим доступом реалізуються централізовано через підсистеми IAM (Identity and Access Management) та ABAC (Attribute-Based Access Control) ЄСІКС з дотриманням принципу найменших привілеїв.\vspace{1mm}\newline\noindent\rule{\linewidth}{0.4pt}\vspace{1mm}
Параметр: Використання портативних пристроїв носіїв інформації уповноваженими особами обмежено у зовнішніх системах за допомогою обмеження \newline Тип: list\newline Значення: \textbf{admin, security\_\allowbreak{}officer}\vspace{1mm}\newline\noindent\rule{\linewidth}{0.4pt}\vspace{1mm}
Параметр: Визначено обмеження на використання авторизованими особами портативних носіїв інформації у зовнішніх системах\newline Тип: list\newline Значення: \textbf{admin, security\_\allowbreak{}officer} \\ \hline
38 & ПУБЛІЧНО ДОСТУПНИЙ КОНТЕНТ (AC-22) & a. Призначити осіб, що уповноважені на розміщення інформації в загальнодоступній системі.\newline b. Навчати уповноважених осіб тому, щоб загальнодоступна інформація не містила інформацію з обмеженим доступом.\newline c. Переглядати запропонований зміст інформації до публікації в загальнодоступній системі, щоб гарантувати, що там не міститься інформація з обмеженим доступом.\newline d. Переглядати вміст загальнодоступної системи на предмет наявності там інформації з обмеженим доступом з [Призначення: визначеною організацією частотою]; така інформація має бути видалена в разі її виявлення. & AC-22 & Політики управління доступом, привілеями та віддаленим доступом реалізуються централізовано через підсистеми IAM (Identity and Access Management) та ABAC (Attribute-Based Access Control) ЄСІКС з дотриманням принципу найменших привілеїв.\vspace{1mm}\newline\noindent\rule{\linewidth}{0.4pt}\vspace{1mm}
Параметр: Визначені особи уповноважені на розміщення інформації в загальнодоступній системі\newline Тип: list\newline Значення: \textbf{admin, security\_\allowbreak{}officer}\vspace{1mm}\newline\noindent\rule{\linewidth}{0.4pt}\vspace{1mm}
Параметр: Уповноважені особи проходять навчання, щоб гарантувати, що загальнодоступна інформація не містить інформацію з обмеженим доступом\newline Тип: list\newline Значення: \textbf{admin, security\_\allowbreak{}officer}\vspace{1mm}\newline\noindent\rule{\linewidth}{0.4pt}\vspace{1mm}
Параметр: Зміст у загальнодоступній системі перевіряється на наявність інформації з обмеженим доступом з частотою\newline Тип: integer\newline Значення: \textbf{30}\vspace{1mm}\newline\noindent\rule{\linewidth}{0.4pt}\vspace{1mm}
Параметр: Визначено частоту, з якою слід переглядати вміст загальнодоступної системи на предмет наявності там інформації з обмеженим доступом\newline Тип: integer\newline Значення: \textbf{30} \\ \hline
\multicolumn{5}{|l|}{\textbf{2. ОБІЗНАНІСТЬ ТА НАВЧАННЯ (AT)}} \\ \hline
39 & ПОЛІТИКА ТА ПРОЦЕДУРИ ПІДВИЩЕННЯ ОБІЗНАНОСТІ ТА НАВЧАННЯ (AT-1) & a. Розробити, задокументувати та поширити [Призначення: серед визначеного організацією персоналу або ролей]:\newline 1. [Вибір (один або декілька): Рівень організації; Рівень місії/бізнес-процесу; рівень системи] політики обізнаності та навчання у сфері забезпечення безпеки та приватності, яка:\newline (a) містить мету, сферу застосування, ролі, обов'язки, відповідальність керівництва, координацію між організаційними підрозділами та систему контролю відповідності (complaince);\newline (b) відповідає чинним законам, нормативним документам, директивам, нормам, політикам, стандартам та керівним документам.\newline 2. Процедури, що сприяють реалізації політики підвищення обізнаності та професійної підготовки в галузі безпеки, приватності, а також пов'язаних з ними заходів захисту інформації та персональних даних.\newline b. Призначити [Призначення: визначену організацією посадову особу] для управління політикою та процедурами підвищення обізнаності та навчання у сфері забезпечення безпеки та приватності.\newline c. Переглядати та оновлювати:\newline 1. Поточну політика [Призначення: частота, визначена організацією] і наступне [Призначення: події, визначені організацією];\newline 2. Процедури [Призначення: частота, визначена організацією] та наступні [Призначення: події, визначені організацією]. & AT-1 & Відповідальність за навчання та обізнаність з питань кібербезпеки несе організація-користувач відповідно до загальних політик ЄСІКС. \\ \hline
40 & НАВЧАННЯ З ПІДВИЩЕННЯ ОБІЗНАНОСТІ (AT-2) & Впровадити базові тренінги з підвищення обізнаності у сфері безпеки та приватності для користувачів системи (включно з менеджерами, керівниками компаній і підрядниками):\newline a. Забезпечити навчання грамотності з питань безпеки та конфіденційності для користувачів системи (включаючи менеджерів, керівників вищої ланки та підрядників):\newline 1. як частину початкового навчання для нових користувачів і [Призначення: частота, визначена організацією] після цього;\newline 2. якщо цього потребують системні зміни або наступні [Призначення: події, визначені організацією].\newline b. Використовувати наведені нижче методи, щоб підвищити рівень безпеки та конфіденційності користувачів системи [Завдання: визначені організацією методи поінформованості];\newline c. Оновлювати навчання грамотності та зміст обізнаності [Завдання: частота, визначена організацією] і наступні [Завдання: події, визначені організацією];\newline d. Включити уроки, отримані з внутрішніх або зовнішніх інцидентів безпеки або порушень, у навчання грамотності та методи підвищення обізнаності. & AT-2 & Програми навчання та інструктажі (див. розділ AT у digital-profile.pdf).\vspace{1mm}\newline\noindent\rule{\linewidth}{0.4pt}\vspace{1mm}
Параметр: security\_\allowbreak{}awareness\_\allowbreak{}training\_\allowbreak{}ref\newline Тип: string\newline Значення: \textbf{digital-profile.pdf} \\ \hline
41 & НАВЧАННЯ З ПІДВИЩЕННЯ ОБІЗНАНОСТІ | НАВЧАННЯ ЩОДО РОЗПІЗНАВАННЯ ВНУТРІШНЬОЇ ЗАГРОЗИ &  & AT-2(2) & Відповідальність за навчання та обізнаність з питань кібербезпеки несе організація-користувач відповідно до загальних політик ЄСІКС. \\ \hline
42 & НАВЧАННЯ З ПІДВИЩЕННЯ ОБІЗНАНОСТІ | НАВЧАННЯ З ПИТАНЬ СОЦІАЛЬНОЇ ІНЖЕНЕРІЇ &  & AT-2(3) & Відповідальність за навчання та обізнаність з питань кібербезпеки несе організація-користувач відповідно до загальних політик ЄСІКС. \\ \hline
43 & РОЛЬОВЕ НАВЧАННЯ (AT-3) & a. Забезпечити проведення навчання з питань безпеки та приватності на основі ролей для працівників з ролями та обов'язками: [Призначення: визначені організацією ролі та обов'язки]:\newline 1. перед авторизацією доступу до системи, інформації або виконанням призначених обов'язків і [Призначення: частота, визначена організацією] після цього;\newline 2. коли цього потребують системні зміни.\newline b. Оновити навчальний контент на основі ролей [Призначення: частота, визначена організацією] і наступні [Призначення: події, визначені організацією];\newline c. Включіть у рольове навчання, інформацію, отриману з внутрішніх або зовнішніх інцидентів та порушень безпеки. & AT-3 & Спеціалізоване навчання для адміністраторів (див. розділ AT у digital-profile.pdf).\vspace{1mm}\newline\noindent\rule{\linewidth}{0.4pt}\vspace{1mm}
Параметр: role\_\allowbreak{}based\_\allowbreak{}training\_\allowbreak{}ref\newline Тип: string\newline Значення: \textbf{digital-profile.pdf} \\ \hline
\multicolumn{5}{|l|}{\textbf{3. АУДИТ ТА ПІДЗВІТНІСТЬ (AU)}} \\ \hline
44 & Політика та процедури аудиту та підзвітності (AU-1) & a. Розробити, задокументувати та поширити [Призначення: серед персоналу або ролей, що їх визначила організація]:\newline 1. [Вибір (один або декілька): Рівень організації; Рівень місії/бізнес-процесу; рівень системи] політика аудиту та підзвітності, яка:\newline (a) містить мету, сферу застосування, ролі, обов'язки, відповідальність керівництва, координацію між організаційними підрозділами та систему контролю відповідності (complaince);\newline (b) відповідає чинним законам, нормативним документам, директивам, нормам, політикам, стандартам та керівним документам.\newline 2. Процедури, що сприяють здійсненню політики аудиту та підзвітності, а також пов'язані з ними заходи аудиту та підзвітності.\newline b. Призначити [Призначення: визначену організацією старшу посадову особу] для управління політикою та процедурами аудиту та підзвітності.\newline c. Переглядати та оновлювати поточний аудит та підзвітність:\newline 1. політику [Призначення: частота, визначена організацією] та наступне [Призначення: події, визначені організацією];\newline 2. процедури аудиту [Призначення: визначеною організацією частотою] та [Завдання: події, визначені організацією]. & AU-1 & Підсистема логування та аудиту автоматично фіксує всі події доступу (автентифікації/авторизації) та зміни даних. Журнали аудиту зберігаються в імутабельному сховищі не менше 3 місяців.\vspace{1mm}\newline\noindent\rule{\linewidth}{0.4pt}\vspace{1mm}
Параметр: Розроблено та задокументовано політику аудиту та підзвітності\newline Тип: list\newline Значення: \textbf{default\_\allowbreak{}deny\_\allowbreak{}rule, abac\_\allowbreak{}rule\_\allowbreak{}1}\vspace{1mm}\newline\noindent\rule{\linewidth}{0.4pt}\vspace{1mm}
Параметр: Політика аудиту та підзвітності доведена до персонал або ролі\newline Тип: list\newline Значення: \textbf{admin, security\_\allowbreak{}officer}\vspace{1mm}\newline\noindent\rule{\linewidth}{0.4pt}\vspace{1mm}
Параметр: Розроблені та задокументовані процедури аудиту та підзвітності, що сприяють впровадженню політики аудиту та підзвітності, а також відповідні заходи контролю аудиту та підзвітності\newline Тип: list\newline Значення: \textbf{default\_\allowbreak{}deny\_\allowbreak{}rule, abac\_\allowbreak{}rule\_\allowbreak{}1}\vspace{1mm}\newline\noindent\rule{\linewidth}{0.4pt}\vspace{1mm}
Параметр: Посадова особа призначається для управління політикою та процедурами аудиту та підзвітності AU-01(c)[01][01] переглядається та оновлюється поточна політика аудиту та підзвітності з частота; AU-01(c)[01][02] переглядається та оновлюється поточна політика аудиту та підзвітності після подій; AU-01(c)[02][01] переглядається та оновлюється поточні процедури аудиту та підзвітності з частота; AU-01(c)[02][02] переглядається та оновлюється поточні процедури аудиту та підзвітності після подій\newline Тип: list\newline Значення: \textbf{admin, security\_\allowbreak{}officer} \\ 
 &  &  &  & Параметр: Визначено персонал або ролі, до яких має бути доведена політика аудиту та підзвітності\newline Тип: list\newline Значення: \textbf{admin, security\_\allowbreak{}officer}\vspace{1mm}\newline\noindent\rule{\linewidth}{0.4pt}\vspace{1mm}
Параметр: Визначено персонал або ролі, на які поширюються процедури аудиту та підзвітності\newline Тип: list\newline Значення: \textbf{admin, security\_\allowbreak{}officer}\vspace{1mm}\newline\noindent\rule{\linewidth}{0.4pt}\vspace{1mm}
Параметр: Визначено посадову особу, яка управлятиме політикою та процедурами аудиту та підзвітності\newline Тип: list\newline Значення: \textbf{admin, security\_\allowbreak{}officer}\vspace{1mm}\newline\noindent\rule{\linewidth}{0.4pt}\vspace{1mm}
Параметр: Визначено частоту, з якою переглядається та оновлюється поточна політика аудиту та підзвітності\newline Тип: list\newline Значення: \textbf{default\_\allowbreak{}deny\_\allowbreak{}rule, abac\_\allowbreak{}rule\_\allowbreak{}1} \\ 
 &  &  &  & Параметр: Визначено події, які потребують перегляду та оновлення поточної політики аудиту та підзвітності\newline Тип: list\newline Значення: \textbf{default\_\allowbreak{}deny\_\allowbreak{}rule, abac\_\allowbreak{}rule\_\allowbreak{}1}\vspace{1mm}\newline\noindent\rule{\linewidth}{0.4pt}\vspace{1mm}
Параметр: Визначено частоту, з якою переглядаються та оновлюються поточні процедури аудиту та підзвітності\newline Тип: integer\newline Значення: \textbf{30}\vspace{1mm}\newline\noindent\rule{\linewidth}{0.4pt}\vspace{1mm}
Параметр: Визначено події, які потребують перегляду та оновлення поточної процедури аудиту та підзвітності\newline Тип: list\newline Значення: \textbf{login, logout, failed\_\allowbreak{}attempt} \\ \hline
45 & ПОДІЇ АУДИТУ (AU-2) & a. Визначити типи подій, які система може реєструвати для підтримки функції аудиту: [Призначення: типи подій, визначені організацією, які система здатна реєструвати];\newline b. Координувати функції аудиту безпеки з іншими організаційними підрозділами, які вимагають інформації, пов'язаної з аудитом, для посилення взаємної підтримки та допомоги у виборі типів подій, що перевіряються;\newline c. Визначити, які типи подій підлягають аудиту: [Призначення: визначені організацією події, що підлягають аудиту (підмножина подій, що підлягають аудиту, визначених в AU-2 a.), а також частота (або ситуація, що вимагає) проведення аудиту для кожної ідентифікованої події]\newline d. Обґрунтувати, чому типи подій, що перевіряються, вважаються достатніми для підтримки розслідувань інцидентів (постфактум), пов'язаних з безпекою та приватністю;\newline e. Перегляньте й оновіть типи подій, вибрані для журналювання [Призначення: частота, визначена організацією]. & AU-2 & Підсистема логування та аудиту автоматично фіксує всі події доступу (автентифікації/авторизації) та зміни даних. Журнали аудиту зберігаються в імутабельному сховищі не менше 3 місяців.\vspace{1mm}\newline\noindent\rule{\linewidth}{0.4pt}\vspace{1mm}
Параметр: Зазначені типи подій реєструються 02\_\allowbreak{}ODP[03] частота або ситуація>; <AU-\newline Тип: string\newline Значення: \textbf{щорічно}\vspace{1mm}\newline\noindent\rule{\linewidth}{0.4pt}\vspace{1mm}
Параметр: Переглядаються та оновлюються типи подій, вибрані для реєстрації, частота. у системі\newline Тип: string\newline Значення: \textbf{щорічно}\vspace{1mm}\newline\noindent\rule{\linewidth}{0.4pt}\vspace{1mm}
Параметр: Визначено частоту або ситуацію, що вимагає проведення аудиту для кожної ідентифікованої події\newline Тип: list\newline Значення: \textbf{login, logout, failed\_\allowbreak{}attempt}\vspace{1mm}\newline\noindent\rule{\linewidth}{0.4pt}\vspace{1mm}
Параметр: Частота перегляду та оновлення типів подій, обраних для журналювання\newline Тип: string\newline Значення: \textbf{щорічно} \\ \hline
46 & ЗМІСТ ЗАПИСІВ АУДИТУ (AU-3) & Переконатися, що записи аудиту містять інформацію, яка встановлює наступне:\newline a. який тип події стався;\newline b. коли відбулася подія;\newline c. де відбулася подія;\newline d. джерело події;\newline e. наслідки події;\newline f. результат події та ідентифікатор будь-яких осіб або суб'єктів, пов'язаних з подією. & AU-3 & Інформація в логах BPMN (особливо при збереженні об'єктів BPE/ABAC) містить вичерпні метадані: тип події, суб'єкт, цільовий об'єкт, результат операції (Permit/Deny) та позначку часу, що повністю закриває вимоги AU-3.\vspace{1mm}\newline\noindent\rule{\linewidth}{0.4pt}\vspace{1mm}
Параметр: audit\_\allowbreak{}content\_\allowbreak{}metadata\newline Тип: boolean\newline Значення: \textbf{true} \\ \hline
47 & ЗМІСТ ЗАПИСІВ АУДИТУ - ДОДАТКОВА ІНФОРМАЦІЯ ПРО АУДИТ (AU-3(1)) &  & AU-3(1) & Підсистема логування та аудиту автоматично фіксує всі події доступу (автентифікації/авторизації) та зміни даних. Журнали аудиту зберігаються в імутабельному сховищі не менше 3 місяців. \\ \hline
48 & РЕАГУВАННЯ НА ВІДМОВИ ОБРОБКИ ДАНИХ АУДИТУ (AU-5) & a. Сповіщати [Призначення: визначені організацією персонал або посади] у разі збою обробки даних аудиту в [Призначення: визначений організацією період часу].\newline b. Виконати наступні додаткові дії: [Призначення: визначені організацією дії, які необхідно зробити]. & AU-5 & Підсистема логування та аудиту автоматично фіксує всі події доступу (автентифікації/авторизації) та зміни даних. Журнали аудиту зберігаються в імутабельному сховищі не менше 3 місяців.\vspace{1mm}\newline\noindent\rule{\linewidth}{0.4pt}\vspace{1mm}
Параметр: Персонал або ролі отримують сповіщення у разі збою процесу обробки даних аудиту періоду часу\newline Тип: list\newline Значення: \textbf{admin, security\_\allowbreak{}officer}\vspace{1mm}\newline\noindent\rule{\linewidth}{0.4pt}\vspace{1mm}
Параметр: Додаткові дії виконуються у разі збою процесу обробки даних аудиту\newline Тип: list\newline Значення: \textbf{login, logout, failed\_\allowbreak{}attempt}\vspace{1mm}\newline\noindent\rule{\linewidth}{0.4pt}\vspace{1mm}
Параметр: Визначено персонал або ролі, які отримують сповіщення про збої в процесі обробки даних аудиту \newline Тип: list\newline Значення: \textbf{admin, security\_\allowbreak{}officer}\vspace{1mm}\newline\noindent\rule{\linewidth}{0.4pt}\vspace{1mm}
Параметр: Визначено період часу, протягом якого персонал або ролі отримують сповіщення про збої в процесі обробки даних аудиту\newline Тип: list\newline Значення: \textbf{admin, security\_\allowbreak{}officer} \\ 
 &  &  &  & Параметр: Визначено додаткові дії, яких слід вжити у випадку збою в процесі обробки даних аудиту \newline Тип: list\newline Значення: \textbf{login, logout, failed\_\allowbreak{}attempt} \\ \hline
49 & ОГЛЯД, АНАЛІЗ І ЗВІТНІСТЬ АУДИТУ (AU-6) & a. Переглядати та аналізувати записи системного аудиту [Призначення: з визначеною організацією частотою] для виявлення [Призначення: визначеної організацією неналежної або незвичайної діяльності].\newline b. Відправляти звіт про аудит [Призначення: визначеним організацією персоналу або посадам].\newline c. Налаштувати рівні огляду аудиту, аналізу та звітності в рамках системи, коли змінюється рівень ризику на основі інформації від правоохоронних органів, розвідувальної інформації або від інших достовірних джерел інформації. & AU-6 & Аналіз аудиту відбувається за допомогою консольних інструментів платформи (наприклад, `kvs:hist/1` або інтерфейсів Erlang Shell). Також доступний експорт журналів для SIEM.\vspace{1mm}\newline\noindent\rule{\linewidth}{0.4pt}\vspace{1mm}
Параметр: audit\_\allowbreak{}analysis\_\allowbreak{}tools\newline Тип: boolean\newline Значення: \textbf{true} \\ \hline
50 & ОГЛЯД, АНАЛІЗ І ЗВІТНІСТЬ АУДИТУ - ЗІСТАВЛЯННЯ СХОВИЩ АУДИТУ (AU-6(3)) &  & AU-6(3) & Підсистема логування та аудиту автоматично фіксує всі події доступу (автентифікації/авторизації) та зміни даних. Журнали аудиту зберігаються в імутабельному сховищі не менше 3 місяців. \\ \hline
51 & СКОРОЧЕННЯ ЗАПИСІВ АУДИТУ ТА ФОРМУВАННЯ ЗВІТУ (AU-7) & Забезпечити та реалізувати можливості скорочення записів перевірок аудитом і звітів, до рівня, який:\newline a. підтримує перевірку, аналіз і звітність аудиту на вимогу та розслідування (постфактум) інцидентів безпеки;\newline b. не змінює оригінальний вміст або час упорядкування записів аудиту. & AU-7 & Параметр: Забезпечено можливість скорочення записів перевірок аудитом та звітів, які не змінюють оригінальний зміст або час упорядкування записів аудиту\newline Тип: integer\newline Значення: \textbf{30}\vspace{1mm}\newline\noindent\rule{\linewidth}{0.4pt}\vspace{1mm}
Параметр: Реалізовано можливість скорочення записів перевірок аудитом та звітів, які не змінюють оригінальний зміст або час упорядкування записів аудиту\newline Тип: integer\newline Значення: \textbf{30} \\ \hline
52 & ПОЗНАЧКА ЧАСУ (AU-8) & a. Використовувати внутрішньосистемний годинник для створення позначок часу для записів аудиту.\newline b. Застосовувати позначки часу, які відповідають [Призначення: деталізація вимірювання часу, визначена організацією] і використовують всесвітній координований час, мають фіксоване зміщення місцевого часу відносно всесвітнього координованого часу або включають зміщення місцевого часу як частину позначки часу. & AU-8 & Підсистема логування та аудиту автоматично фіксує всі події доступу (автентифікації/авторизації) та зміни даних. Журнали аудиту зберігаються в імутабельному сховищі не менше 3 місяців.\vspace{1mm}\newline\noindent\rule{\linewidth}{0.4pt}\vspace{1mm}
Параметр: Внутрішній системний годинник використовується для створення позначок часу для записів аудиту\newline Тип: integer\newline Значення: \textbf{30}\vspace{1mm}\newline\noindent\rule{\linewidth}{0.4pt}\vspace{1mm}
Параметр: Позначки часу застосовуються для записів аудиту, які відповідають деталізація вимірювання часу і які використовують всесвітній координований час, мають фіксоване місцеве зміщення місцевого часу від всесвітнього координованого часу або включають зміщення місцевого часу як частину позначки часу\newline Тип: integer\newline Значення: \textbf{30}\vspace{1mm}\newline\noindent\rule{\linewidth}{0.4pt}\vspace{1mm}
Параметр: Визначено деталізацію вимірювання часу для часових позначок записів аудиту\newline Тип: integer\newline Значення: \textbf{30} \\ \hline
53 & ЗАХИСТ ІНФОРМАЦІЇ АУДИТУ (AU-9) & a. Захист інформації аудиту та інструментів несанкціонованого доступу, зміни та видалення; журналювання аудиту від\newline b. Сповіщення [Призначення: персонал або ролі, визначені організацією] у разі виявлення несанкціонованого доступу, зміни або видалення інформації аудиту. & AU-9 & Підсистема логування та аудиту автоматично фіксує всі події доступу (автентифікації/авторизації) та зміни даних. Журнали аудиту зберігаються в імутабельному сховищі не менше 3 місяців.\vspace{1mm}\newline\noindent\rule{\linewidth}{0.4pt}\vspace{1mm}
Параметр: Персонал або ролі отримують сповіщення при виявленні несанкціонованого доступу, зміни або видалення інформації аудиту\newline Тип: list\newline Значення: \textbf{admin, security\_\allowbreak{}officer}\vspace{1mm}\newline\noindent\rule{\linewidth}{0.4pt}\vspace{1mm}
Параметр: Визначено персонал або ролі, які мають бути сповіщені при виявленні несанкціонованого доступу, зміни або видалення інформації аудиту\newline Тип: list\newline Значення: \textbf{admin, security\_\allowbreak{}officer} \\ \hline
54 & ЗАХИСТ ІНФОРМАЦІЇ АУДИТУ - ДОСТУП, ЯКИЙ НАДАЄТЬСЯ ЧЕРЕЗ ЧЛЕНСТВО В ПІДМНОЖИНИ ПРИВІЛЕЙОВАНИХ КОРИСТУВАЧІВ (AU-9(4)) &  & AU-9(4) & Підсистема логування та аудиту автоматично фіксує всі події доступу (автентифікації/авторизації) та зміни даних. Журнали аудиту зберігаються в імутабельному сховищі не менше 3 місяців. \\ \hline
55 & НЕСПРОСТОВНІСТЬ (AU-10) & Надавайте неспростовні докази того, що особа (або процес, який діє від імені особи) виконала [Призначення: дії, визначені організацією, на які поширюється принцип неспростовності]. & AU-10 & Підсистема логування та аудиту автоматично фіксує всі події доступу (автентифікації/авторизації) та зміни даних. Журнали аудиту зберігаються в імутабельному сховищі не менше 3 місяців.\vspace{1mm}\newline\noindent\rule{\linewidth}{0.4pt}\vspace{1mm}
Параметр: Надаються неспростовні докази того, що особа (або процес, що діє від імені особи) виконала дії\newline Тип: list\newline Значення: \textbf{admin, security\_\allowbreak{}officer}\vspace{1mm}\newline\noindent\rule{\linewidth}{0.4pt}\vspace{1mm}
Параметр: Визначено дії, на які поширюється принцип неспростовності\newline Тип: list\newline Значення: \textbf{login, logout, failed\_\allowbreak{}attempt} \\ \hline
56 & НЕСПРОСТОВНІСТЬ - ЦИФРОВІ ПІДПИСИ (AU-10(5)) &  & AU-10(5) & Підсистема логування та аудиту автоматично фіксує всі події доступу (автентифікації/авторизації) та зміни даних. Журнали аудиту зберігаються в імутабельному сховищі не менше 3 місяців. \\ \hline
57 & ЗБЕРЕЖЕННЯ ЗАПИСІВ АУДИТУ (AU-11) & Зберігати записи аудиту впродовж [Призначення: визначеного організацією періоду часу, відповідно політиці зберігання записів], щоб забезпечити підтримку розслідувань (постфактум) інцидентів безпеки та приватності, а також для задоволення вимог нормативних і документів організації щодо збереження даних аудиту. & AU-11 & Підсистема логування та аудиту автоматично фіксує всі події доступу (автентифікації/авторизації) та зміни даних. Журнали аудиту зберігаються в імутабельному сховищі не менше 3 місяців.\vspace{1mm}\newline\noindent\rule{\linewidth}{0.4pt}\vspace{1mm}
Параметр: Записи аудиту зберігаються впродовж період часу, щоб забезпечити підтримку розслідування (постфактум) інцидентів безпеки та конфіденційності, а також відповідати нормативним та вимогам організації щодо збереження даних аудиту\newline Тип: integer\newline Значення: \textbf{30}\vspace{1mm}\newline\noindent\rule{\linewidth}{0.4pt}\vspace{1mm}
Параметр: Визначено період часу для зберігання записів аудиту, який узгоджується з політикою зберігання записів\newline Тип: list\newline Значення: \textbf{default\_\allowbreak{}deny\_\allowbreak{}rule, abac\_\allowbreak{}rule\_\allowbreak{}1} \\ \hline
58 & ГЕНЕРАЦІЯ ДАНИХ АУДИТУ (AU-12) & a. Забезпечити генерацію даних аудиту для типів подій, що перевіряються в AU-2а в [Призначення: визначених організацією компонентах системи].\newline b. Дозволити [Призначення: визначеному організацією персоналу або посадам] вибирати, які типи подій, що перевіряються, повинні перевірятися окремими компонентами системи;\newline c. Генерувати записи аудиту для типів подій, визначених в AU-2с. з вмістом згідно з AU-3. & AU-12 & KVS забезпечує генерацію записів аудиту для всіх операцій створення, оновлення (через створення нової ревізії) та видалення, гарантуючи фіксацію критичних подій безпеки.\vspace{1mm}\newline\noindent\rule{\linewidth}{0.4pt}\vspace{1mm}
Параметр: audit\_\allowbreak{}generation\newline Тип: boolean\newline Значення: \textbf{true} \\ \hline
\multicolumn{5}{|l|}{\textbf{4. ОЦІНЮВАННЯ, АВТОРИЗАЦІЯ ТА МОНІТОРИНГ (CA)}} \\ \hline
59 & ПОЛІТИКА І ПРОЦЕДУРИ ОЦІНЮВАННЯ, АКРЕДИТАЦІЯ ТА МОНІТОРИНГ БЕЗПЕКИ (CA-1) & a. Розробити, задокументувати та поширити серед [Призначення: визначеного організацією персоналу або посад]:\newline 1. [Вибір (один або декілька): Рівень організації; Рівень місії/бізнес-процесу; рівень системи] політика оцінювання, авторизації та моніторингу, яка:\newline (a) містить мету, сферу застосування, ролі, обов'язки, відповідальність керівництва, координацію між організаційними підрозділами та систему контролю відповідності (complaince);\newline (b) відповідає чинним законам, нормативним документам, наказам, положенням, політиці, стандартам і керівним принципам.\newline 2. Процедури, що сприяють реалізації політики оцінювання, авторизації та моніторингу безпеки та приватності, а також пов'язаних з ними заходів оцінювання, авторизації та моніторингу безпеки та приватності.\newline b. Призначити [Призначення: посадова особа, визначена організацією] для управління розробкою, документуванням і розповсюдженням політики та процедур оцінювання, авторизації та моніторингу;\newline c. Переглядати та оновлювати поточне оцінювання, авторизацію та моніторинг:\newline 1. Політику [Призначення: частота, визначена організацією] та наступне [Призначення: події, визначені організацією];\newline 2. Процедури [Призначення: частота, визначена організацією] та наступні [Призначення: події, визначені організацією]. & CA-1 & Моніторинг та авторизація компонентів ЄСІКС здійснюється на рівні інфраструктури безпеки, з регулярною перевіркою відповідності стандартам (OWASP Top 10, ДСТУ).\vspace{1mm}\newline\noindent\rule{\linewidth}{0.4pt}\vspace{1mm}
Параметр: Розроблено та задокументовано політику оцінювання, авторизації та моніторингу\newline Тип: list\newline Значення: \textbf{default\_\allowbreak{}deny\_\allowbreak{}rule, abac\_\allowbreak{}rule\_\allowbreak{}1}\vspace{1mm}\newline\noindent\rule{\linewidth}{0.4pt}\vspace{1mm}
Параметр: Політика оцінювання, авторизації та моніторингу поширюється серед персоналу або ролей\newline Тип: list\newline Значення: \textbf{admin, security\_\allowbreak{}officer}\vspace{1mm}\newline\noindent\rule{\linewidth}{0.4pt}\vspace{1mm}
Параметр: Розроблені та задокументовані процедури оцінювання, авторизації та моніторингу, що сприяють впровадженню політики оцінювання, авторизації та моніторингу, а також пов'язані з ними засоби контролю оцінювання, авторизації та моніторингу\newline Тип: list\newline Значення: \textbf{default\_\allowbreak{}deny\_\allowbreak{}rule, abac\_\allowbreak{}rule\_\allowbreak{}1}\vspace{1mm}\newline\noindent\rule{\linewidth}{0.4pt}\vspace{1mm}
Параметр: Посадова особа призначається для управління розробкою, документуванням та розповсюдженням політики та процедур оцінювання, авторизації та моніторингу; CA-01(c)[01][01] переглядається та оновлюється поточна політика оцінки, авторизації та моніторингу з частотою; CA-01(c)[01][02] переглядається та оновлюється поточна політика оцінки, авторизації та моніторингу після подій; CA-01(c)[02][01] переглядаються та оновлюються поточні процедури оцінки, авторизації та моніторингу з частотою; CA-01(c)[02][02] переглядаються та оновлюються поточні процедури оцінки, авторизації та моніторингу після подій; ПОТЕНЦІЙНІ МЕТОДИ ТА ОБ'ЄКТИ ОЦІНКИ:\newline Тип: list\newline Значення: \textbf{admin, security\_\allowbreak{}officer} \\ 
 &  &  &  & Параметр: Визначено персонал або ролі, серед яких має бути поширена політика оцінювання, авторизації та моніторингу\newline Тип: list\newline Значення: \textbf{admin, security\_\allowbreak{}officer}\vspace{1mm}\newline\noindent\rule{\linewidth}{0.4pt}\vspace{1mm}
Параметр: Визначено персонал або ролі, серед яких мають бути поширені процедури оцінювання, авторизації та моніторингу\newline Тип: list\newline Значення: \textbf{admin, security\_\allowbreak{}officer}\vspace{1mm}\newline\noindent\rule{\linewidth}{0.4pt}\vspace{1mm}
Параметр: Визначено посадову особу, яка управлятиме розробкою, документуванням і розповсюдженням політики та процедур оцінювання, авторизації та моніторингу\newline Тип: list\newline Значення: \textbf{admin, security\_\allowbreak{}officer}\vspace{1mm}\newline\noindent\rule{\linewidth}{0.4pt}\vspace{1mm}
Параметр: Визначено частоту, з якою переглядається та оновлюється поточна політика оцінювання, авторизації та моніторингу\newline Тип: list\newline Значення: \textbf{default\_\allowbreak{}deny\_\allowbreak{}rule, abac\_\allowbreak{}rule\_\allowbreak{}1} \\ 
 &  &  &  & Параметр: Визначено події, які потребують перегляду та оновлення поточної політики оцінки, авторизації та моніторингу\newline Тип: list\newline Значення: \textbf{default\_\allowbreak{}deny\_\allowbreak{}rule, abac\_\allowbreak{}rule\_\allowbreak{}1}\vspace{1mm}\newline\noindent\rule{\linewidth}{0.4pt}\vspace{1mm}
Параметр: Визначено частоту, з якою переглядається та оновлюється поточні процедури оцінювання, авторизації та моніторингу\newline Тип: integer\newline Значення: \textbf{30}\vspace{1mm}\newline\noindent\rule{\linewidth}{0.4pt}\vspace{1mm}
Параметр: Визначено події, які потребують перегляду та оновлення поточні процедури оцінки, авторизації та моніторингу\newline Тип: list\newline Значення: \textbf{login, logout, failed\_\allowbreak{}attempt} \\ \hline
60 & ОЦІНЮВАННЯ (CA-2) & a. Виберіть відповідного оцінювача або команду з оцінки для типу оцінювання, яке буде проводитися;\newline b. Розробіть план контрольної оцінки, який описує обсяг оцінки, в тому числі:\newline 1. заходи захисту та посилені заходи, що підлягають оцінюванню;\newline 2. процедури оцінювання, ефективності заходів; які використовуватимуться для визначення\newline 3. середовище оцінювання, групу оцінювання, ролі й обов'язки з оцінювання.\newline c. Забезпечити розгляд і затвердження плану оцінювання уповноваженою офіційною особою або призначеним для проведення оцінювання представником;\newline d. Оцінити заходи захисту в системі та в її середовищі функціонування з [Призначення: визначеною організацією частотою] для визначення, наскільки коректно реалізовані заходи безпеки і чи працюють вони за призначенням і дають бажаний результат щодо дотримання встановлених вимог безпеки та приватності;\newline e. Підготовити звіт оцінювання безпеки, який документує результати оцінювання;\newline f. Надати результати оцінювання з безпеки [Призначення: особам або ролям, визначеним організацією]. & CA-2 & Моніторинг та авторизація компонентів ЄСІКС здійснюється на рівні інфраструктури безпеки, з регулярною перевіркою відповідності стандартам (OWASP Top 10, ДСТУ).\vspace{1mm}\newline\noindent\rule{\linewidth}{0.4pt}\vspace{1mm}
Параметр: Розроблено план контрольної оцінки, який описує обсяг оцінки, включаючи заходи захисту та посилені заходи, що підлягають оцінюванню. CA-02(b)[02] розроблено план контрольної оцінки, який описує обсяг оцінки, включаючи процедури оцінювання, які використовуватимуться для визначення ефективності заходів. CA-02(b)[03][01] розроблено план контрольної оцінки, який описує обсяг оцінки, включаючи середовище оцінювання. CA-02(b)[03][02] розроблено план контрольної оцінки, який описує обсяг оцінки, включаючи групу оцінювання. CA-02(b)[03][03] розроблено план контрольної оцінки, який описує обсяг оцінки, включаючи ролі й обов'язки з оцінювання\newline Тип: list\newline Значення: \textbf{admin, security\_\allowbreak{}officer}\vspace{1mm}\newline\noindent\rule{\linewidth}{0.4pt}\vspace{1mm}
Параметр: Розроблено план контрольної оцінки, який описує обсяг оцінки, включаючи процедури оцінювання, які використовуватимуться для визначення ефективності заходів. CA-02(b)[03][01] розроблено план контрольної оцінки, який описує обсяг оцінки, включаючи середовище оцінювання. CA-02(b)[03][02] розроблено план контрольної оцінки, який описує обсяг оцінки, включаючи групу оцінювання. CA-02(b)[03][03] розроблено план контрольної оцінки, який описує обсяг оцінки, включаючи ролі й обов'язки з оцінювання\newline Тип: list\newline Значення: \textbf{admin, security\_\allowbreak{}officer}\vspace{1mm}\newline\noindent\rule{\linewidth}{0.4pt}\vspace{1mm}
Параметр: План оцінки заходів захисту розглядається та затверджується уповноваженою посадовою особою або призначеним представником перед проведенням оцінки\newline Тип: list\newline Значення: \textbf{admin, security\_\allowbreak{}officer}\vspace{1mm}\newline\noindent\rule{\linewidth}{0.4pt}\vspace{1mm}
Параметр: Заходи захисту оцінюються в системі та середовищі її функціонування частота оцінки, щоб визначити, наскільки коректно реалізовані заходи захисту, чи працюють вони за призначенням і чи дають бажаний результат щодо дотримання встановлених вимог до безпеки\newline Тип: string\newline Значення: \textbf{щорічно} \\ 
 &  &  &  & Параметр: Заходи захисту оцінюються в системі та середовищі її функціонування частота оцінки, щоб визначити, наскільки коректно реалізовані заходи захисту, чи працюють вони за призначенням і чи дають бажаний результат щодо дотримання встановлених вимог до конфіденційності; CA-02(e) готується звіт оцінювання\newline Тип: string\newline Значення: \textbf{щорічно}\vspace{1mm}\newline\noindent\rule{\linewidth}{0.4pt}\vspace{1mm}
Параметр: Результати оцінювання з безпеки надаються особам або ролям. оцінювання , який документує результати\newline Тип: list\newline Значення: \textbf{admin, security\_\allowbreak{}officer}\vspace{1mm}\newline\noindent\rule{\linewidth}{0.4pt}\vspace{1mm}
Параметр: Визначено частоту, з якою слід оцінювати засоби контролю в системі та середовищі її функціонування\newline Тип: integer\newline Значення: \textbf{30}\vspace{1mm}\newline\noindent\rule{\linewidth}{0.4pt}\vspace{1mm}
Параметр: Визначені особи або ролі, яким мають бути надані результати оцінювання з безпеки\newline Тип: list\newline Значення: \textbf{admin, security\_\allowbreak{}officer} \\ \hline
61 & ВЗАЄМОДІЯ СИСТЕМ (CA-3) & a. схвалити та керувати обміном інформацією між системою та іншими системами за допомогою [Вибір (один або кілька): угоди безпеки взаємозв'язку; договори безпеки обміну інформацією; меморандуми про взаєморозуміння; угоди про рівень обслуговування; угоди користувача; угоди про нерозголошення; [Доручення: тип договору, визначений організацією]];.\newline b. документувати, як частину угоди про обмін, характеристики інтерфейсу, вимоги до безпеки та приватності, засоби контролю та відповідальність для кожної системи, а також характер переданої інформації;\newline c. здійснювати перегляд та оновлення угод з [Призначення: визначеною організацією частотою]. & CA-3 & Моніторинг та авторизація компонентів ЄСІКС здійснюється на рівні інфраструктури безпеки, з регулярною перевіркою відповідності стандартам (OWASP Top 10, ДСТУ).\vspace{1mm}\newline\noindent\rule{\linewidth}{0.4pt}\vspace{1mm}
Параметр: Характеристики інтерфейсу задокументовані як частина кожної угоди про обмін\newline Тип: integer\newline Значення: \textbf{30}\vspace{1mm}\newline\noindent\rule{\linewidth}{0.4pt}\vspace{1mm}
Параметр: Вимоги до безпеки задокументовані як частина кожної угоди про обмін\newline Тип: integer\newline Значення: \textbf{30}\vspace{1mm}\newline\noindent\rule{\linewidth}{0.4pt}\vspace{1mm}
Параметр: Вимоги щодо конфіденційності задокументовані як частина кожної угоди про обмін\newline Тип: integer\newline Значення: \textbf{30}\vspace{1mm}\newline\noindent\rule{\linewidth}{0.4pt}\vspace{1mm}
Параметр: Заходи захисту задокументовані як частина кожної угоди про обмін\newline Тип: integer\newline Значення: \textbf{30} \\ 
 &  &  &  & Параметр: Відповідальність за кожну систему задокументована як частина кожної угоди про обмін\newline Тип: integer\newline Значення: \textbf{30}\vspace{1mm}\newline\noindent\rule{\linewidth}{0.4pt}\vspace{1mm}
Параметр: Характер переданої інформації документується як частина кожної угоди про обмін\newline Тип: integer\newline Значення: \textbf{30}\vspace{1mm}\newline\noindent\rule{\linewidth}{0.4pt}\vspace{1mm}
Параметр: Угоди переглядаються та оновлюються частота\newline Тип: string\newline Значення: \textbf{щорічно}\vspace{1mm}\newline\noindent\rule{\linewidth}{0.4pt}\vspace{1mm}
Параметр: Визначено частоту, з якою необхідно переглядати та оновлювати угоди\newline Тип: integer\newline Значення: \textbf{30} \\ \hline
62 & ПЛАН УСУНЕННЯ НЕДОЛІКІВ ТА КОНТРОЛЬНІ ПОКАЗНИКИ (CA-5) & a. Розробити для системи план усунення недоліків та контрольні показники з метою документування запланованих коригувальних дій організації для усунення недоліків і зауважень, які виявлені в ході оцінювання заходів захисту, а також для зменшення або усунення відомих вразливостей у системі.\newline b. Оновлювати чинний план усунення недоліків та контрольні показники з [Призначення: визначеною організацією частотою] на основі результатів оцінювання заходів, незалежних аудитів та постійного моніторингу. & CA-5 & Моніторинг та авторизація компонентів ЄСІКС здійснюється на рівні інфраструктури безпеки, з регулярною перевіркою відповідності стандартам (OWASP Top 10, ДСТУ).\vspace{1mm}\newline\noindent\rule{\linewidth}{0.4pt}\vspace{1mm}
Параметр: Розроблено план усунення недоліків та контрольні показники для системи, та задокументувано заплановані дії організації з коригування, спрямовані на усунення недоліків та зауважень, виявлених під час оцінки заходів захисту, а також на зменшення або усунення відомих вразливостей в системі\newline Тип: list\newline Значення: \textbf{login, logout, failed\_\allowbreak{}attempt}\vspace{1mm}\newline\noindent\rule{\linewidth}{0.4pt}\vspace{1mm}
Параметр: Існуючий план оновлюються частота на основі результатів оцінювання заходів, незалежних аудитів та постійного моніторингу\newline Тип: string\newline Значення: \textbf{щорічно}\vspace{1mm}\newline\noindent\rule{\linewidth}{0.4pt}\vspace{1mm}
Параметр: Визначено частоту оновлення чинного плану усунення недоліків та контрольних показників на основі результатів оцінювання заходів захисту, незалежних аудитів та постійного моніторингу\newline Тип: integer\newline Значення: \textbf{30} \\ \hline
63 & БЕЗПЕРЕРВНИЙ МОНІТОРИНГ (CA-7) & Розробити стратегію безперервного моніторингу безпеки та приватності й упровадити програму безперервного моніторингу безпеки та приватності, яка охоплює:\newline a. встановлення показників безпеки та приватності, які необхідно відстежувати: [Призначення: визначені організацією метрики];\newline b. встановлення [Призначення: визначена організацією частота] для моніторингу та [Призначення: визначена організацією частота] для безперервного оцінювання ефективності заходів захисту;\newline c. поточні оцінювання заходів захисту відповідно до стратегії безперервного моніторингу організації;\newline d. постійний моніторинг стану безпеки та приватності відповідно до встановлених організацією метрик і відповідно до стратегії безперервного моніторингу організації;\newline e. зіставлення та аналіз інформації, отриманої в результаті оцінювання та моніторингу безпеки та приватності;\newline f. дії реагування за результатами аналізу інформації, пов'язаної з безпекою та приватністю;\newline g. повідомлення про статус безпеки та приватності системи [Призначення: визначені організацією персонал або ролі] з [Призначення: визначеною організацією частотою]. & CA-7 & Моніторинг та авторизація компонентів ЄСІКС здійснюється на рівні інфраструктури безпеки, з регулярною перевіркою відповідності стандартам (OWASP Top 10, ДСТУ).\vspace{1mm}\newline\noindent\rule{\linewidth}{0.4pt}\vspace{1mm}
Параметр: Безперервний моніторинг на рівні системи включає встановлені частоти для моніторингу ефективності заходів захисту\newline Тип: integer\newline Значення: \textbf{30}\vspace{1mm}\newline\noindent\rule{\linewidth}{0.4pt}\vspace{1mm}
Параметр: Безперервний моніторинг на рівні системи включає встановлені частоти для оцінки ефективності заходів захисту\newline Тип: integer\newline Значення: \textbf{30}\vspace{1mm}\newline\noindent\rule{\linewidth}{0.4pt}\vspace{1mm}
Параметр: Безперервний моніторинг на рівні системи включає в себе дії з реагування на результати аналізу інформації, пов'язаної з безпекою та приватністю\newline Тип: list\newline Значення: \textbf{login, logout, failed\_\allowbreak{}attempt}\vspace{1mm}\newline\noindent\rule{\linewidth}{0.4pt}\vspace{1mm}
Параметр: Безперервний моніторинг на рівні системи включає повідомлення про статус безпеки системи для персоналу або ролей частота\newline Тип: list\newline Значення: \textbf{admin, security\_\allowbreak{}officer} \\ 
 &  &  &  & Параметр: Безперервний моніторинг на рівні системи включає повідомлення про стан конфіденційності системи персоналу або ролям частота\newline Тип: list\newline Значення: \textbf{admin, security\_\allowbreak{}officer}\vspace{1mm}\newline\noindent\rule{\linewidth}{0.4pt}\vspace{1mm}
Параметр: Визначено частоту, з якою слід моніторити ефективність заходів захисту\newline Тип: integer\newline Значення: \textbf{30}\vspace{1mm}\newline\noindent\rule{\linewidth}{0.4pt}\vspace{1mm}
Параметр: Визначено частоту, з якою слід оцінювати ефективність заходів захисту\newline Тип: integer\newline Значення: \textbf{30}\vspace{1mm}\newline\noindent\rule{\linewidth}{0.4pt}\vspace{1mm}
Параметр: Визначено персонал або ролі, яким повідомляється про стан безпеки системи\newline Тип: list\newline Значення: \textbf{admin, security\_\allowbreak{}officer} \\ 
 &  &  &  & Параметр: Визначено частоту, з якою повідомляється про стан безпеки системи\newline Тип: integer\newline Значення: \textbf{30}\vspace{1mm}\newline\noindent\rule{\linewidth}{0.4pt}\vspace{1mm}
Параметр: Визначено персонал або ролі, яким повідомляється про стан конфіденційності системи\newline Тип: list\newline Значення: \textbf{admin, security\_\allowbreak{}officer}\vspace{1mm}\newline\noindent\rule{\linewidth}{0.4pt}\vspace{1mm}
Параметр: Визначено частоту, з якою повідомляється про стан конфіденційності системи\newline Тип: integer\newline Значення: \textbf{30} \\ \hline
\multicolumn{5}{|l|}{\textbf{5. УПРАВЛІННЯ КОНФІГУРАЦІЄЮ (CM)}} \\ \hline
64 & Політика та процедури управління конфігурацією (CM-1) & a. Розробити, задокументувати та поширити серед [Призначення: визначених організацією персоналу або ролей]:\newline 1. [Вибір (один або декілька): Рівень організації; Рівень місії/бізнес-процесу; рівень системи] політики управління конфігурацією, яка:\newline 2.\newline (a) містить мету, сферу застосування, ролі, обов'язки, відповідальність керівництва, координацію між організаційними підрозділами та систему контролю відповідності (complaince);\newline (b) відповідає чинним законам, нормативним документам, наказам, положенням, політикам, стандартам і керівним принципам; процедури, що сприяють реалізації політики управління конфігурацією та пов'язаних з нею заходів управління конфігурацією.\newline b. Призначити [Призначення: посадова особа, визначена організацією] для управління розробкою, документуванням і розповсюдженням політики та процедур керування конфігурацією.\newline c. Переглядати та оновлювати поточну політику управління конфігурацією:\newline 1. Політика [Призначення: частота, визначена організацією] та наступні [Призначення: події, визначені організацією];\newline 2. Процедури [Призначення: частота, визначена організацією] та наступні [Призначення: події, визначені організацією]. & CM-1 & Управління конфігураціями та оновленнями (CM) здійснюється згідно з політиками ІСС. Вразливі компоненти оновлюються протягом 5 днів після виявлення в exploit-db. Конфігураційні файли зберігаються в зашифрованому вигляді.\vspace{1mm}\newline\noindent\rule{\linewidth}{0.4pt}\vspace{1mm}
Параметр: Розроблено та задокументовано політику управління конфігурацією\newline Тип: list\newline Значення: \textbf{default\_\allowbreak{}deny\_\allowbreak{}rule, abac\_\allowbreak{}rule\_\allowbreak{}1}\vspace{1mm}\newline\noindent\rule{\linewidth}{0.4pt}\vspace{1mm}
Параметр: Політика управління конфігурацією поширюється на персонал або ролі\newline Тип: list\newline Значення: \textbf{admin, security\_\allowbreak{}officer}\vspace{1mm}\newline\noindent\rule{\linewidth}{0.4pt}\vspace{1mm}
Параметр: Розроблені та задокументовані процедури управління , що сприяють реалізації політики управління конфігурацією та пов'язаних з нею заходів управління конфігурацією\newline Тип: list\newline Значення: \textbf{default\_\allowbreak{}deny\_\allowbreak{}rule, abac\_\allowbreak{}rule\_\allowbreak{}1}\vspace{1mm}\newline\noindent\rule{\linewidth}{0.4pt}\vspace{1mm}
Параметр: Посадова особа призначається для управління розробкою, документуванням і розповсюдженням політики та процедур керування конфігурацією; CM-01(c)[01][01] переглядається та оновлюється поточна політика управління конфігурацією частота; CM-01(c)[01][02] переглядається та оновлюється поточна політика управління конфігурацією після події; CM-01(c)[02][01] переглядаються та оновлюються поточні процедури управління конфігурацією частота; CM-01(c)[02][02] переглядаються та оновлюються поточні процедури управління конфігурацією після події; ЗНАЧЕННЯ конфігурацією\newline Тип: list\newline Значення: \textbf{admin, security\_\allowbreak{}officer} \\ 
 &  &  &  & Параметр: Визначено персонал або ролі, на яких поширюється політика управління конфігурацією\newline Тип: list\newline Значення: \textbf{admin, security\_\allowbreak{}officer}\vspace{1mm}\newline\noindent\rule{\linewidth}{0.4pt}\vspace{1mm}
Параметр: Визначено персонал або ролі, на яких поширюється процедури управління конфігурацією\newline Тип: list\newline Значення: \textbf{admin, security\_\allowbreak{}officer}\vspace{1mm}\newline\noindent\rule{\linewidth}{0.4pt}\vspace{1mm}
Параметр: Визначено посадову особу, яка управляє розробкою, документуванням і розповсюдженням політики та процедур керування конфігурацією\newline Тип: list\newline Значення: \textbf{admin, security\_\allowbreak{}officer}\vspace{1mm}\newline\noindent\rule{\linewidth}{0.4pt}\vspace{1mm}
Параметр: Визначено частоту, з якою переглядається та оновлюється поточна політика управління конфігурацією\newline Тип: list\newline Значення: \textbf{default\_\allowbreak{}deny\_\allowbreak{}rule, abac\_\allowbreak{}rule\_\allowbreak{}1} \\ 
 &  &  &  & Параметр: Визначено події, після яких переглядається та оновлюється поточна політика управління конфігурацією\newline Тип: list\newline Значення: \textbf{default\_\allowbreak{}deny\_\allowbreak{}rule, abac\_\allowbreak{}rule\_\allowbreak{}1}\vspace{1mm}\newline\noindent\rule{\linewidth}{0.4pt}\vspace{1mm}
Параметр: Визначено частоту, з якою переглядаються та оновлюються поточні процедури управління конфігурацією\newline Тип: integer\newline Значення: \textbf{30}\vspace{1mm}\newline\noindent\rule{\linewidth}{0.4pt}\vspace{1mm}
Параметр: Визначено події, після яких переглядаються та оновлюються поточні процедури управління конфігурацією\newline Тип: list\newline Значення: \textbf{login, logout, failed\_\allowbreak{}attempt} \\ \hline
65 & БАЗОВА КОНФІГУРАЦІЯ (CM-2) & a. Розробити, задокументувати та підтримувати за допомогою заходів конфігурації поточні базові налаштування системи.\newline b. Переглядати та оновлювати базові налаштування системи:\newline 1. з [Призначення: визначеною організацією частотою];\newline 2. за потреби внаслідок [Призначення: визначених організацією обставин];\newline 3. коли встановлені нові або оновлені компоненти системи. & CM-2 & Файл `erpuno.ex` фіксує єдину базову конфігурацію (Baseline) безпеки для всієї системи. Зміни в конфігурації застосовуються виключно через оновлення коду, що забезпечує строгий контроль та аудит змін.\vspace{1mm}\newline\noindent\rule{\linewidth}{0.4pt}\vspace{1mm}
Параметр: baseline\_\allowbreak{}configuration\_\allowbreak{}code\newline Тип: boolean\newline Значення: \textbf{true} \\ \hline
66 & БАЗОВА КОНФІГУРАЦІЯ - КОНФІГУРАЦІЯ СИСТЕМ ТА КОМПОНЕНТІВ ДЛЯ СФЕР З ВИСОКИМ РИЗИКОМ (CM-2(7)) &  & CM-2(7) & Управління конфігураціями та оновленнями (CM) здійснюється згідно з політиками ІСС. Вразливі компоненти оновлюються протягом 5 днів після виявлення в exploit-db. Конфігураційні файли зберігаються в зашифрованому вигляді.\vspace{1mm}\newline\noindent\rule{\linewidth}{0.4pt}\vspace{1mm}
Параметр: Системи або компоненти системи з конфігураціями видаються особам, що перебувають у місцях, які організація вважає, становлять значний ризик\newline Тип: list\newline Значення: \textbf{admin, security\_\allowbreak{}officer}\vspace{1mm}\newline\noindent\rule{\linewidth}{0.4pt}\vspace{1mm}
Параметр: Заходи безпеки застосовуються до систем або компонентів системи, коли особи повертаються з поїздки\newline Тип: list\newline Значення: \textbf{admin, security\_\allowbreak{}officer}\vspace{1mm}\newline\noindent\rule{\linewidth}{0.4pt}\vspace{1mm}
Параметр: Визначено системи або компоненти систем, які мають видаватися особам, що перебувають у місцях зі значним ризиком\newline Тип: list\newline Значення: \textbf{admin, security\_\allowbreak{}officer} \\ \hline
67 & УПРАВЛІННЯ ЗМІНАМИ КОНФІГУРАЦІЇ (CM-3) & a. Визначити типи змін у системі, які контролюються конфігурацією.\newline b. Переглядати запропоновані зміни в конфігурації, контрольовані системою, і схвалити або відхилити ці зміни з явним урахуванням аналізу наслідків безпеки.\newline c. Документувати рішення зі зміни конфігурації системи.\newline d. Впровадити схвалені зміни конфігурації в систему.\newline e. Зберігати записи змін конфігурації системі впродовж [Призначення: певного періоду часу, визначеного організацією].\newline f. Здійснювати моніторинг і аналіз дій, пов'язаних зі змінами конфігурації системи.\newline g. Координувати та впроваджувати нагляд за діяльністю з управління змінами конфігурації за допомогою [Призначення: елементу управління змінами конфігурації, визначеного організацією], який викликається [Вибір (один або кілька): [Призначення: з визначеною організацією частотою]; [Призначення: визначені організацією умови зміни конфігурації]]. & CM-3 & Оскільки CMDB є частиною вихідного коду (Elixir модуль), всі зміни проходять через стандартний процес розробки (Code Review, CI/CD, контроль версій Git), унеможливлюючи несанкціоноване втручання в налаштування на \textquotedbl{}живій\textquotedbl{} системі.\vspace{1mm}\newline\noindent\rule{\linewidth}{0.4pt}\vspace{1mm}
Параметр: configuration\_\allowbreak{}change\_\allowbreak{}control\newline Тип: boolean\newline Значення: \textbf{true} \\ \hline
68 & АНАЛІЗ ВПЛИВУ НА БЕЗПЕКУ ТА ПРИВАТНІСТЬ (CM-4) & Аналізувати зміни в системі, щоб визначити потенційну загрозу безпеці та приватності перед реалізацією змін. & CM-4 & Управління конфігураціями та оновленнями (CM) здійснюється згідно з політиками ІСС. Вразливі компоненти оновлюються протягом 5 днів після виявлення в exploit-db. Конфігураційні файли зберігаються в зашифрованому вигляді. \\ \hline
69 & АНАЛІЗ ВПЛИВУ НА БЕЗПЕКУ ТА ПРИВАТНІСТЬ - ВЕРИФІКАЦІЯ ФУНКЦІЙ БЕЗПЕКИ ТА ПРИВАТНОСТІ (CM-4(2)) &  & CM-4(2) & Управління конфігураціями та оновленнями (CM) здійснюється згідно з політиками ІСС. Вразливі компоненти оновлюються протягом 5 днів після виявлення в exploit-db. Конфігураційні файли зберігаються в зашифрованому вигляді. \\ \hline
70 & ОБМЕЖЕННЯ ДОСТУПУ ДО ЗМІНИ (CM-5) & Визначити, задокументувати, затвердити та забезпечити застосування фізичних і логічних обмежень доступу, пов'язаних зі змінами в системі. & CM-5 & Управління конфігураціями та оновленнями (CM) здійснюється згідно з політиками ІСС. Вразливі компоненти оновлюються протягом 5 днів після виявлення в exploit-db. Конфігураційні файли зберігаються в зашифрованому вигляді. \\ \hline
71 & НАЛАШТУВАННЯ КОНФІГУРАЦІЇ (CM-6) & a. Встановити та задокументувати параметри конфігурації компонентів, які застосовуються в системі, які відображають найбільш обмежений режим, що відповідає експлуатаційним вимогам, використовуючи [Призначення: визначені організацією загальні безпечні конфігурації].\newline b. Реалізувати конфігураційні установки.\newline c. Визначити, задокументувати та затвердити будь-які відхилення від встановлених конфігураційних параметрів конфігурації для [Призначення: визначених організацією компонентів системи] на основі [Призначення: визначених організацією експлуатаційних вимог].\newline d. Відстежувати та керувати змінами конфігураційних параметрів відповідно до організаційної політики та процедур. & CM-6 & Параметри політик безпеки та їхні значення за замовчуванням жорстко зашиті в `erpuno.ex`. Компілятор генерує незмінний об'єктний файл, який гарантує виконання затверджених налаштувань під час роботи.\vspace{1mm}\newline\noindent\rule{\linewidth}{0.4pt}\vspace{1mm}
Параметр: compile\_\allowbreak{}time\_\allowbreak{}config\newline Тип: boolean\newline Значення: \textbf{true} \\ \hline
72 & МІНІМАЛЬНО НЕОБХІДНА ФУНКЦІОНАЛЬНІСТЬ (CM-7) & a. Налаштуйте систему для забезпечення лише [Призначення: основні функції, визначені організацією для місії];\newline b. Заборонити або обмежити використання таких функцій, портів, протоколів, програмного забезпечення та/або служб: [Призначення: визначені організацією заборонені або обмежені функції, системні порти, протоколи, програмне забезпечення та/або служби]. & CM-7 & Управління конфігураціями та оновленнями (CM) здійснюється згідно з політиками ІСС. Вразливі компоненти оновлюються протягом 5 днів після виявлення в exploit-db. Конфігураційні файли зберігаються в зашифрованому вигляді.\vspace{1mm}\newline\noindent\rule{\linewidth}{0.4pt}\vspace{1mm}
Параметр: Використання протоколи заборонено або обмежено\newline Тип: string\newline Значення: \textbf{TLS 1.3}\vspace{1mm}\newline\noindent\rule{\linewidth}{0.4pt}\vspace{1mm}
Параметр: Визначено протоколи, які необхідно заборонити або обмежити; CM-07\_\allowbreak{}ODP[05] визначено програмне забезпечення, яке необхідно заборонити або обмежити\newline Тип: string\newline Значення: \textbf{TLS 1.3} \\ \hline
73 & МІНІМАЛЬНО НЕОБХІДНА ФУНКЦІОНАЛЬНІСТЬ - ПЕРІОДИЧНИЙ ПЕРЕГЛЯД (CM-7(1)) &  & CM-7(1) & Управління конфігураціями та оновленнями (CM) здійснюється згідно з політиками ІСС. Вразливі компоненти оновлюються протягом 5 днів після виявлення в exploit-db. Конфігураційні файли зберігаються в зашифрованому вигляді.\vspace{1mm}\newline\noindent\rule{\linewidth}{0.4pt}\vspace{1mm}
Параметр: Система переглядається частота для виявлення непотрібних та/або незахищених функцій, портів, протоколів і послуг\newline Тип: string\newline Значення: \textbf{TLS 1.3}\vspace{1mm}\newline\noindent\rule{\linewidth}{0.4pt}\vspace{1mm}
Параметр: Протоколи, які вважаються непотрібними та/або незахищеними, вимкнено\newline Тип: string\newline Значення: \textbf{TLS 1.3}\vspace{1mm}\newline\noindent\rule{\linewidth}{0.4pt}\vspace{1mm}
Параметр: Визначено частоту, з якою слід переглядати систему для виявлення непотрібних та/або незахищених функцій, портів, протоколів і послуг\newline Тип: string\newline Значення: \textbf{TLS 1.3}\vspace{1mm}\newline\noindent\rule{\linewidth}{0.4pt}\vspace{1mm}
Параметр: Визначено протоколи, які слід вимкнути, якщо вони вважаються непотрібними або незахищеними\newline Тип: string\newline Значення: \textbf{TLS 1.3} \\ \hline
74 & МІНІМАЛЬНО НЕОБХІДНА ФУНКЦІОНАЛЬНІСТЬ - АВТОРИЗОВАНЕ ПРОГРАМНЕ ЗАБЕЗПЕЧЕННЯ -- БІЛИЙ СПИСОК (CM-7(5)) &  & CM-7(5) & Управління конфігураціями та оновленнями (CM) здійснюється згідно з політиками ІСС. Вразливі компоненти оновлюються протягом 5 днів після виявлення в exploit-db. Конфігураційні файли зберігаються в зашифрованому вигляді.\vspace{1mm}\newline\noindent\rule{\linewidth}{0.4pt}\vspace{1mm}
Параметр: Політика \textquotedbl{}заборонити все, дозволити за винятком\textquotedbl{} застосовуэться, щоб дозволити виконання авторизованих програм у системі\newline Тип: list\newline Значення: \textbf{default\_\allowbreak{}deny\_\allowbreak{}rule, abac\_\allowbreak{}rule\_\allowbreak{}1}\vspace{1mm}\newline\noindent\rule{\linewidth}{0.4pt}\vspace{1mm}
Параметр: Переглядається та оновлюється список авторизованих програм частота\newline Тип: list\newline Значення: \textbf{}\vspace{1mm}\newline\noindent\rule{\linewidth}{0.4pt}\vspace{1mm}
Параметр: Визначено частоту перегляду та оновлення списку авторизованих програм\newline Тип: integer\newline Значення: \textbf{30} \\ \hline
75 & ІНВЕНТАРИЗАЦІЯ КОМПОНЕНТІВ СИСТЕМИ (CM-8) & a.\newline b. Розробити та задокументувати процес інвентаризації компонентів системи, який:\newline 1. точно описує поточну систему;\newline 2. охоплює всі компоненти в межах акредитації системи;\newline 3. не включає повторний облік компонентів або компонентів, будь-якої іншої системи;\newline 4. визначає рівень деталізації, який є необхідним для відстеження та звітування;\newline 5. включає інформацію для досягнення підзвітності компонентів системи: [Призначення: визначена організацією інформація, необхідна для досягнення ефективної підзвітності компонентів системи]. Переглядати та оновлювати опис компонентів системи з [Призначення: визначеною організацією частотою]. & CM-8 & Профіль автоматично фіксує перелік використовуваних компонентів системи (бібліотеки `synrc/kvs`, `synrc/ca` тощо) та їхні параметри, забезпечуючи інтегральну інвентаризацію модулів захисту.\vspace{1mm}\newline\noindent\rule{\linewidth}{0.4pt}\vspace{1mm}
Параметр: component\_\allowbreak{}inventory\newline Тип: boolean\newline Значення: \textbf{true} \\ \hline
76 & ІНВЕНТАРИЗАЦІЯ КОМПОНЕНТІВ СИСТЕМИ - ОНОВЛЕННЯ ПІД ЧАС ВСТАНОВЛЕННЯ ТА ВИДАЛЕННЯ (CM-8(1)) &  & CM-8(1) & Управління конфігураціями та оновленнями (CM) здійснюється згідно з політиками ІСС. Вразливі компоненти оновлюються протягом 5 днів після виявлення в exploit-db. Конфігураційні файли зберігаються в зашифрованому вигляді. \\ \hline
77 & РОЗТАШУВАННЯ ІНФОРМАЦІЇ (CM-12) &  & CM-12 & Управління конфігураціями та оновленнями (CM) здійснюється згідно з політиками ІСС. Вразливі компоненти оновлюються протягом 5 днів після виявлення в exploit-db. Конфігураційні файли зберігаються в зашифрованому вигляді. \\ \hline
\multicolumn{5}{|l|}{\textbf{6. ПЛАНУВАННЯ БЕЗПЕРЕРВНОЇ РОБОТИ (CP)}} \\ \hline
78 & РЕЗЕРВНЕ КОПІЮВАННЯ (CP-9) &  & CP-9 & Система здійснює автоматичне створення резервних копій бази даних KVS. Бекапи можуть зберігатися локально або відправлятися до хмарних сховищ, гарантуючи відновлення даних відповідно до RTO/RPO.\vspace{1mm}\newline\noindent\rule{\linewidth}{0.4pt}\vspace{1mm}
Параметр: system\_\allowbreak{}backup\newline Тип: boolean\newline Значення: \textbf{true} \\ \hline
79 & РЕЗЕРВНЕ КОПІЮВАННЯ - КРИПТОГРАФІЧНИЙ ЗАХИСТ (CP-9(8)) &  & CP-9(8) & Безперервність роботи та резервне копіювання забезпечується кластеризацією сховищ даних (LSM-дерева, реплікація PostgreSQL) та регулярним копіюванням БД і журналів згідно з політиками ЄСІКС.\vspace{1mm}\newline\noindent\rule{\linewidth}{0.4pt}\vspace{1mm}
Параметр: Реалізовано криптографічні механізми для запобігання несанкціонованому розкриттю та зміні резервної інформації\newline Тип: string\newline Значення: \textbf{AES-256-GCM} \\ \hline
\multicolumn{5}{|l|}{\textbf{7. ІДЕНТИФІКАЦІЯ ТА АВТЕНТИФІКАЦІЯ (IA)}} \\ \hline
80 & Політика та процедури ідентифікації та автентифікації (IA-1) &  & IA-1 & Ідентифікація та автентифікація в ЄСІКС вимагає обов'язкового використання КЕП на захищеному носії (ДСТУ 4145:2002) з форматом CAdES-X Long та SSO інтеграцією.\vspace{1mm}\newline\noindent\rule{\linewidth}{0.4pt}\vspace{1mm}
Параметр: Розроблено та задокументовано політику ідентифікації та автентифікації\newline Тип: list\newline Значення: \textbf{default\_\allowbreak{}deny\_\allowbreak{}rule, abac\_\allowbreak{}rule\_\allowbreak{}1}\vspace{1mm}\newline\noindent\rule{\linewidth}{0.4pt}\vspace{1mm}
Параметр: Політика ідентифікації та автентифікації поширюється на персонал або ролі\newline Тип: list\newline Значення: \textbf{admin, security\_\allowbreak{}officer}\vspace{1mm}\newline\noindent\rule{\linewidth}{0.4pt}\vspace{1mm}
Параметр: Розроблені та задокументовані процедури ідентифікації та автентифікації, що сприяють впровадженню політики ідентифікації та автентифікації, а також відповідні заходи ідентифікації та перевірки автентичності\newline Тип: list\newline Значення: \textbf{default\_\allowbreak{}deny\_\allowbreak{}rule, abac\_\allowbreak{}rule\_\allowbreak{}1}\vspace{1mm}\newline\noindent\rule{\linewidth}{0.4pt}\vspace{1mm}
Параметр: Посадова особа призначається для управління політикою та процедурами ідентифікації та автентифікації; IA-01(c)[01][01] переглядається та оновлюється поточна політика ідентифікації та автентифікації частота; IA-01(c)[01][02] переглядається та оновлюється поточна політика ідентифікації та автентифікації після подій; IA-01(c)[02][01] переглядаються та оновлюються поточні процедури ідентифікації та автентифікації частота; IA-01(c)[02][02] переглядаються та оновлюються поточні процедури ідентифікації та автентифікації після подій\newline Тип: list\newline Значення: \textbf{admin, security\_\allowbreak{}officer} \\ 
 &  &  &  & Параметр: Визначено персонал або ролі, на які поширюється політика ідентифікації та автентифікації\newline Тип: list\newline Значення: \textbf{admin, security\_\allowbreak{}officer}\vspace{1mm}\newline\noindent\rule{\linewidth}{0.4pt}\vspace{1mm}
Параметр: Визначено персонал або ролі, на які поширюються процедури ідентифікації та автентифікації\newline Тип: list\newline Значення: \textbf{admin, security\_\allowbreak{}officer}\vspace{1mm}\newline\noindent\rule{\linewidth}{0.4pt}\vspace{1mm}
Параметр: Визначено посадову особу, яка управляє політикою та процедурами ідентифікації та автентифікації\newline Тип: list\newline Значення: \textbf{admin, security\_\allowbreak{}officer}\vspace{1mm}\newline\noindent\rule{\linewidth}{0.4pt}\vspace{1mm}
Параметр: Визначено частоту, з якою переглядається та оновлюється поточна політика ідентифікації та автентифікації\newline Тип: list\newline Значення: \textbf{default\_\allowbreak{}deny\_\allowbreak{}rule, abac\_\allowbreak{}rule\_\allowbreak{}1} \\ 
 &  &  &  & Параметр: Визначено події, які потребують перегляду та оновлення поточної політики ідентифікації та автентифікації\newline Тип: list\newline Значення: \textbf{default\_\allowbreak{}deny\_\allowbreak{}rule, abac\_\allowbreak{}rule\_\allowbreak{}1}\vspace{1mm}\newline\noindent\rule{\linewidth}{0.4pt}\vspace{1mm}
Параметр: Визначено частоту, з якою переглядаються та оновлюються поточні процедури ідентифікації та автентифікації\newline Тип: integer\newline Значення: \textbf{30}\vspace{1mm}\newline\noindent\rule{\linewidth}{0.4pt}\vspace{1mm}
Параметр: Визначено події, які потребують перегляду та оновлення процедур ідентифікації та автентифікації\newline Тип: list\newline Значення: \textbf{login, logout, failed\_\allowbreak{}attempt} \\ \hline
81 & ІДЕНТИФІКАЦІЯ ТА АВТЕНТИФІКАЦІЯ КОРИСТУВАЧІВ (IA-2) &  & IA-2 & Ідентифікація та автентифікація в ЄСІКС вимагає обов'язкового використання КЕП на захищеному носії (ДСТУ 4145:2002) з форматом CAdES-X Long та SSO інтеграцією.\vspace{1mm}\newline\noindent\rule{\linewidth}{0.4pt}\vspace{1mm}
Параметр: org\_\allowbreak{}user\_\allowbreak{}auth\newline Тип: boolean\newline Значення: \textbf{true} \\ \hline
82 & ІДЕНТИФІКАЦІЯ ТА АВТЕНТИФІКАЦІЯ (КОРИСТУВАЧІВ ОРГАНІЗАЦІЇ) - БАГАТОФАКТОРНА АВТЕНТИФІКАЦІЯ ПРИВІЛЕЙОВАНИХ ОБЛІКОВИХ ЗАПИСІВ (IA-2(1)) &  & IA-2(1) & Ідентифікація та автентифікація в ЄСІКС вимагає обов'язкового використання КЕП на захищеному носії (ДСТУ 4145:2002) з форматом CAdES-X Long та SSO інтеграцією. \\ \hline
83 & ІДЕНТИФІКАЦІЯ ТА АВТЕНТИФІКАЦІЯ (КОРИСТУВАЧІВ ОРГАНІЗАЦІЇ) - БАГАТОФАКТОРНА АВТЕНТИФІКАЦІЯ НЕПРИВІЛЕЙОВАНИХ (IA-2(2)) &  & IA-2(2) & Ідентифікація та автентифікація в ЄСІКС вимагає обов'язкового використання КЕП на захищеному носії (ДСТУ 4145:2002) з форматом CAdES-X Long та SSO інтеграцією. \\ \hline
84 & Ідентифікація та автентифікація (користувачів організації) - Доступ до облікових записів --- стійкість до відтворення (IA-2(8)) &  & IA-2(8) & Ідентифікація та автентифікація в ЄСІКС вимагає обов'язкового використання КЕП на захищеному носії (ДСТУ 4145:2002) з форматом CAdES-X Long та SSO інтеграцією. \\ \hline
85 & ІДЕНТИФІКАЦІЯ ТА АВТЕНТИФІКАЦІЯ (КОРИСТУВАЧІВ ОРГАНІЗАЦІЇ) - ПРИЙНЯТТЯ ПОВНОВАЖЕНЬ ДЛЯ ВЕРИФІКАЦІЇ ОСОБИСТОЇ ІНФОРМАЦІЇ (PIV CREDENTIALS) (IA-2(12)) &  & IA-2(12) & Ідентифікація та автентифікація в ЄСІКС вимагає обов'язкового використання КЕП на захищеному носії (ДСТУ 4145:2002) з форматом CAdES-X Long та SSO інтеграцією.\vspace{1mm}\newline\noindent\rule{\linewidth}{0.4pt}\vspace{1mm}
Параметр: Приймаються та електронним шляхом підтверджуються повноваження облікових даних особистої ідентифікації\newline Тип: list\newline Значення: \textbf{admin, security\_\allowbreak{}officer} \\ \hline
86 & ІДЕНТИФІКАЦІЯ ТА АВТЕНТИФІКАЦІЯ ПРИСТРОЇВ (IA-3) &  & IA-3 & Ідентифікація та автентифікація в ЄСІКС вимагає обов'язкового використання КЕП на захищеному носії (ДСТУ 4145:2002) з форматом CAdES-X Long та SSO інтеграцією. \\ \hline
87 & УПРАВЛІННЯ ІДЕНТИФІКАЦІЄЮ (IA-4) &  & IA-4 & Ідентифікація та автентифікація в ЄСІКС вимагає обов'язкового використання КЕП на захищеному носії (ДСТУ 4145:2002) з форматом CAdES-X Long та SSO інтеграцією.\vspace{1mm}\newline\noindent\rule{\linewidth}{0.4pt}\vspace{1mm}
Параметр: Управління ідентифікаторами здійснюється шляхом отримання дозволу від персоналу або ролей на призначення ідентифікатора особі, групі, ролі або пристрою\newline Тип: list\newline Значення: \textbf{admin, security\_\allowbreak{}officer}\vspace{1mm}\newline\noindent\rule{\linewidth}{0.4pt}\vspace{1mm}
Параметр: Управління ідентифікаторами здійснюється шляхом вибору ідентифікатора, який ідентифікує окрему особу, групу, ролі або пристрій\newline Тип: list\newline Значення: \textbf{admin, security\_\allowbreak{}officer}\vspace{1mm}\newline\noindent\rule{\linewidth}{0.4pt}\vspace{1mm}
Параметр: Управління ідентифікаторами здійснюється шляхом призначення ідентифікатора особі, групі, ролі або пристрою; IA-04(d) ідентифікатори управляються шляхом запобігання повторному використанню ідентифікаторів впродовж період часу\newline Тип: list\newline Значення: \textbf{admin, security\_\allowbreak{}officer}\vspace{1mm}\newline\noindent\rule{\linewidth}{0.4pt}\vspace{1mm}
Параметр: Визначено персонал або ролі, від яких необхідно отримати дозвіл на призначення ідентифікатора\newline Тип: list\newline Значення: \textbf{admin, security\_\allowbreak{}officer} \\ 
 &  &  &  & Параметр: Визначено період часу для запобігання повторному використанню ідентифікаторів\newline Тип: integer\newline Значення: \textbf{30} \\ \hline
88 & УПРАВЛІННЯ ІДЕНТИФІКАЦІЄЮ - ІДЕНТИФІКАЦІЯ СТАТУСУ КОРИСТУВАЧА (IA-4(4)) &  & IA-4(4) & Ідентифікація та автентифікація в ЄСІКС вимагає обов'язкового використання КЕП на захищеному носії (ДСТУ 4145:2002) з форматом CAdES-X Long та SSO інтеграцією. \\ \hline
89 & АВТЕНТИФІКАТОР УПРАВЛІННЯ &  & IA-5 & Ідентифікація та автентифікація в ЄСІКС вимагає обов'язкового використання КЕП на захищеному носії (ДСТУ 4145:2002) з форматом CAdES-X Long та SSO інтеграцією.\vspace{1mm}\newline\noindent\rule{\linewidth}{0.4pt}\vspace{1mm}
Параметр: Управління системними автентифікаторами здійснюється шляхом перевірки, як частини початкового розподілу автентифікатора, особи, групи, ролі або пристрою, який отримує автентифікатор\newline Тип: list\newline Значення: \textbf{admin, security\_\allowbreak{}officer}\vspace{1mm}\newline\noindent\rule{\linewidth}{0.4pt}\vspace{1mm}
Параметр: Управління системними автентифікаторами здійснюється шляхом зміни/оновлення автентифікаторів у встановлений період часу або коли відбуваються події\newline Тип: list\newline Значення: \textbf{login, logout, failed\_\allowbreak{}attempt}\vspace{1mm}\newline\noindent\rule{\linewidth}{0.4pt}\vspace{1mm}
Параметр: Визначено період часу для зміни або оновлення автентифікаторів за типом автентифікатора\newline Тип: integer\newline Значення: \textbf{30} \\ \hline
90 & УПРАВЛІННЯ АВТЕНТИФІКАТОРОМ - АВТЕНТИФІКАЦІЯ НА ОСНОВІ ПАРОЛЯ (IA-5(1)) &  & IA-5(1) & Ідентифікація та автентифікація в ЄСІКС вимагає обов'язкового використання КЕП на захищеному носії (ДСТУ 4145:2002) з форматом CAdES-X Long та SSO інтеграцією.\vspace{1mm}\newline\noindent\rule{\linewidth}{0.4pt}\vspace{1mm}
Параметр: Для автентифікації на основі паролів підтримується та оновлюється список часто використовуваних, очікуваних або скомпрометованих паролів частота, а також коли є підозра, що паролі організації були скомпрометовані прямо чи опосередковано\newline Тип: list\newline Значення: \textbf{admin, security\_\allowbreak{}officer}\vspace{1mm}\newline\noindent\rule{\linewidth}{0.4pt}\vspace{1mm}
Параметр: Для автентифікації на основі паролів, коли паролі створюються або оновлюються користувачами, паролі перевіряються на відсутність у списку загальновживаних, очікуваних або скомпрометованих паролів в IA-05(01)(a)\newline Тип: list\newline Значення: \textbf{admin, security\_\allowbreak{}officer}\vspace{1mm}\newline\noindent\rule{\linewidth}{0.4pt}\vspace{1mm}
Параметр: Для автентифікації на основі паролів, паролі передаються лише криптографічно захищеними каналами\newline Тип: list\newline Значення: \textbf{admin, security\_\allowbreak{}officer}\vspace{1mm}\newline\noindent\rule{\linewidth}{0.4pt}\vspace{1mm}
Параметр: Для автентифікації на основі паролів паролі зберігаються за допомогою затвердженого алгоритму гешування, переважно використовуючи ключову геш-функцію\newline Тип: list\newline Значення: \textbf{admin, security\_\allowbreak{}officer} \\ 
 &  &  &  & Параметр: Для автентифікації на основі пароля дозволяється вибір користувачем довгих паролів і фраз, що включають пробіли та всі друковані символи\newline Тип: list\newline Значення: \textbf{admin, security\_\allowbreak{}officer}\vspace{1mm}\newline\noindent\rule{\linewidth}{0.4pt}\vspace{1mm}
Параметр: Для автентифікації на основі пароля використовуються автоматизовані інструменти, які допомагають користувачеві у виборі надійних автентифікаторів паролів\newline Тип: list\newline Значення: \textbf{admin, security\_\allowbreak{}officer}\vspace{1mm}\newline\noindent\rule{\linewidth}{0.4pt}\vspace{1mm}
Параметр: Для автентифікації на основі пароля застосовуються склад та правила складності\newline Тип: list\newline Значення: \textbf{default\_\allowbreak{}deny\_\allowbreak{}rule, abac\_\allowbreak{}rule\_\allowbreak{}1}\vspace{1mm}\newline\noindent\rule{\linewidth}{0.4pt}\vspace{1mm}
Параметр: Визначено частоту оновлення списку часто використовуваних, очікуваних або скомпрометованих паролів\newline Тип: list\newline Значення: \textbf{admin, security\_\allowbreak{}officer} \\ 
 &  &  &  & Параметр: Визначено склад та правила складності автентифікатора\newline Тип: list\newline Значення: \textbf{default\_\allowbreak{}deny\_\allowbreak{}rule, abac\_\allowbreak{}rule\_\allowbreak{}1} \\ \hline
91 & УПРАВЛІННЯ АВТЕНТИФІКАТОРОМ - АВТЕНТИФІКАЦІЯ НА ОСНОВІ ВІДКРИТОГО КЛЮЧА (IA-5(2)) &  & IA-5(2) & Ідентифікація та автентифікація в ЄСІКС вимагає обов'язкового використання КЕП на захищеному носії (ДСТУ 4145:2002) з форматом CAdES-X Long та SSO інтеграцією. \\ \hline
92 & УПРАВЛІННЯ АВТЕНТИФІКАТОРОМ - АВТЕНТИФІКАЦІЯ НА ОСНОВІ АПАРАТНИХ ТОКЕНІВ (IA-5(11)) &  & IA-5(11) & Ідентифікація та автентифікація в ЄСІКС вимагає обов'язкового використання КЕП на захищеному носії (ДСТУ 4145:2002) з форматом CAdES-X Long та SSO інтеграцією. \\ \hline
93 & УПРАВЛІННЯ АВТЕНТИФІКАТОРОМ - ЕФЕКТИВНІСТЬ БІОМЕТРИЧНОЇ АВТЕНТИФІКАЦІЇ (IA-5(12)) &  & IA-5(12) & Ідентифікація та автентифікація в ЄСІКС вимагає обов'язкового використання КЕП на захищеному носії (ДСТУ 4145:2002) з форматом CAdES-X Long та SSO інтеграцією.\vspace{1mm}\newline\noindent\rule{\linewidth}{0.4pt}\vspace{1mm}
Параметр: Визначено вимоги до якості біометрії; для біометричної автентифікації використовувати механізми, які задовольняють вимоги\newline Тип: string\newline Значення: \textbf{автоматизований засіб моніторингу} \\ \hline
94 & УПРАВЛІННЯ АВТЕНТИФІКАТОРОМ - ПРОДУКТИ ТА ПОСЛУГИ, ЗАТВЕРДЖЕНІ УПОВНОВАЖЕНИМ ОРГАНОМ (IA-5(15)) &  & IA-5(15) & Ідентифікація та автентифікація в ЄСІКС вимагає обов'язкового використання КЕП на захищеному носії (ДСТУ 4145:2002) з форматом CAdES-X Long та SSO інтеграцією. \\ \hline
95 & УПРАВЛІННЯ АВТЕНТИФІКАТОРОМ - АВТОМАТИЗОВАНІ ЗАСОБИ ВИЯВЛЕННЯ АТАК ІЗ ВИКОРИСТАННЯМ БІОМЕТРИЧНИХ АВТЕНТИФІКАТОРІВ (IA-5(17)) &  & IA-5(17) & Ідентифікація та автентифікація в ЄСІКС вимагає обов'язкового використання КЕП на захищеному носії (ДСТУ 4145:2002) з форматом CAdES-X Long та SSO інтеграцією.\vspace{1mm}\newline\noindent\rule{\linewidth}{0.4pt}\vspace{1mm}
Параметр: Використовуються механізми виявлення атак із використанням штучно виготовлених артефактів для біометричних автентифікаторів\newline Тип: string\newline Значення: \textbf{автоматизований засіб моніторингу} \\ \hline
96 & ПРИХОВУВАННЯ ЗВОРОТНОГО ЗВ'ЯЗКУ АВТЕНТИФІКАТОРА &  & IA-6 & Ідентифікація та автентифікація в ЄСІКС вимагає обов'язкового використання КЕП на захищеному носії (ДСТУ 4145:2002) з форматом CAdES-X Long та SSO інтеграцією.\vspace{1mm}\newline\noindent\rule{\linewidth}{0.4pt}\vspace{1mm}
Параметр: Забезпечино приховану зворотну передачу інформації автентифікації в про- цесі автентифікації для забезпечення захисту інформації від можливої експлуатації та використання неавторизованими особами\newline Тип: list\newline Значення: \textbf{admin, security\_\allowbreak{}officer} \\ \hline
97 & АВТЕНТИФІКАЦІЯ КРИПТОГРАФІЧНОГО МОДУЛЯ (IA-7) &  & IA-7 & Ідентифікація та автентифікація в ЄСІКС вимагає обов'язкового використання КЕП на захищеному носії (ДСТУ 4145:2002) з форматом CAdES-X Long та SSO інтеграцією.\vspace{1mm}\newline\noindent\rule{\linewidth}{0.4pt}\vspace{1mm}
Параметр: Впроваджено механізми автентифікації в криптографічний модуль, який відповідає вимогам чинних законів, виконавчих розпоряджень, директив, політик, правил, стандартів та рекомендацій для такої автентифікації\newline Тип: list\newline Значення: \textbf{default\_\allowbreak{}deny\_\allowbreak{}rule, abac\_\allowbreak{}rule\_\allowbreak{}1} \\ \hline
98 & ПОВТОРНА АВТЕНТИФІКАЦІЯ (IA-11) &  & IA-11 & Ідентифікація та автентифікація в ЄСІКС вимагає обов'язкового використання КЕП на захищеному носії (ДСТУ 4145:2002) з форматом CAdES-X Long та SSO інтеграцією. \\ \hline
\multicolumn{5}{|l|}{\textbf{8. РЕАГУВАННЯ НА ІНЦИДЕНТИ (IR)}} \\ \hline
99 & Політика та процедури реагування на інциденти (IR-1) &  & IR-1 & Політика реагування на інциденти закодована безпосередньо у бізнес-процесах (BPMN) `erpuno/itsm`, забезпечуючи неухильне виконання процедур на кожному кроці життєвого циклу інциденту.\vspace{1mm}\newline\noindent\rule{\linewidth}{0.4pt}\vspace{1mm}
Параметр: incident\_\allowbreak{}policy\_\allowbreak{}enforcement\newline Тип: boolean\newline Значення: \textbf{true} \\ \hline
100 & НАВЧАННЯ З РЕАГУВАННЯ НА ІНЦИДЕНТИ (IR-2) &  & IR-2 & Інтерфейс `erpuno/itsm` надає контекстні підказки та інструкції безпосередньо у формах обробки інцидентів, що виконує функцію інтерактивного навчання та підтримки персоналу в кризових ситуаціях.\vspace{1mm}\newline\noindent\rule{\linewidth}{0.4pt}\vspace{1mm}
Параметр: incident\_\allowbreak{}response\_\allowbreak{}training\newline Тип: boolean\newline Значення: \textbf{true} \\ \hline
101 & ПЕРЕВІРКА РЕАГУВАНЬ НА ІНЦИДЕНТИ (IR-3) &  & IR-3 & Можливості BPE-движка дозволяють запускати симуляції інцидентів у тестовому середовищі `itsm` для регулярного тестування планів реагування без впливу на продуктивну систему.\vspace{1mm}\newline\noindent\rule{\linewidth}{0.4pt}\vspace{1mm}
Параметр: incident\_\allowbreak{}testing\_\allowbreak{}simulation\newline Тип: boolean\newline Значення: \textbf{true} \\ \hline
102 & Обробка інциденту (IR-4) &  & IR-4 & Будь-який інцидент безпеки автоматично реєструється як тикет у модулі `itsm`. Процес обробки (`incident\_\allowbreak{}management\_\allowbreak{}process`) гарантує, що інцидент пройде всі необхідні етапи: виявлення, ізоляція, усунення, відновлення.\vspace{1mm}\newline\noindent\rule{\linewidth}{0.4pt}\vspace{1mm}
Параметр: incident\_\allowbreak{}handling\newline Тип: boolean\newline Значення: \textbf{true} \\ \hline
103 & Моніторинг інциденту (IR-5) &  & IR-5 & Статуси інцидентів у реальному часі відображаються на дашбордах `erpuno/itsm`. BPE-движок дозволяє призначати SLA-таймаути (Boundary Events) для ескалації інцидентів, якщо час реакції перевищено.\vspace{1mm}\newline\noindent\rule{\linewidth}{0.4pt}\vspace{1mm}
Параметр: incident\_\allowbreak{}monitoring\newline Тип: boolean\newline Значення: \textbf{true} \\ \hline
104 & ЗВІТНІСТЬ ПРО ІНЦИДЕНТИ (IR-6) &  & IR-6 & Система `erpuno/itsm` автоматично формує звіти про інциденти, надаючи вичерпну інформацію про порушені сервіси, час простою (SLA) та дії команди, що закриває вимоги внутрішньої та зовнішньої звітності.\vspace{1mm}\newline\noindent\rule{\linewidth}{0.4pt}\vspace{1mm}
Параметр: automated\_\allowbreak{}incident\_\allowbreak{}reporting\newline Тип: boolean\newline Значення: \textbf{true} \\ \hline
105 & ПІДТРИМКА РЕАГУВАННЯ НА ІНЦИДЕНТИ (IR-7) &  & IR-7 & Вбудовані засоби комунікації в рамках тикету (коментарі, призначення експертів, інтеграція з базою знань) забезпечують оперативну допомогу та координацію групи безпеки під час інциденту.\vspace{1mm}\newline\noindent\rule{\linewidth}{0.4pt}\vspace{1mm}
Параметр: incident\_\allowbreak{}assistance\newline Тип: boolean\newline Значення: \textbf{true} \\ \hline
106 & План реагування на інциденти (IR-8) &  & IR-8 & Сам конфігураційний код процесів `erpuno/itsm` виступає в ролі актуального та виконуваного плану реагування на інциденти (Executable Incident Response Plan), усуваючи розбіжність між документацією та реальністю.\vspace{1mm}\newline\noindent\rule{\linewidth}{0.4pt}\vspace{1mm}
Параметр: executable\_\allowbreak{}ir\_\allowbreak{}plan\newline Тип: boolean\newline Значення: \textbf{true} \\ \hline
\multicolumn{5}{|l|}{\textbf{9. ТЕХНІЧНЕ ОБСЛУГОВУВАННЯ (MA)}} \\ \hline
107 & ПОЛІТИКА ТА ПРОЦЕДУРИ ТЕХНІЧНОГО ОБСЛУГОВУВАННЯ (MA-1) & a. Планувати, документувати та переглядати записи з технічного обслуговування, ремонту або заміни компонентів системи відповідно до вимог виробника та постачальників та/або вимог організації.\newline b. Затвердити та здійснювати моніторинг усіх заходів з технічного обслуговування, незалежно від того, виконуються вони на місці або віддалено, а також чи обслуговуються системи або системні компоненти на місці, чи переміщуються в інше місце.\newline c. Вимагати, щоб [Призначення: визначені організацією персонал чи ролі] явно схвалили видалення системи або компоненту системи з організаційного обладнання для технічного обслуговування, ремонту чи заміни поза об'єктами експлуатації.\newline d. Очищати обладнання з погляду видалення всієї інформації з носіїв до вилучення обладнання організації для технічного обслуговування, ремонту чи заміни поза об'єктами експлуатації.\newline e. Перевірити всі потенційно порушені заходи захисту, щоб переконатися, що вони, як і раніше, працюють належним чином після дій з обслуговування, ремонту або заміни.\newline f. Вносити [Призначення: визначену організацією інформацію, пов'язану з технічним обслуговуванням] до записів з технічного обслуговування. & MA-1 & Технічне обслуговування компонентів ЄСІКС здійснюється авторизованим персоналом (ІСС) через безпечні канали з відповідним логуванням дій.\vspace{1mm}\newline\noindent\rule{\linewidth}{0.4pt}\vspace{1mm}
Параметр: Розроблено та задокументовано політику технічного обслуговування\newline Тип: list\newline Значення: \textbf{default\_\allowbreak{}deny\_\allowbreak{}rule, abac\_\allowbreak{}rule\_\allowbreak{}1}\vspace{1mm}\newline\noindent\rule{\linewidth}{0.4pt}\vspace{1mm}
Параметр: Політика технічного обслуговування поширюється на персонал або ролі\newline Тип: list\newline Значення: \textbf{admin, security\_\allowbreak{}officer}\vspace{1mm}\newline\noindent\rule{\linewidth}{0.4pt}\vspace{1mm}
Параметр: Розроблені та задокументовані процедури технічного обслуговування, що сприяють впровадженню політики технічного обслуговування та пов'язаних з нею заходів технічного обслуговування\newline Тип: list\newline Значення: \textbf{default\_\allowbreak{}deny\_\allowbreak{}rule, abac\_\allowbreak{}rule\_\allowbreak{}1}\vspace{1mm}\newline\noindent\rule{\linewidth}{0.4pt}\vspace{1mm}
Параметр: Посадова особа призначається для управління розробкою, документуванням та розповсюдженням політики та процедур технічного обслуговування; MA-01(c)[01][01] переглядається та оновлюється поточна політика обслуговування частота; MA-01(c)[01][02] переглядається та оновлюється поточна політика обслуговування після подій; MA-01(c)[02][01] переглядаються та оновлюються поточні процедури технічного обслуговування частота; MA-01(c)[02][02] переглядаються та оновлюються поточні процедури технічного обслуговування після подій\newline Тип: list\newline Значення: \textbf{admin, security\_\allowbreak{}officer} \\ 
 &  &  &  & Параметр: Визначено персонал або ролі, на які поширюється політика технічного обслуговування\newline Тип: list\newline Значення: \textbf{admin, security\_\allowbreak{}officer}\vspace{1mm}\newline\noindent\rule{\linewidth}{0.4pt}\vspace{1mm}
Параметр: Визначено персонал або ролі, на які поширюються процедури технічного обслуговування\newline Тип: list\newline Значення: \textbf{admin, security\_\allowbreak{}officer}\vspace{1mm}\newline\noindent\rule{\linewidth}{0.4pt}\vspace{1mm}
Параметр: Визначено посадову особу, яка керуватиме політикою та процедурами технічного обслуговування\newline Тип: list\newline Значення: \textbf{admin, security\_\allowbreak{}officer}\vspace{1mm}\newline\noindent\rule{\linewidth}{0.4pt}\vspace{1mm}
Параметр: Визначено частоту, з якою переглядається та оновлюється поточна політика технічного обслуговування\newline Тип: list\newline Значення: \textbf{default\_\allowbreak{}deny\_\allowbreak{}rule, abac\_\allowbreak{}rule\_\allowbreak{}1} \\ 
 &  &  &  & Параметр: Визначено події, які потребують перегляду та оновлення поточної політики технічного обслуговування\newline Тип: list\newline Значення: \textbf{default\_\allowbreak{}deny\_\allowbreak{}rule, abac\_\allowbreak{}rule\_\allowbreak{}1}\vspace{1mm}\newline\noindent\rule{\linewidth}{0.4pt}\vspace{1mm}
Параметр: Визначено частоту, з якою переглядаються та оновлюються поточні процедури технічного обслуговування\newline Тип: integer\newline Значення: \textbf{30}\vspace{1mm}\newline\noindent\rule{\linewidth}{0.4pt}\vspace{1mm}
Параметр: Визначено події, які потребують перегляду та оновлення процедур технічного обслуговування\newline Тип: list\newline Значення: \textbf{login, logout, failed\_\allowbreak{}attempt} \\ \hline
108 & ІНСТРУМЕНТИ ДЛЯ ОБСЛУГОВУВАННЯ (MA-3) & a. Затвердити, контролювати та відстежувати використання засобів технічного обслуговування.\newline b. Переглядати раніше затверджені інструменти технічного [Призначення: з частотою, визначеною організацією]. обслуговування & MA-3 & Технічне обслуговування компонентів ЄСІКС здійснюється авторизованим персоналом (ІСС) через безпечні канали з відповідним логуванням дій.\vspace{1mm}\newline\noindent\rule{\linewidth}{0.4pt}\vspace{1mm}
Параметр: Переглядаються раніше затверджені інструменти технічного обслуговування частота\newline Тип: string\newline Значення: \textbf{щорічно}\vspace{1mm}\newline\noindent\rule{\linewidth}{0.4pt}\vspace{1mm}
Параметр: Визначено частоту, з якою слід переглядати раніше затверджені інструменти технічного обслуговування\newline Тип: integer\newline Значення: \textbf{30} \\ \hline
109 & ІНСТРУМЕНТИ ДЛЯ ОБСЛУГОВУВАННЯ - ПЕРЕВІРКА ІНСТРУМЕНТІВ (MA-3(1)) &  & MA-3(1) & Технічне обслуговування компонентів ЄСІКС здійснюється авторизованим персоналом (ІСС) через безпечні канали з відповідним логуванням дій.\vspace{1mm}\newline\noindent\rule{\linewidth}{0.4pt}\vspace{1mm}
Параметр: Оглядаються інструменти для технічного обслуговування, які доставлені на об'єкт обслуговуючим персоналом, на предмет неправильних або несанкціонованих модифікацій\newline Тип: list\newline Значення: \textbf{admin, security\_\allowbreak{}officer} \\ \hline
110 & ІНСТРУМЕНТИ ДЛЯ ОБСЛУГОВУВАННЯ - ПЕРЕВІРКА НОСІЇВ ІНФОРМАЦІЇ (MA-3(2)) &  & MA-3(2) & Технічне обслуговування компонентів ЄСІКС здійснюється авторизованим персоналом (ІСС) через безпечні канали з відповідним логуванням дій. \\ \hline
111 & ІНСТРУМЕНТИ ДЛЯ ОБСЛУГОВУВАННЯ - ЗАПОБІГАННЯ НЕСАНКЦІОНОВАНОМУ ПЕРЕМІЩЕННЮ (MA-3(3)) &  & MA-3(3) & Технічне обслуговування компонентів ЄСІКС здійснюється авторизованим персоналом (ІСС) через безпечні канали з відповідним логуванням дій.\vspace{1mm}\newline\noindent\rule{\linewidth}{0.4pt}\vspace{1mm}
Параметр: Переміщення обладнання для технічного обслуговування, що містить інформацію організації, запобігається шляхом отримання дозволу від персонал або ролі\newline Тип: list\newline Значення: \textbf{admin, security\_\allowbreak{}officer}\vspace{1mm}\newline\noindent\rule{\linewidth}{0.4pt}\vspace{1mm}
Параметр: Визначено персонал або ролі, які можуть надавати дозвіл на переміщення обладнання з об'єкту\newline Тип: list\newline Значення: \textbf{admin, security\_\allowbreak{}officer} \\ \hline
112 & ВІДДАЛЕНЕ ОБСЛУГОВУВАННЯ (MA-4) &  & MA-4 & Технічне обслуговування компонентів ЄСІКС здійснюється авторизованим персоналом (ІСС) через безпечні канали з відповідним логуванням дій.\vspace{1mm}\newline\noindent\rule{\linewidth}{0.4pt}\vspace{1mm}
Параметр: Впроваджено віддалені дії з обслуговування та діагностики\newline Тип: list\newline Значення: \textbf{login, logout, failed\_\allowbreak{}attempt}\vspace{1mm}\newline\noindent\rule{\linewidth}{0.4pt}\vspace{1mm}
Параметр: Відстежуються віддалені дії з обслуговування та діагностики\newline Тип: list\newline Значення: \textbf{login, logout, failed\_\allowbreak{}attempt}\vspace{1mm}\newline\noindent\rule{\linewidth}{0.4pt}\vspace{1mm}
Параметр: Використання віддалених засобів технічного обслуговування та діагностики дозволено лише відповідно до політики організації\newline Тип: list\newline Значення: \textbf{default\_\allowbreak{}deny\_\allowbreak{}rule, abac\_\allowbreak{}rule\_\allowbreak{}1} \\ \hline
113 & ТЕХНІЧНИЙ ПЕРСОНАЛ (MA-5) & a. Встановити процедуру авторизації технічного персоналу та вести перелік авторизованих організацій технічного обслуговування або персоналу.\newline b. Перевіряти, що персонал, який не супроводжується та виконує технічне обслуговування в системі, має необхідні дозволи на доступ.\newline c. Визначити персонал організації з необхідними повноваженнями щодо доступу та технічною компетенцією для нагляду за персоналом з технічного обслуговування, який не має необхідних дозволів на доступ. & MA-5 & Технічне обслуговування компонентів ЄСІКС здійснюється авторизованим персоналом (ІСС) через безпечні канали з відповідним логуванням дій.\vspace{1mm}\newline\noindent\rule{\linewidth}{0.4pt}\vspace{1mm}
Параметр: Запроваджено процес авторизації технічного персоналу\newline Тип: list\newline Значення: \textbf{admin, security\_\allowbreak{}officer}\vspace{1mm}\newline\noindent\rule{\linewidth}{0.4pt}\vspace{1mm}
Параметр: Ведеться перелік авторизованих організацій або персоналу з технічного обслуговування\newline Тип: list\newline Значення: \textbf{admin, security\_\allowbreak{}officer}\vspace{1mm}\newline\noindent\rule{\linewidth}{0.4pt}\vspace{1mm}
Параметр: Персонал без супроводу, який виконує технічне обслуговування системи, має необхідні дозволи на доступ\newline Тип: list\newline Значення: \textbf{admin, security\_\allowbreak{}officer}\vspace{1mm}\newline\noindent\rule{\linewidth}{0.4pt}\vspace{1mm}
Параметр: Персонал організації з необхідними повноваженнями доступу та технічною компетентністю призначений/призначені для нагляду за діяльністю з технічного обслуговування персоналу, який не має необхідних дозволів на доступ\newline Тип: list\newline Значення: \textbf{admin, security\_\allowbreak{}officer} \\ \hline
\multicolumn{5}{|l|}{\textbf{10. ЗАХИСТ НОСІЇВ ІНФОРМАЦІЇ (MP)}} \\ \hline
114 & ПОЛІТИКА ТА ПРОЦЕДУРИ ЩОДО ЗАХИСТУ НОСІЇВ ІНФОРМАЦІЇ (MP-1) & a. Розробити, задокументувати та поширити серед [Призначення: визначеного організацією персоналу або посад]:\newline 1. 2. політику захисту носіїв інформації, яка:\newline (a) містить мету, сферу застосування, ролі, обов'язки, відповідальність керівництва, координацію між організаційними підрозділами та систему контролю відповідності (complaince);\newline (b) відповідає чинному законодавству, виконавчим наказам, директивам, нормам, політикам, стандартам та керівним принципам; процедури, які сприяють здійсненню політики та заходів захисту носіїв інформації.\newline b. Призначити [Призначення: визначену організацією посадову особу] для управління розробкою, документування, та розповсюдження політики та процедурами захисту носіїв інформації.\newline c. Переглядати та оновлювати чинну систему захисту носіїв інформації:\newline 1. поточну політику захисту носіїв інформації [Призначення: з визначеною організацією частотою];\newline 2. поточні процедури захисту носіїв інформації [Призначення: з визначеною організацією частотою]. & MP-1 & Криптографічні ключі та сертифікати зберігаються на захищених апаратних носіях (HSM) або в зашифрованих сховищах, доступ до яких обмежено.\vspace{1mm}\newline\noindent\rule{\linewidth}{0.4pt}\vspace{1mm}
Параметр: Розроблено та задокументовано політику захисту носіїв інформації\newline Тип: list\newline Значення: \textbf{default\_\allowbreak{}deny\_\allowbreak{}rule, abac\_\allowbreak{}rule\_\allowbreak{}1}\vspace{1mm}\newline\noindent\rule{\linewidth}{0.4pt}\vspace{1mm}
Параметр: Політика захисту носіїв інформації поширюється на персонал або ролі\newline Тип: list\newline Значення: \textbf{admin, security\_\allowbreak{}officer}\vspace{1mm}\newline\noindent\rule{\linewidth}{0.4pt}\vspace{1mm}
Параметр: Розроблено та задокументовано процедури захисту носіїв інформації, що сприятимуть реалізації політики захисту носіїв інформації та заходів захисту носіїв інформації\newline Тип: list\newline Значення: \textbf{default\_\allowbreak{}deny\_\allowbreak{}rule, abac\_\allowbreak{}rule\_\allowbreak{}1}\vspace{1mm}\newline\noindent\rule{\linewidth}{0.4pt}\vspace{1mm}
Параметр: Посадова особа призначається для управління розробкою, документуванням та розповсюдженням політики та процедур захисту носіїв інформації. MP-01(c)[01][01] переглядається та оновлюється поточна політика захисту носіїв інформації частота; MP-01(c)[01][02] переглядається та оновлюється поточна політика захисту носіїв інформації після подій; MP-01(c)[02][01] переглядаються та оновлюються поточні процедури захисту носіїв інформації частота; MP-01(c)[02][02] переглядаються та оновлюються поточні процедури захисту носіїв інформації після подій\newline Тип: list\newline Значення: \textbf{admin, security\_\allowbreak{}officer} \\ 
 &  &  &  & Параметр: Визначено персонал або ролі, серед яких має бути поширена політика захисту носіїв інформації\newline Тип: list\newline Значення: \textbf{admin, security\_\allowbreak{}officer}\vspace{1mm}\newline\noindent\rule{\linewidth}{0.4pt}\vspace{1mm}
Параметр: Визначено персонал або ролі, серед яких мають бути поширені процедури захисту носіїв інформації\newline Тип: list\newline Значення: \textbf{admin, security\_\allowbreak{}officer}\vspace{1mm}\newline\noindent\rule{\linewidth}{0.4pt}\vspace{1mm}
Параметр: Визначено посадову особу, яка керуватиме політикою та процедурами захисту носіїв інформації\newline Тип: list\newline Значення: \textbf{admin, security\_\allowbreak{}officer}\vspace{1mm}\newline\noindent\rule{\linewidth}{0.4pt}\vspace{1mm}
Параметр: Визначено частоту, з якою переглядається та оновлюється поточна політика захисту носіїв інформації\newline Тип: list\newline Значення: \textbf{default\_\allowbreak{}deny\_\allowbreak{}rule, abac\_\allowbreak{}rule\_\allowbreak{}1} \\ 
 &  &  &  & Параметр: Визначено події, які потребують перегляду та оновлення чинної політики захисту носіїв інформації\newline Тип: list\newline Значення: \textbf{default\_\allowbreak{}deny\_\allowbreak{}rule, abac\_\allowbreak{}rule\_\allowbreak{}1}\vspace{1mm}\newline\noindent\rule{\linewidth}{0.4pt}\vspace{1mm}
Параметр: Визначено частоту, з якою переглядаються та оновлюються чинні процедури захисту носіїв інформації\newline Тип: integer\newline Значення: \textbf{30}\vspace{1mm}\newline\noindent\rule{\linewidth}{0.4pt}\vspace{1mm}
Параметр: Визначено події, які потребують перегляду та оновлення процедур захисту носіїв інформації\newline Тип: list\newline Значення: \textbf{login, logout, failed\_\allowbreak{}attempt} \\ \hline
115 & ДОСТУП ДО НОСІЇВ ІНФОРМАЦІЇ (MP-2) & Обмежити доступ до [Призначення: визначених організацією типів цифрових та/або нецифрових носіїв інформації] [Призначення: визначеним організацією персоналом або ролями]. & MP-2 & Криптографічні ключі та сертифікати зберігаються на захищених апаратних носіях (HSM) або в зашифрованих сховищах, доступ до яких обмежено.\vspace{1mm}\newline\noindent\rule{\linewidth}{0.4pt}\vspace{1mm}
Параметр: Доступ до типів цифрових носіїв інформації обмежено для персоналу або ролей\newline Тип: list\newline Значення: \textbf{admin, security\_\allowbreak{}officer}\vspace{1mm}\newline\noindent\rule{\linewidth}{0.4pt}\vspace{1mm}
Параметр: Доступ до типів нецифрових носіїв інформації обмежено для персоналу або ролей\newline Тип: list\newline Значення: \textbf{admin, security\_\allowbreak{}officer}\vspace{1mm}\newline\noindent\rule{\linewidth}{0.4pt}\vspace{1mm}
Параметр: Визначено персонал або ролі, уповноважені на доступ до цифрових носіїв інформації\newline Тип: list\newline Значення: \textbf{admin, security\_\allowbreak{}officer}\vspace{1mm}\newline\noindent\rule{\linewidth}{0.4pt}\vspace{1mm}
Параметр: Визначено персонал або ролі, уповноважені на доступ до нецифрових носіїв інформації\newline Тип: list\newline Значення: \textbf{admin, security\_\allowbreak{}officer} \\ \hline
116 & МАРКУВАННЯ НОСІЇВ ІНФОРМАЦІЇ (MP-3) & a. Наносити на носії інформації маркування, що вказують на обмеження поширення, обробки, а також застереження та відповідні мітки безпеки (якщо такі є) інформації.\newline b. Звільнити [Призначення: визначені організацією типи носіїв системи] від маркування, якщо носії залишаються в межах [Призначення: визначених організацією контрольованих зон]. & MP-3 & Криптографічні ключі та сертифікати зберігаються на захищених апаратних носіях (HSM) або в зашифрованих сховищах, доступ до яких обмежено.\vspace{1mm}\newline\noindent\rule{\linewidth}{0.4pt}\vspace{1mm}
Параметр: Визначено типи носіїв інформації, які звільняються від маркування під час перебування на контрольованих зонах\newline Тип: integer\newline Значення: \textbf{30} \\ \hline
117 & ЗБЕРІГАННЯ НОСІЇВ ІНФОРМАЦІЇ (MP-4) & a. Фізично контролювати та безпечно зберігати [Призначення: визначені організацією типи цифрових та/або нецифрових носіїв інформації] в межах [Призначення: визначених організацією контрольованих зон].\newline b. Захищати системні носії, які визначені в МР-4 до того часу, як носії знищуються або очищаються, з використанням затвердженого обладнання, методів та процедур. & MP-4 & Криптографічні ключі та сертифікати зберігаються на захищених апаратних носіях (HSM) або в зашифрованих сховищах, доступ до яких обмежено. \\ \hline
118 & Транспортування носіїв інформації (MP-5) & a. Захищати та контролювати [Призначення: визначені організацією типи носіїв системи] під час транспортування за межами контрольованих зон, використовуючи [Призначення: визначені організацією заходи безпеки].\newline b. Вести облік носіїв системи інформації під час транспортування за межами контрольованих зон.\newline c. Документувати дії, пов'язані з транспортуванням носіїв системи.\newline d. Обмежити діяльність уповноваженого персоналу, пов'язану з транспортуванням носіїв системи. & MP-5 & Криптографічні ключі та сертифікати зберігаються на захищених апаратних носіях (HSM) або в зашифрованих сховищах, доступ до яких обмежено.\vspace{1mm}\newline\noindent\rule{\linewidth}{0.4pt}\vspace{1mm}
Параметр: Типи носіїв інформації системи захищаються під час транспортування за межі контрольованих зон за допомогою заходів безпеки\newline Тип: integer\newline Значення: \textbf{30}\vspace{1mm}\newline\noindent\rule{\linewidth}{0.4pt}\vspace{1mm}
Параметр: Типи носіїв інформації системи контролюються під час транспортування за межі контрольованих зон за допомогою заходів безпеки\newline Тип: integer\newline Значення: \textbf{30}\vspace{1mm}\newline\noindent\rule{\linewidth}{0.4pt}\vspace{1mm}
Параметр: Під час транспортування за межі контрольованих зон ведеться облік носіїв інформації системи\newline Тип: integer\newline Значення: \textbf{30}\vspace{1mm}\newline\noindent\rule{\linewidth}{0.4pt}\vspace{1mm}
Параметр: Визначено персонал, уповноважений здійснювати діяльність з транспортування носіїв інформації\newline Тип: list\newline Значення: \textbf{admin, security\_\allowbreak{}officer} \\ 
 &  &  &  & Параметр: Діяльність, пов'язана з транспортуванням носіїв інформації системи, обмежується визначеним уповноваженим персоналом\newline Тип: list\newline Значення: \textbf{admin, security\_\allowbreak{}officer}\vspace{1mm}\newline\noindent\rule{\linewidth}{0.4pt}\vspace{1mm}
Параметр: Визначено типи носіїв інформації системи для захисту та контролю під час транспортування за межі контрольованих зон\newline Тип: integer\newline Значення: \textbf{30} \\ \hline
119 & ЗНИЩЕННЯ ІНФОРМАЦІЇ НА НОСІЯХ ІНФОРМАЦІЇ (MP-6) & a. Очищувати [Призначення: визначені організацією системні носії] перед утилізацією, випуском за межі організаційного контролю, або перед повторним використанням [Призначення: методами та процедурами очищення, визначеними організацією].\newline b. Використовувати механізми очищення зі стійкістю та цілісністю, що відповідає категорії безпеки або рівню секретності інформації. & MP-6 & Криптографічні ключі та сертифікати зберігаються на захищених апаратних носіях (HSM) або в зашифрованих сховищах, доступ до яких обмежено.\vspace{1mm}\newline\noindent\rule{\linewidth}{0.4pt}\vspace{1mm}
Параметр: Застосовуються механізми очищення, надійність і цілісність яких відповідає категорії безпеки або рівню секретності інформації\newline Тип: string\newline Значення: \textbf{автоматизований засіб моніторингу} \\ \hline
120 & ВИКОРИСТАННЯ НОСІЇВ ІНФОРМАЦІЇ (MP-7) & a. [Вибір: обмежити; заборонити] використання [Призначення: визначених організацією типів носіїв системи] на [Призначення: визначені організацією системи або компоненти системи], використовуючи [Призначення: визначені організацією заходи безпеки].\newline b. Заборонити використання портативних пристроїв зберігання даних в системах організації, якщо такі пристрої не мають визначеного власника. & MP-7 & Криптографічні ключі та сертифікати зберігаються на захищених апаратних носіях (HSM) або в зашифрованих сховищах, доступ до яких обмежено. \\ \hline
121 & ЗНИЖЕННЯ КАТЕГОРІЇ БЕЗПЕКИ НОСІЇВ ІНФОРМАЦІЇ - ТАЄМНА ІНФОРМАЦІЯ (MP-8(4)) &  & MP-8(4) & Криптографічні ключі та сертифікати зберігаються на захищених апаратних носіях (HSM) або в зашифрованих сховищах, доступ до яких обмежено.\vspace{1mm}\newline\noindent\rule{\linewidth}{0.4pt}\vspace{1mm}
Параметр: Носії інформації, що містять інформацію з обмеженим доступом, знижуються в класі перед передачею особам, які не мають необхідних дозволів на доступ\newline Тип: list\newline Значення: \textbf{admin, security\_\allowbreak{}officer} \\ \hline
\multicolumn{5}{|l|}{\textbf{11. ФІЗИЧНИЙ ЗАХИСТ ТА ЗАХИСТ НАВКОЛИШНЬОГО СЕРЕДОВИЩА (PE)}} \\ \hline
122 & РЕ-1 фізичного захисту та захисту робочого середовища Авторизація фізичного (PE-1) &  & PE-1 & Фізичний захист серверного обладнання та носіїв інформації забезпечується центрами обробки даних (ЦОД) ЄСІКС з багаторівневим контролем доступу.\vspace{1mm}\newline\noindent\rule{\linewidth}{0.4pt}\vspace{1mm}
Параметр: Розроблено та задокументовано політику фізичного захисту та захисту робочого середовища\newline Тип: list\newline Значення: \textbf{default\_\allowbreak{}deny\_\allowbreak{}rule, abac\_\allowbreak{}rule\_\allowbreak{}1}\vspace{1mm}\newline\noindent\rule{\linewidth}{0.4pt}\vspace{1mm}
Параметр: Політика фізичного захисту та захисту робочого середовища розповсюджується серед персоналу або ролей\newline Тип: list\newline Значення: \textbf{admin, security\_\allowbreak{}officer}\vspace{1mm}\newline\noindent\rule{\linewidth}{0.4pt}\vspace{1mm}
Параметр: Посадова особа призначається для управління розробкою, документуванням та розповсюдженням політики та процедур фізичного захисту та захисту робочого середовища; PE-01(c)[01][01] переглядається та оновлюється поточна політика фізичного захисту та захисту робочого середовища частота; PE-01(c)[01][02] поточна політика фізичного захисту та захисту робочого середовища переглядається та оновлюється після подій; PE-01(c)[02][01] переглядаються та оновлюються поточні процедури фізичного захисту та захисту робочого середовища частота; поточні процедури фізичного та екологічного захисту переглядаються та оновлюються після подій. PE-01(c)[02][02]\newline Тип: list\newline Значення: \textbf{admin, security\_\allowbreak{}officer}\vspace{1mm}\newline\noindent\rule{\linewidth}{0.4pt}\vspace{1mm}
Параметр: Визначено персонал або ролі, до яких має бути доведена політика фізичного захисту та захисту робочого середовища\newline Тип: list\newline Значення: \textbf{admin, security\_\allowbreak{}officer} \\ 
 &  &  &  & Параметр: Визначено персонал або ролі, на які поширюються процедури фізичного захисту та захисту робочого середовища\newline Тип: list\newline Значення: \textbf{admin, security\_\allowbreak{}officer}\vspace{1mm}\newline\noindent\rule{\linewidth}{0.4pt}\vspace{1mm}
Параметр: Визначено посадову особу, яка керуватиме політикою та процедурами фізичного захисту та захисту робочого середовища\newline Тип: list\newline Значення: \textbf{admin, security\_\allowbreak{}officer}\vspace{1mm}\newline\noindent\rule{\linewidth}{0.4pt}\vspace{1mm}
Параметр: Визначено частоту, з якою переглядається та оновлюється поточна політика фізичного захисту та захисту робочого середовища\newline Тип: list\newline Значення: \textbf{default\_\allowbreak{}deny\_\allowbreak{}rule, abac\_\allowbreak{}rule\_\allowbreak{}1}\vspace{1mm}\newline\noindent\rule{\linewidth}{0.4pt}\vspace{1mm}
Параметр: Визначено події, які потребують перегляду та оновлення поточної політики фізичного захисту та захисту робочого середовища\newline Тип: list\newline Значення: \textbf{default\_\allowbreak{}deny\_\allowbreak{}rule, abac\_\allowbreak{}rule\_\allowbreak{}1} \\ 
 &  &  &  & Параметр: Визначено частоту, з якою переглядаються та оновлюються поточні процедури фізичного захисту та захисту робочого середовища\newline Тип: integer\newline Значення: \textbf{30}\vspace{1mm}\newline\noindent\rule{\linewidth}{0.4pt}\vspace{1mm}
Параметр: Визначено події, які потребують перегляду та оновлення процедур фізичного захисту та захисту робочого середовища\newline Тип: list\newline Значення: \textbf{login, logout, failed\_\allowbreak{}attempt} \\ \hline
123 & РЕ-2 доступу (PE-2) &  & PE-2 & Регламентується інструкціями з режиму ЦОД (див. розділ PE у digital-profile.pdf).\vspace{1mm}\newline\noindent\rule{\linewidth}{0.4pt}\vspace{1mm}
Параметр: physical\_\allowbreak{}access\_\allowbreak{}doc\_\allowbreak{}ref\newline Тип: string\newline Значення: \textbf{digital-profile.pdf} \\ \hline
124 & ФІЗИЧНИЙ ДОСТУП ДО СИСТЕМИ &  & PE-3 & Регламентується системами СКУД будівлі (див. розділ PE у digital-profile.pdf).\vspace{1mm}\newline\noindent\rule{\linewidth}{0.4pt}\vspace{1mm}
Параметр: physical\_\allowbreak{}access\_\allowbreak{}control\_\allowbreak{}ref\newline Тип: string\newline Значення: \textbf{digital-profile.pdf} \\ \hline
125 & РЕ-4 ліній електроживлення (PE-4) &  & PE-4 & Фізичний захист серверного обладнання та носіїв інформації забезпечується центрами обробки даних (ЦОД) ЄСІКС з багаторівневим контролем доступу. \\ \hline
126 & КОНТРОЛЬ ДОСТУПУ В ПРИМІЩЕННЯ ДЛЯ ВІДОБРАЖЕННЯ ІНФОРМАЦІЇ &  & PE-5 & Фізичний захист серверного обладнання та носіїв інформації забезпечується центрами обробки даних (ЦОД) ЄСІКС з багаторівневим контролем доступу. \\ \hline
127 & МОНІТОРИНГ ФІЗИЧНОГО ДОСТУПУ &  & PE-6 & Фізичний захист серверного обладнання та носіїв інформації забезпечується центрами обробки даних (ЦОД) ЄСІКС з багаторівневим контролем доступу.\vspace{1mm}\newline\noindent\rule{\linewidth}{0.4pt}\vspace{1mm}
Параметр: Переглядаються журнали фізичного доступу частота\newline Тип: string\newline Значення: \textbf{щорічно}\vspace{1mm}\newline\noindent\rule{\linewidth}{0.4pt}\vspace{1mm}
Параметр: Визначено частоту перегляду журналів фізичного доступу\newline Тип: integer\newline Значення: \textbf{30}\vspace{1mm}\newline\noindent\rule{\linewidth}{0.4pt}\vspace{1mm}
Параметр: Визначено події або потенційні ознаки подій, що вимагають перегляду журналів фізичного доступу\newline Тип: list\newline Значення: \textbf{login, logout, failed\_\allowbreak{}attempt} \\ \hline
128 & АЛЬТЕРНАТИВНЕ РОБОЧЕ МІСЦЕ (PE-17) &  & PE-17 & Фізичний захист серверного обладнання та носіїв інформації забезпечується центрами обробки даних (ЦОД) ЄСІКС з багаторівневим контролем доступу.\vspace{1mm}\newline\noindent\rule{\linewidth}{0.4pt}\vspace{1mm}
Параметр: Визначаються заходи захисту, які будуть застосовуватися на альтернативних робочих місцях; РЕ-17(a) альтернативні робочі місця визначені та задокументовані; РЕ-17(b) заходи захисту впроваджені на альтернативних робочих місцях; РЕ-17(c) оцінюється ефективність заходів захисту на альтернативних робочих місцях; РЕ-17(d) працівникам надаються засоби комунікації з персоналом служби інформаційної безпеки на випадок інцидентів\newline Тип: list\newline Значення: \textbf{admin, security\_\allowbreak{}officer} \\ \hline
129 & МАРКУВАННЯ КОМПОНЕНТІВ (PE-22) &  & PE-22 & Фізичний захист серверного обладнання та носіїв інформації забезпечується центрами обробки даних (ЦОД) ЄСІКС з багаторівневим контролем доступу. \\ \hline
\multicolumn{5}{|l|}{\textbf{12. ПЛАНУВАННЯ (PL)}} \\ \hline
130 & ПОЛІТИКИ ТА ПРОЦЕДУРИ ПЛАНУВАННЯ БЕЗПЕКИ (PL-1) & a. Розробити, задокументувати та поширити серед [Призначення: визначеного організацією персоналу або ролей]:\newline 1. 2. [вибір (один або декілька): рівень організації; рівень місії/бізнес процесу; рівень системи] політику планування, яка:\newline (a) містить мету, сферу застосування, ролі, обов'язки, відповідальність керівництва, координацію між організаційними підрозділами та системою контролю (complaince);\newline (b) відповідає чинному законодавству, виконавчим наказам, директивам, нормам, політикам, стандартам та керівним принципам; процедури, що полегшують здійснення планування політики безпеки та приватності й пов'язані з ними заходи.\newline b. Призначити [Призначення: визначену організацією посадову особу] для управління політикою та процедурами планування політики безпеки та приватності.\newline c. Переглядати та оновлювати поточне планування:\newline 1. політики планування безпеки та приватності [Призначення:з визначеною організацією частотою] та наступні [Призначення: визначені організацією події];\newline 2. поточні процедури планування політики [Призначення: з визначеною організацією [Призначення: визначені організацією події]. безпеки та приватності частотою] та наступні & PL-1 & Планування архітектури та безпеки базується на методології ISO/IEC/IEEE 42010, фреймворку Закмана та ДСТУ.\vspace{1mm}\newline\noindent\rule{\linewidth}{0.4pt}\vspace{1mm}
Параметр: Визначено персонал або ролі, до яких має бути доведена політика планування безпеки\newline Тип: list\newline Значення: \textbf{admin, security\_\allowbreak{}officer}\vspace{1mm}\newline\noindent\rule{\linewidth}{0.4pt}\vspace{1mm}
Параметр: Визначено персонал або ролі, на які поширюватимуться процедури планування безпеки\newline Тип: list\newline Значення: \textbf{admin, security\_\allowbreak{}officer}\vspace{1mm}\newline\noindent\rule{\linewidth}{0.4pt}\vspace{1mm}
Параметр: Визначено посадову особу, яка керуватиме політикою та процедурами планування безпеки\newline Тип: list\newline Значення: \textbf{admin, security\_\allowbreak{}officer}\vspace{1mm}\newline\noindent\rule{\linewidth}{0.4pt}\vspace{1mm}
Параметр: Визначено періодичність перегляду та оновлення поточної політики планування безпеки\newline Тип: list\newline Значення: \textbf{default\_\allowbreak{}deny\_\allowbreak{}rule, abac\_\allowbreak{}rule\_\allowbreak{}1} \\ 
 &  &  &  & Параметр: Є події, які потребують перегляду та оновлення поточної політики планування безпеки\newline Тип: list\newline Значення: \textbf{default\_\allowbreak{}deny\_\allowbreak{}rule, abac\_\allowbreak{}rule\_\allowbreak{}1}\vspace{1mm}\newline\noindent\rule{\linewidth}{0.4pt}\vspace{1mm}
Параметр: Визначена частота, з якою переглядаються та оновлюються поточні процедури планування безпеки\newline Тип: string\newline Значення: \textbf{щорічно} \\ \hline
131 & ПЛАНИ ЗАХИСТУ ІНФОРМАЦІЇ ТА ПЕРСОНАЛЬНИХ ДАНИХ (PL-2) & a. Розробити план захисту інформації та персональних даних для інформаційної системи, який:\newline 1. узгоджується з архітектурою підприємства організації;\newline 2. чітко визначає складові компоненти системи;\newline 3. описує оперативний контекст інформаційної системи з точки зору завдань та процесів;\newline 4. визначає осіб, які виконують системні ролі та обов'язки;\newline 5. визначає тип інформації, яка обробляється, зберігається та передається системою;\newline 6. надає огляд вимог безпеки та приватності інформаційної системи;\newline 7. описує будь які конкретні загрози системи, які викликають стурбованість організації;\newline 8. надає результати оцінки ризику конфіденційності для систем, в яких обробляються персональні дані;\newline 9. описує робоче середовище інформаційної системи та будь які залежності від систем або компонентів систем або підключень до таких систем та їх компонентів;\newline 10. надає огляд вимог безпеки та конфіденційності системи;\newline 11. визначає будь які відповідні контрольні базові рівні або накладання, якщо вони застосовуються;\newline 12. описує чинні або заплановані заходи щодо забезпечення безпеки та приватності, включно з обґрунтуванням рішень щодо налаштування\newline 13. включає виявлення ризиків для архітектури безпеки і приватності, а також проєктних рішень;\newline 14. включає дії, пов'язані з безпекою та конфіденційністю, які впливають на систему, виконання яких вимагає планування та координацію з [Призначення: визначені організацією окремі особи або групи];\newline 15. розглядається та затверджується уповноваженою посадовою особою або призначеним представником до початку реалізації плану.\newline b. Поширити копії планів захисту інформації та персональних даних і повідомляти про подальші зміни планів серед [Призначення:визначеного організацією персоналу або ролей].\newline c. Переглядати плани захисту інформації та персональних даних [Призначення: з визначеною організацією частотою].\newline d. Оновлювати плани захисту інформації та персональних даних для врахування змін в інформаційній системі й робочому середовищі або проблем, виявлених у ході реалізації або оцінювання заходів безпеки та приватності.\newline e. забезпечити захист планів захисту інформації та персональних даних від несанкціонованого розголошення та змін. & PL-2 & Планування архітектури та безпеки базується на методології ISO/IEC/IEEE 42010, фреймворку Закмана та ДСТУ.\vspace{1mm}\newline\noindent\rule{\linewidth}{0.4pt}\vspace{1mm}
Параметр: system\_\allowbreak{}security\_\allowbreak{}plan\newline Тип: boolean\newline Значення: \textbf{true} \\ \hline
132 & ПРАВИЛА ПОВЕДІНКИ (PL-4) & a. Створити та надати особам, що потребують доступу до інформаційної системи, правила, які описують їхні обов'язки й очікувану поведінку щодо інформації та використання інформаційної системи, безпеки та приватності.\newline b. Отримати документальне підтвердження від таких осіб про те, що вони прочитали, зрозуміли та погодилися дотримуватися правил поведінки, перш ніж дозволяти доступ до інформації та інформаційної системи.\newline c. Переглядати й оновлювати правила поведінки [Призначення: з визначеною організацією частотою].\newline d. Вимагати від осіб, які підписали попередню версію правил поведінки, перечитати та повторно підписати правила [Вибір (один або декілька): [Призначення: з визначеною організацією частотою]; коли правила переглядаються чи оновлюються]. & PL-4 & Планування архітектури та безпеки базується на методології ISO/IEC/IEEE 42010, фреймворку Закмана та ДСТУ.\vspace{1mm}\newline\noindent\rule{\linewidth}{0.4pt}\vspace{1mm}
Параметр: rules\_\allowbreak{}of\_\allowbreak{}behavior\newline Тип: boolean\newline Значення: \textbf{true} \\ \hline
\multicolumn{5}{|l|}{\textbf{13. БЕЗПЕКА ПЕРСОНАЛУ (PS)}} \\ \hline
133 & Політика та процедури кадрової безпеки (PS-1) &  & PS-1 & Вимоги до безпеки персоналу визначаються політиками управління кадрами та процедурами допуску до інформації в системі ЄСІКС.\vspace{1mm}\newline\noindent\rule{\linewidth}{0.4pt}\vspace{1mm}
Параметр: Визначено персонал або ролі, на які поширюється політика кадрової безпеки\newline Тип: list\newline Значення: \textbf{admin, security\_\allowbreak{}officer}\vspace{1mm}\newline\noindent\rule{\linewidth}{0.4pt}\vspace{1mm}
Параметр: Визначено персонал або ролі, на які поширюються процедури кадрової безпеки\newline Тип: list\newline Значення: \textbf{admin, security\_\allowbreak{}officer}\vspace{1mm}\newline\noindent\rule{\linewidth}{0.4pt}\vspace{1mm}
Параметр: Визначено посадову особу, яка керуватиме політикою та процедурами кадрової безпеки\newline Тип: list\newline Значення: \textbf{admin, security\_\allowbreak{}officer}\vspace{1mm}\newline\noindent\rule{\linewidth}{0.4pt}\vspace{1mm}
Параметр: Визначена періодичність перегляду та оновлення поточної політики кадрової безпеки\newline Тип: list\newline Значення: \textbf{default\_\allowbreak{}deny\_\allowbreak{}rule, abac\_\allowbreak{}rule\_\allowbreak{}1} \\ 
 &  &  &  & Параметр: Є події, які вимагають перегляду та оновлення поточної політики кадрової безпеки\newline Тип: list\newline Значення: \textbf{default\_\allowbreak{}deny\_\allowbreak{}rule, abac\_\allowbreak{}rule\_\allowbreak{}1}\vspace{1mm}\newline\noindent\rule{\linewidth}{0.4pt}\vspace{1mm}
Параметр: Визначено періодичність перегляду та оновлення поточних процедур кадрової безпеки\newline Тип: string\newline Значення: \textbf{щорічно} \\ \hline
134 & Перевірка персоналу (PS-3) &  & PS-3 & Вимоги до безпеки персоналу визначаються політиками управління кадрами та процедурами допуску до інформації в системі ЄСІКС.\vspace{1mm}\newline\noindent\rule{\linewidth}{0.4pt}\vspace{1mm}
Параметр: Визначені умови, що вимагають повторної перевірки осіб\newline Тип: list\newline Значення: \textbf{}\vspace{1mm}\newline\noindent\rule{\linewidth}{0.4pt}\vspace{1mm}
Параметр: Визначена частота повторної перевірки осіб, для яких це показано; PS-03a. проходять особи перевірку перед тим, як надати їм доступ до системи; PS-03b.[01] проходять особи повторну перевірку відповідно до умови, що вимагають повторної перевірки; PS-03b.[02] проводиться повторна перевірка у випадках, коли це зазначено, частота\newline Тип: list\newline Значення: \textbf{admin, security\_\allowbreak{}officer} \\ \hline
135 & Звільнення персоналу (PS-4) &  & PS-4 & Вимоги до безпеки персоналу визначаються політиками управління кадрами та процедурами допуску до інформації в системі ЄСІКС.\vspace{1mm}\newline\noindent\rule{\linewidth}{0.4pt}\vspace{1mm}
Параметр: Визначено період часу, протягом якого забороняється доступ до системи\newline Тип: integer\newline Значення: \textbf{30}\vspace{1mm}\newline\noindent\rule{\linewidth}{0.4pt}\vspace{1mm}
Параметр: Визначені теми інформаційної безпеки для обговорення під час проведення співбесід; PS-04a. при звільненні працівника доступ до системи вимикається протягом часового періоду; PS-04b. припиняють дію або анулюють будь-які автентифікатори та облікові дані після припинення трудових відносин з окремими особами; PS-04c. проводяться при звільненні окремих працівників співбесіди, які включають обговорення питань інформаційної безпеки; PS-04d. отримується після звільнення особи все майно, пов'язане з безпекою організаційної системи; PS-04e. зберігається доступ до організаційної інформації та систем, які раніше перебували під контролем особи, що звільняється, після припинення нею трудових відносин\newline Тип: list\newline Значення: \textbf{admin, security\_\allowbreak{}officer} \\ \hline
136 & Переведення персоналу (PS-5) &  & PS-5 & Вимоги до безпеки персоналу визначаються політиками управління кадрами та процедурами допуску до інформації в системі ЄСІКС.\vspace{1mm}\newline\noindent\rule{\linewidth}{0.4pt}\vspace{1mm}
Параметр: Визначені дії, які мають бути ініційовані після переведення або перепризначення\newline Тип: list\newline Значення: \textbf{login, logout, failed\_\allowbreak{}attempt}\vspace{1mm}\newline\noindent\rule{\linewidth}{0.4pt}\vspace{1mm}
Параметр: Визначено період часу, протягом якого мають бути здійснені дії з переведення або перепризначення після переведення або перепризначення\newline Тип: list\newline Значення: \textbf{login, logout, failed\_\allowbreak{}attempt}\vspace{1mm}\newline\noindent\rule{\linewidth}{0.4pt}\vspace{1mm}
Параметр: Визначено персонал або ролі, про які необхідно повідомляти, коли осіб призначають на інші посади або переводять на інші посади в організації\newline Тип: list\newline Значення: \textbf{admin, security\_\allowbreak{}officer}\vspace{1mm}\newline\noindent\rule{\linewidth}{0.4pt}\vspace{1mm}
Параметр: Визначено період часу, протягом якого необхідно повідомляти визначений організацією персонал або ролі, коли осіб перепризначають або переводять на інші посади в межах організації; PS-05a. переглядаються та підтверджуються поточні потреби в логічних та фізичних дозволах на доступ до систем та об'єктів при перепризначенні або переведенні осіб на інші посади в організації; PS-05b. були ініційовані дії з переведення або перепризначення протягом періоду часу після формальної дії з переведення; PS-05c. змінюється авторизація доступу за необхідності, щоб відповідати будь-яким змінам в оперативних потребах у зв'язку з перепризначенням або переведенням; PS-05d. було повідомлено персонал або ролі протягом часового періоду\newline Тип: list\newline Значення: \textbf{admin, security\_\allowbreak{}officer} \\ \hline
\multicolumn{5}{|l|}{\textbf{14. ОЦІНКА РИЗИКІВ (RA)}} \\ \hline
137 & ПОЛІТИКА ТА ПРОЦЕДУРИ ОЦІНЮВАННЯ РИЗИКУ (RA-1) & a. Розробити, задокументувати та поширити серед [Призначення: визначеного організацією персоналу або посадових осіб]:\newline 1. 2. Політику оцінювання ризику, яка:\newline (a) містить мету, сферу застосування, ролі, обов'язки, відповідальність керівництва, координацію між організаційними підрозділами та систему контролю відповідності (complaince);\newline (b) відповідає чинному законодавству, виконавчим наказам, директивам, нормам, політикам, стандартам і керівним принципам. Процедури, що сприяють здійсненню політики оцінювання ризику та пов'язаних з ними заходів оцінювання ризику.\newline b. Призначити [Призначення: визначена організацією посадову особу] для управління політикою та процедурами оцінювання ризику.\newline c. Переглядати й оновлювати:\newline 1. Поточну політику оцінювання організацією частотою]. ризику [Призначення: з визначеною\newline 2. Поточні процедури оцінювання ризику [Призначення: з визначеною організацією частотою]. & RA-1 & Оцінка ризиків здійснюється з урахуванням вимог OWASP Top 10 та чинного законодавства України з кібербезпеки.\vspace{1mm}\newline\noindent\rule{\linewidth}{0.4pt}\vspace{1mm}
Параметр: Визначено персонал або ролі, на які поширюється політика оцінювання ризику\newline Тип: list\newline Значення: \textbf{admin, security\_\allowbreak{}officer}\vspace{1mm}\newline\noindent\rule{\linewidth}{0.4pt}\vspace{1mm}
Параметр: Визначено персонал або ролі, на які поширюються процедури, що сприяють здійсненню політики оцінювання ризику\newline Тип: list\newline Значення: \textbf{admin, security\_\allowbreak{}officer}\vspace{1mm}\newline\noindent\rule{\linewidth}{0.4pt}\vspace{1mm}
Параметр: Визначена посадова особа, відповідальна за управління політикою та процедурами оцінювання ризику\newline Тип: list\newline Значення: \textbf{admin, security\_\allowbreak{}officer}\vspace{1mm}\newline\noindent\rule{\linewidth}{0.4pt}\vspace{1mm}
Параметр: Визначена періодичність перегляду та оновлення поточної політики оцінювання ризику\newline Тип: list\newline Значення: \textbf{default\_\allowbreak{}deny\_\allowbreak{}rule, abac\_\allowbreak{}rule\_\allowbreak{}1} \\ 
 &  &  &  & Параметр: Є події, які вимагають перегляду та оновлення поточної політики оцінювання ризику\newline Тип: list\newline Значення: \textbf{default\_\allowbreak{}deny\_\allowbreak{}rule, abac\_\allowbreak{}rule\_\allowbreak{}1}\vspace{1mm}\newline\noindent\rule{\linewidth}{0.4pt}\vspace{1mm}
Параметр: Визначено періодичність перегляду поточних процедур оцінювання ризику\newline Тип: string\newline Значення: \textbf{щорічно} \\ \hline
138 & ОЦІНЮВАННЯ РИЗИКУ (RA-3) &  & RA-3 & Всі ідентифіковані вразливості або загрози фіксуються в реєстрі `itsm`, де їм присвоюється рівень критичності. На базі цього рівня розраховується пріоритет та призначаються відповідальні особи.\vspace{1mm}\newline\noindent\rule{\linewidth}{0.4pt}\vspace{1mm}
Параметр: risk\_\allowbreak{}assessment\newline Тип: boolean\newline Значення: \textbf{true} \\ \hline
139 & ОЦІНЮВАННЯ ПОСТАЧАННЯ (RA-3(1)) &  & RA-3(1) & Оцінка ризиків здійснюється з урахуванням вимог OWASP Top 10 та чинного законодавства України з кібербезпеки. \\ \hline
140 & СКАНУВАННЯ ВРАЗЛИВОСТЕЙ (RA-5) & Використовувати заходи ПДТР за [Призначення: визначені організацією місця] [Вибір (один або кілька): [Призначення: з визначеною організацією частотою]; [Призначення: за визначеними організацією подіями або показниками]]. & RA-5 & Система передбачає можливість інтеграції зовнішніх сканерів вразливостей, звіти яких автоматично перетворюються на тикети усунення ризиків в системі `itsm`.\vspace{1mm}\newline\noindent\rule{\linewidth}{0.4pt}\vspace{1mm}
Параметр: vulnerability\_\allowbreak{}scanning\newline Тип: boolean\newline Значення: \textbf{true} \\ \hline
141 & СКАНУВАННЯ ВРАЗЛИВОСТЕЙ - ОНОВЛЕННЯ ЗА ЧАСТОТОЮ, ПЕРЕД НОВИМ СКАНУВАННЯМ АБО ПРИ ІДЕНТИФІКАЦІЇ (RA-5(2)) &  & RA-5(2) & Оцінка ризиків здійснюється з урахуванням вимог OWASP Top 10 та чинного законодавства України з кібербезпеки.\vspace{1mm}\newline\noindent\rule{\linewidth}{0.4pt}\vspace{1mm}
Параметр: Визначено час реагування на усунення законних вразливостей відповідно до організаційної оцінення ризику\newline Тип: integer\newline Значення: \textbf{30} \\ \hline
142 & РЕАГУВАННЯ НА РИЗИК (RA-7) & Реагувати на результати оцінювання, моніторингу й аудиту безпеки та приватності. & RA-7 & Оцінка ризиків здійснюється з урахуванням вимог OWASP Top 10 та чинного законодавства України з кібербезпеки. \\ \hline
\multicolumn{5}{|l|}{\textbf{15. ПРИДБАННЯ СИСТЕМ ТА ПОСЛУГ (SA)}} \\ \hline
143 & ПОЛІТИКИ ТА ПРОЦЕДУРИ ПРИДБАННЯ СИСТЕМ ТА ПОСЛУГ (SA-1) & a. Розробити, задокументувати та поширити серед [Призначення: визначеного організацією персоналу або ролей]:\newline 1. 2.\newline b. [Вибір (один або декілька): рівень організації; рівень місії/бізнес-процесу; рівень системи] політики придбання систем і послуг, яка:\newline (a) містить мету, сферу застосування, ролі, обов'язки, відповідальність керівництва, координацію між організаційними підрозділами та систему контролю (complaince);\newline (b) відповідає чинному законодавству, виконавчим наказам, директивам, нормам, політикам, стандартам і рекомендаціям. Процедури, що полегшують впровадження політики та заходів придбання систем і послуг. Призначити [Призначення: визначену організацією посадову особу] для управління політикою та процедурами придбання системи та послуг.\newline c. Переглядати й оновлювати поточні політику та процедури придбання систем та послуг:\newline 1. Поточну політику придбання системи та послуг [Призначення: з визначеною організацією частотою].\newline 2. Поточні процедури придбання системи та послуг [Призначення: з визначеною організацією частотою] та наступні [Призначення: події, визначені організацією]. & SA-1 & Інтеграція нових підсистем (SA) вимагає використання open-source компонентів та безкоштовного ПЗ, сумісного з мікросервісною архітектурою та стандартами ЄСІКС.\vspace{1mm}\newline\noindent\rule{\linewidth}{0.4pt}\vspace{1mm}
Параметр: Визначено персонал або ролі, на які поширюватиметься політика придбання систем і послуг\newline Тип: list\newline Значення: \textbf{admin, security\_\allowbreak{}officer}\vspace{1mm}\newline\noindent\rule{\linewidth}{0.4pt}\vspace{1mm}
Параметр: Визначено персонал або ролі, на які поширюються процедури придбання систем і послуг\newline Тип: list\newline Значення: \textbf{admin, security\_\allowbreak{}officer}\vspace{1mm}\newline\noindent\rule{\linewidth}{0.4pt}\vspace{1mm}
Параметр: Визначено посадову особу, яка керуватиме політикою та процедурами придбання систем і послуг\newline Тип: list\newline Значення: \textbf{admin, security\_\allowbreak{}officer}\vspace{1mm}\newline\noindent\rule{\linewidth}{0.4pt}\vspace{1mm}
Параметр: Визначено періодичність перегляду та оновлення поточної політики придбання систем і послуг\newline Тип: list\newline Значення: \textbf{default\_\allowbreak{}deny\_\allowbreak{}rule, abac\_\allowbreak{}rule\_\allowbreak{}1} \\ 
 &  &  &  & Параметр: Є події, які вимагають перегляду та оновлення поточної політики придбання систем і послуг\newline Тип: list\newline Значення: \textbf{default\_\allowbreak{}deny\_\allowbreak{}rule, abac\_\allowbreak{}rule\_\allowbreak{}1}\vspace{1mm}\newline\noindent\rule{\linewidth}{0.4pt}\vspace{1mm}
Параметр: Визначено частоту, з якою переглядаються та оновлюються поточні процедури придбання систем і послуг\newline Тип: integer\newline Значення: \textbf{30} \\ \hline
144 & БЕЗПЕКА ТА ПРИВАТНІСТЬ ПРИНЦИПІВ ІНЖИНІРИНГУ (SA-8) & Застосовувати [Призначення: визначені організацією принципи інжинірингу безпеки та конфіденційності системи] в специфікації, проєктуванні, розробці, впровадженні та зміні системи й компонентів системи. & SA-8 & Інтеграція нових підсистем (SA) вимагає використання open-source компонентів та безкоштовного ПЗ, сумісного з мікросервісною архітектурою та стандартами ЄСІКС. \\ \hline
145 & ЗОВНІШНІ ПОСЛУГИ ДЛЯ СИСТЕМИ (SA-9) & a. Вимагати, щоб постачальники зовнішніх послуг для системи відповідали вимогам безпеки та приватності в організації та застосовували такі заходи захисту [Призначення: встановлені організацією заходи безпеки та приватності].\newline b. Визначити та задокументувати нагляд організаціїповноваження та обов'язки користувачів щодо зовнішніх послуг для системи.\newline c. Використовувати наступні процеси, методи та техніки для постійного моніторингу дотримання контролю зовнішніми постачальниками послуг: [Призначення: визначені організацією процеси, методи та техніки]. & SA-9 & Інтеграція нових підсистем (SA) вимагає використання open-source компонентів та безкоштовного ПЗ, сумісного з мікросервісною архітектурою та стандартами ЄСІКС.\vspace{1mm}\newline\noindent\rule{\linewidth}{0.4pt}\vspace{1mm}
Параметр: Визначені засоби контролю, які будуть застосовуватися зовнішніми постачальниками системних послуг\newline Тип: string\newline Значення: \textbf{автоматизований засіб моніторингу}\vspace{1mm}\newline\noindent\rule{\linewidth}{0.4pt}\vspace{1mm}
Параметр: Визначені процеси, методи та техніки, що застосовуються для моніторингу дотримання вимог контролю зовнішніми постачальниками послуг; SA-09a.[01] дотримуються постачальники зовнішніх системних послуг вимог організаційної безпеки; SA-09a.[02] дотримуються постачальники зовнішніх системних послуг вимог приватності організації; SA-09a.[03] використовують постачальники зовнішніх системних послуг елементи керування; SA-09b.[01] визначено та задокументовано організаційний нагляд за зовнішніми системними послугами; SA-09b.[02] визначені та задокументовані ролі та обов'язки користувачів щодо зовнішніх системних сервісів; SA-09c. визначені та задокументовані ролі та обов'язки користувачів щодо зовнішніх системних послуг\newline Тип: list\newline Значення: \textbf{admin, security\_\allowbreak{}officer} \\ \hline
146 & КОМПОНЕНТИ СИСТЕМИ, ЩО НЕ ПІДТРИМУЮТЬСЯ (SA-22) & a. Замінювати компоненти системи, якщо підтримка компонентів більше не доступна розробнику, постачальнику або виробнику.\newline b. Надавати такі варіанти альтернативних джерел для подальшої підтримки непідтримуваних компонентів [Вибір (один або більше): внутрішня підтримка; [Призначення: підтримка, визначена організацією від зовнішніх постачальників]]. & SA-22 & Інтеграція нових підсистем (SA) вимагає використання open-source компонентів та безкоштовного ПЗ, сумісного з мікросервісною архітектурою та стандартами ЄСІКС. \\ \hline
\multicolumn{5}{|l|}{\textbf{16. ЗАХИСТ СИСТЕМ ТА КОМУНІКАЦІЙ (SC)}} \\ \hline
147 & ПОЛІТИКА ТА ПРОЦЕДУРИ ЗАХИСТУ СИСТЕМИ ТА КОМУНІКАЦІЙ (SC-1) & a. Розробити, задокументувати та поширити серед [Призначення: визначеного організацією персоналу або посадових осіб]:\newline 1. 2. Політику захисту системи та комунікацій, яка:\newline (a) містить мету, сферу застосування, ролі, обов'язки, відповідальність керівництва, координацію між організаційними підрозділами та систему контролю відповідності (complaince);\newline (b) відповідає чинному законодавству, виконавчим наказам, директивам, нормам, політикам, стандартам та керівним принципам. Процедури для сприяння впровадженню політики в області захисту систем і комунікацій, а також пов'язаних з ними систем і засобів захисту зв'язку.\newline b. Призначити [Призначення: визначена організацією посадову особу] для управління політикою та процедурами захисту системи та комунікацій.\newline c. Переглядати та оновлювати:\newline 1. поточну політику захисту системи та комунікацій [Призначення: визначена організацією частота];\newline 2. поточні процедури захисту системи та комунікацій [Призначення: визначена організацією частота]. & SC-1 & Комунікації захищені TLS версії не нижче 1.3 на порту 443. Заборонено відкривати інші порти назовні. Взаємодія сервісів використовує брокери повідомлень та REST/RPC.\vspace{1mm}\newline\noindent\rule{\linewidth}{0.4pt}\vspace{1mm}
Параметр: Визначено персонал або ролі, до яких має бути доведена політика захисту системи та комунікацій\newline Тип: list\newline Значення: \textbf{admin, security\_\allowbreak{}officer}\vspace{1mm}\newline\noindent\rule{\linewidth}{0.4pt}\vspace{1mm}
Параметр: Визначено персонал або ролі, на які поширюються процедури захисту системи та комунікацій\newline Тип: list\newline Значення: \textbf{admin, security\_\allowbreak{}officer}\vspace{1mm}\newline\noindent\rule{\linewidth}{0.4pt}\vspace{1mm}
Параметр: Визначено посадову особу, яка керуватиме політикою та процедурами захисту системи та комунікацій\newline Тип: list\newline Значення: \textbf{admin, security\_\allowbreak{}officer}\vspace{1mm}\newline\noindent\rule{\linewidth}{0.4pt}\vspace{1mm}
Параметр: Визначена періодичність перегляду та оновлення поточної політики захисту системи та комунікацій\newline Тип: list\newline Значення: \textbf{default\_\allowbreak{}deny\_\allowbreak{}rule, abac\_\allowbreak{}rule\_\allowbreak{}1} \\ 
 &  &  &  & Параметр: Є події, які вимагають перегляду та оновлення поточної політики захисту системи та комунікацій\newline Тип: list\newline Значення: \textbf{default\_\allowbreak{}deny\_\allowbreak{}rule, abac\_\allowbreak{}rule\_\allowbreak{}1}\vspace{1mm}\newline\noindent\rule{\linewidth}{0.4pt}\vspace{1mm}
Параметр: Визначена періодичність перегляду та оновлення поточних процедур захисту системи та засобів зв'язку\newline Тип: string\newline Значення: \textbf{щорічно} \\ \hline
148 & ІНФОРМАЦІЯ В ЗАГАЛЬНИХ СИСТЕМНИХ РЕСУРСАХ (SC-4) & Запобігати несанкціонованій та ненавмисній передачі інформації через спільні системні ресурси. & SC-4 & Erlang VM не використовує спільну пам'ять (shared memory) між процесами сесій N2O. Пам'ять звільняється автоматично (Garbage Collection) після обробки запиту, запобігаючи витоку залишкової інформації між різними користувачами.\vspace{1mm}\newline\noindent\rule{\linewidth}{0.4pt}\vspace{1mm}
Параметр: shared\_\allowbreak{}resource\_\allowbreak{}protection\newline Тип: boolean\newline Значення: \textbf{true} \\ \hline
149 & ДОСТУПНІСТЬ РЕСУРСІВ (SC-7) & a. Контролювати та управляти зв'язком на зовнішньому периметрі системи та на ключових внутрішніх периметрах всередині системи.\newline b. Реалізувати підмережі для загальнодоступних компонентів системи, які є [Вибір: фізично; логічно] відділені від внутрішніх мереж організації.\newline c. Підключатися до зовнішніх мереж або систем тільки через керовані інтерфейси, що складаються з пристроїв захисту периметру, і розташованих відповідно до архітектури безпеки та приватності організації. & SC-7 & Всі зовнішні підключення приймаються виключно через єдиний порт (керований інтерфейс), що обробляється веб-сервером (Cowboy). Доступ до внутрішніх API і вузлів кластера блокується правилами мережевого екранування за замовчуванням (Default Deny).\vspace{1mm}\newline\noindent\rule{\linewidth}{0.4pt}\vspace{1mm}
Параметр: boundary\_\allowbreak{}protection\newline Тип: boolean\newline Значення: \textbf{true} \\ \hline
150 & ЗАХИСТ ПЕРИМЕТРА - ВІДМОВА ЗА ЗАМОВЧУВАННЯМ - ДОЗВІЛ ЗА ВИНЯТКОМ (SC-7(5)) &  & SC-7(5) & Комунікації захищені TLS версії не нижче 1.3 на порту 443. Заборонено відкривати інші порти назовні. Взаємодія сервісів використовує брокери повідомлень та REST/RPC. \\ \hline
151 & КОНФІДЕНЦІЙНІСТЬ ТА ЦІЛІСНІСТЬ ПЕРЕДАЧІ (SC-8) & Забезпечити [Вибір (один або кілька): конфіденційність; цілісність] інформації, що передається. & SC-8 & Передача інформації між клієнтом і сервером здійснюється виключно через зашифровані канали (WSS/TLS), що унеможливлює перехоплення (Man-in-the-Middle) або несанкціоновану модифікацію даних на льоту.\vspace{1mm}\newline\noindent\rule{\linewidth}{0.4pt}\vspace{1mm}
Параметр: transmission\_\allowbreak{}confidentiality\newline Тип: boolean\newline Значення: \textbf{true} \\ \hline
152 & КОНФІДЕНЦІЙНІСТЬ ТА КРИПТОГРАФІЧНИЙ ЗАХИСТ (SC-8(1)) &  & SC-8(1) & Комунікації захищені TLS версії не нижче 1.3 на порту 443. Заборонено відкривати інші порти назовні. Взаємодія сервісів використовує брокери повідомлень та REST/RPC. \\ \hline
153 & ВІДКЛЮЧЕННЯ МЕРЕЖІ (SC-10) & Завершити з'єднання з мережею, яке пов'язане із сеансом зв'язку в кінці сеансу або після [Призначення: визначений організацією період часу] бездіяльності. & SC-10 & Протокол N2O автоматично розриває мережеве з'єднання після завершення комунікаційного сеансу або перевищення часу очікування (Inactivity Timeout), звільняючи ресурси та сесію.\vspace{1mm}\newline\noindent\rule{\linewidth}{0.4pt}\vspace{1mm}
Параметр: network\_\allowbreak{}disconnect\_\allowbreak{}timeout\newline Тип: integer\newline Значення: \textbf{3600} \\ \hline
154 & ВСТАНОВЛЕННЯ КЛЮЧАМИ (SC-12) & Встановити та управляти криптографічними ключами для криптографічних засобів, які використовуються в системі відповідно до [Призначення: визначені організацією вимоги до генерації, поширення, зберігання, доступу та знищення ключів]. & SC-12 & Керування ключами та сертифікатами для встановлення захищених з'єднань делегується бібліотеці `synrc/ca` та інтегрованим механізмам TLS ОС, що відповідають чинним криптографічним стандартам.\vspace{1mm}\newline\noindent\rule{\linewidth}{0.4pt}\vspace{1mm}
Параметр: cryptographic\_\allowbreak{}key\_\allowbreak{}management\newline Тип: boolean\newline Значення: \textbf{true} \\ \hline
155 & КРИПТОГРАФІЧНИЙ ЗАХИСТ (SC-13) & a. Визначити [Призначення: використання криптографічних засобів, визначених організацією];\newline b. Впровадити [Завдання: визначені організацією види криптографії для кожного визначеного криптографічного використання]. & SC-13 & Система підтримує використання схвалених криптографічних алгоритмів (AES-GCM, SHA-256/512, ДСТУ 4145-2002 через відповідні модулі), які забезпечують необхідний рівень стійкості.\vspace{1mm}\newline\noindent\rule{\linewidth}{0.4pt}\vspace{1mm}
Параметр: cryptographic\_\allowbreak{}protection\newline Тип: boolean\newline Значення: \textbf{true} \\ \hline
156 & СПІЛЬНІ ОБЧИСЛЮВАЛЬНІ ПРИСТРОЇ ТА ЗАСТОСУНКИ (SC-15) & a. Заборонити віддалену активацію спільних обчислювальних пристроїв (хмар) та застосунків з такими виключеннями: [Призначення: визначені організацією виключення, у яких дозволена віддалена активація].\newline b. Надати явну вказівку щодо використання користувачами фізично присутніми пристроями. & SC-15 & Комунікації захищені TLS версії не нижче 1.3 на порту 443. Заборонено відкривати інші порти назовні. Взаємодія сервісів використовує брокери повідомлень та REST/RPC. \\ \hline
157 & СЕРТИФІКАТИ ІНФРАСТРУКТУРИ ВІДКРИТИХ КЛЮЧІВ (SC-17) & a. Випускати сертифікати відкритого ключа відповідно до [Призначення: визначеної організацією політики сертифікації];\newline b. Отримувати сертифікати відкритого ключа від затвердженого постачальника послуг. & SC-17 & Комунікації захищені TLS версії не нижче 1.3 на порту 443. Заборонено відкривати інші порти назовні. Взаємодія сервісів використовує брокери повідомлень та REST/RPC. \\ \hline
158 & МОБІЛЬНИЙ КОД (SC-18) & a. Визначати прийнятні та неприйнятні мобільні коди та технології мобільних кодів.\newline b. Проводити авторизацію, відстежувати та контролювати використання мобільного коду всередині системи. & SC-18 &  \\ \hline
159 & АВТЕНТИФІКАЦІЯ СЕСІЇ (SC-23) & Забезпечити автентифікацію сеансів зв'язку. & SC-23 & Автентичність WebSocket-сесій N2O захищається токенами з криптографічним підписом (HMAC). Це гарантує, що повідомлення в рамках сесії не можуть бути підроблені або інжектовані сторонньою особою.\vspace{1mm}\newline\noindent\rule{\linewidth}{0.4pt}\vspace{1mm}
Параметр: session\_\allowbreak{}authenticity\newline Тип: boolean\newline Значення: \textbf{true} \\ \hline
160 & ЗАХИСТ ІНФОРМАЦІЇ В СТАНІ СПОКОЮ (SC-28) & Забезпечити [Вибір (один або кілька): конфіденційність; цілісність] [Призначення: визначена організацією інформація] в стані спокою. & SC-28 & Комунікації захищені TLS версії не нижче 1.3 на порту 443. Заборонено відкривати інші порти назовні. Взаємодія сервісів використовує брокери повідомлень та REST/RPC. \\ \hline
161 & ЗАХИСТ ІНФОРМАЦІЇ В СТАНІ СПОКОЮ | КРИПТОГРАФІЧНИЙ ЗАХИСТ &  & SC-28(1) & Комунікації захищені TLS версії не нижче 1.3 на порту 443. Заборонено відкривати інші порти назовні. Взаємодія сервісів використовує брокери повідомлень та REST/RPC.\vspace{1mm}\newline\noindent\rule{\linewidth}{0.4pt}\vspace{1mm}
Параметр: Реалізовані криптографічні механізми для запобігання несанкціонованому розкриттю інформації, що знаходиться в стані спокою на системних компонентах або носіях\newline Тип: string\newline Значення: \textbf{AES-256-GCM}\vspace{1mm}\newline\noindent\rule{\linewidth}{0.4pt}\vspace{1mm}
Параметр: Реалізовані криптографічні механізми для запобігання несанкціонованій модифікації інформації, що знаходиться в стані спокою на системних компонентах або носіях\newline Тип: string\newline Значення: \textbf{AES-256-GCM} \\ \hline
162 & ЕКРАНОВАНІ КАМЕРИ (SC-44) & Впровадити екрановані камери в [Призначення: визначену організацією систему, компонент системи або місце розташування]. & SC-44 & Комунікації захищені TLS версії не нижче 1.3 на порту 443. Заборонено відкривати інші порти назовні. Взаємодія сервісів використовує брокери повідомлень та REST/RPC. \\ \hline
163 & ПРИМУСОВЕ АПАРАТНЕ ЗАБЕЗПЕЧЕННЯ ВИКОНАННЯ (SC-49) & Впровадити механізми апаратного поділу та застосування політики між [Призначення: домени безпеки, визначені організацією]. & SC-49 & Комунікації захищені TLS версії не нижче 1.3 на порту 443. Заборонено відкривати інші порти назовні. Взаємодія сервісів використовує брокери повідомлень та REST/RPC.\vspace{1mm}\newline\noindent\rule{\linewidth}{0.4pt}\vspace{1mm}
Параметр: Впроваджено механізми апаратного розділення та застосування політик між доменами безпеки\newline Тип: list\newline Значення: \textbf{default\_\allowbreak{}deny\_\allowbreak{}rule, abac\_\allowbreak{}rule\_\allowbreak{}1}\vspace{1mm}\newline\noindent\rule{\linewidth}{0.4pt}\vspace{1mm}
Параметр: Визначені домени безпеки, які потребують апаратного розділення та механізмів забезпечення дотримання політики\newline Тип: list\newline Значення: \textbf{default\_\allowbreak{}deny\_\allowbreak{}rule, abac\_\allowbreak{}rule\_\allowbreak{}1} \\ \hline
164 & АПАРАТНИЙ ЗАХИСТ (SC-51) & a. Перевіряти правильність роботи [Призначення: визначені організацією функції безпеки та приватності].\newline b. Виконувати перевірку [Вибір (один або кілька): [Призначення: визначені організацією системні перехідні стани]; за командою користувача з відповідними повноваженнями; [Призначення: визначена організацією частота]].\newline c. Повідомляти [Призначення: визначені організацією персонал або посадові особи] про невдалі перевірки безпеки та приватності.\newline d. [Вибір (один або кілька): Вимкнути систему; Перезапустити систему; [Призначення: визначені організацією альтернативні дії]], коли виявляються аномалії. & SC-51 & Комунікації захищені TLS версії не нижче 1.3 на порту 443. Заборонено відкривати інші порти назовні. Взаємодія сервісів використовує брокери повідомлень та REST/RPC.\vspace{1mm}\newline\noindent\rule{\linewidth}{0.4pt}\vspace{1mm}
Параметр: Визначені уповноважені особи, які повинні виконувати процедури вимкнення та повторного ввімкнення апаратного захисту від запису; SC-51a. використовується апаратний захист від запису для компонентів мікропрограми системи; SC-51b.[01] впроваджено спеціальні процедури для уповноважених осіб для ручного вимкнення апаратного захисту від запису для модифікацій мікропрограми; SC-51b.[02] реалізовано спеціальні процедури для уповноважених осіб для повторного увімкнення захисту від запису перед поверненням до робочого режиму\newline Тип: list\newline Значення: \textbf{admin, security\_\allowbreak{}officer} \\ \hline
\multicolumn{5}{|l|}{\textbf{17. ЦІЛІСНІСТЬ СИСТЕМИ ТА ІНФОРМАЦІЇ (SI)}} \\ \hline
165 & ПОЛІТИКА І ПРОЦЕДУРИ ЦІЛІСНОСТІ ІНФОРМАЦІЇ (SI-1) &  & SI-1 & Цілісність даних забезпечується захистом від SQL та XSS ін'єкцій, екрануванням спецсимволів, валідацією JSON Schema та імутабельністю історичних даних (HL7, BPMN).\vspace{1mm}\newline\noindent\rule{\linewidth}{0.4pt}\vspace{1mm}
Параметр: Визначено персонал або ролі, до яких має бути доведена політика цілісності системи та інформації\newline Тип: list\newline Значення: \textbf{admin, security\_\allowbreak{}officer}\vspace{1mm}\newline\noindent\rule{\linewidth}{0.4pt}\vspace{1mm}
Параметр: Визначено персонал або ролі, на які поширюються процедури цілісності системи та інформації \newline Тип: list\newline Значення: \textbf{admin, security\_\allowbreak{}officer}\vspace{1mm}\newline\noindent\rule{\linewidth}{0.4pt}\vspace{1mm}
Параметр: Визначено посадову особу, відповідальну за управління системою та політикою і процедурами цілісності інформації\newline Тип: list\newline Значення: \textbf{admin, security\_\allowbreak{}officer}\vspace{1mm}\newline\noindent\rule{\linewidth}{0.4pt}\vspace{1mm}
Параметр: Визначено періодичність перегляду та оновлення поточної політики цілісності системи та інформації\newline Тип: list\newline Значення: \textbf{default\_\allowbreak{}deny\_\allowbreak{}rule, abac\_\allowbreak{}rule\_\allowbreak{}1} \\ 
 &  &  &  & Параметр: Є події, які вимагають перегляду та оновлення поточної політики цілісності системи та інформації\newline Тип: list\newline Значення: \textbf{default\_\allowbreak{}deny\_\allowbreak{}rule, abac\_\allowbreak{}rule\_\allowbreak{}1}\vspace{1mm}\newline\noindent\rule{\linewidth}{0.4pt}\vspace{1mm}
Параметр: Визначено частоту, з якою переглядаються та оновлюються поточні цілісності системи та інформації\newline Тип: integer\newline Значення: \textbf{30} \\ \hline
166 & ВИПРАВЛЕННЯ ДЕФЕКТІВ (SI-2) &  & SI-2 & Цілісність даних забезпечується захистом від SQL та XSS ін'єкцій, екрануванням спецсимволів, валідацією JSON Schema та імутабельністю історичних даних (HL7, BPMN).\vspace{1mm}\newline\noindent\rule{\linewidth}{0.4pt}\vspace{1mm}
Параметр: Визначено період часу, протягом якого необхідно встановити оновлення програмного забезпечення, пов'язані з безпекою, після виходу оновлень; SI-02a.[01] виявлено недоліки системи; SI-02a.[02] повідомляється про недоліки системи; SI-02a.[03] виправлені недоліки системи; SI-02b.[01] перевіряються оновлення програмного забезпечення, пов'язані з усуненням недоліків, на ефективність перед встановленням; SI-02b.[02] перевіряються оновлення програмного забезпечення, пов'язані з виправленням дефектів, на наявність потенційних побічних ефектів перед встановленням; SI-02b.[03] перевіряються оновлення прошивки, пов'язані з усуненням недоліків, на ефективність перед встановленням; SI-02b.[04] перевіряються оновлення прошивки, пов'язані з усуненням недоліків, на наявність потенційних побічних ефектів перед встановленням; SI-02c.[01] встановлено оновлення програмного забезпечення, що стосуються безпеки, протягом часовий проміжок з моменту випуску оновлень; SI-02c.[02] встановлено оновлення мікропрограми, що стосуються безпеки, протягом часового періоду з моменту випуску оновлень; SI-02d. включено відновлення порушених прав у процес управління організаційною конфігурацією\newline Тип: integer\newline Значення: \textbf{30} \\ \hline
167 & ЗАХИСТ ВІД ШКІДЛИВОГО КОДУ (SI-3) &  & SI-3 & Цілісність даних забезпечується захистом від SQL та XSS ін'єкцій, екрануванням спецсимволів, валідацією JSON Schema та імутабельністю історичних даних (HL7, BPMN).\vspace{1mm}\newline\noindent\rule{\linewidth}{0.4pt}\vspace{1mm}
Параметр: Визначено частоту, з якою механізми шкідливого коду виконують сканування\newline Тип: integer\newline Значення: \textbf{30}\vspace{1mm}\newline\noindent\rule{\linewidth}{0.4pt}\vspace{1mm}
Параметр: Визначено дії, яких слід вжити у відповідь на виявлення шкідливого коду (якщо вибрано)\newline Тип: list\newline Значення: \textbf{login, logout, failed\_\allowbreak{}attempt} \\ \hline
168 & МОНІТОРИНГ СИСТЕМИ (SI-4) &  & SI-4 & Цілісність даних забезпечується захистом від SQL та XSS ін'єкцій, екрануванням спецсимволів, валідацією JSON Schema та імутабельністю історичних даних (HL7, BPMN).\vspace{1mm}\newline\noindent\rule{\linewidth}{0.4pt}\vspace{1mm}
Параметр: Визначені методи та способи, що використовуються для виявлення несанкціонованого використання системи\newline Тип: list\newline Значення: \textbf{admin, security\_\allowbreak{}officer}\vspace{1mm}\newline\noindent\rule{\linewidth}{0.4pt}\vspace{1mm}
Параметр: Визначена інформація про моніторинг системи, яка повинна надаватися персоналу або ролям\newline Тип: list\newline Значення: \textbf{admin, security\_\allowbreak{}officer}\vspace{1mm}\newline\noindent\rule{\linewidth}{0.4pt}\vspace{1mm}
Параметр: Визначено персонал або ролі, яким має надаватися інформація про моніторинг системи\newline Тип: list\newline Значення: \textbf{admin, security\_\allowbreak{}officer}\vspace{1mm}\newline\noindent\rule{\linewidth}{0.4pt}\vspace{1mm}
Параметр: Вибрано одне або декілька з наступних ЗНАЧЕНЬ ПАРАМЕТРІВ: \{за потребою; частота\}\newline Тип: string\newline Значення: \textbf{щорічно} \\ \hline
169 & МОНІТОРИНГ СИСТЕМИ КОМУНІКАЦІЙ (SI-4(4)) &  & SI-4(4) & Цілісність даних забезпечується захистом від SQL та XSS ін'єкцій, екрануванням спецсимволів, валідацією JSON Schema та імутабельністю історичних даних (HL7, BPMN).\vspace{1mm}\newline\noindent\rule{\linewidth}{0.4pt}\vspace{1mm}
Параметр: Визначені критерії незвичної або несанкціонованої діяльності або умови для вхідного трафіку зв'язку\newline Тип: list\newline Значення: \textbf{}\vspace{1mm}\newline\noindent\rule{\linewidth}{0.4pt}\vspace{1mm}
Параметр: Визначені критерії незвичайної або несанкціонованої діяльності або умови для вихідного трафіку зв'язку\newline Тип: list\newline Значення: \textbf{}\vspace{1mm}\newline\noindent\rule{\linewidth}{0.4pt}\vspace{1mm}
Параметр: Здійснюється моніторинг вхідного комунікаційного трафіку частота на предмет незвичних або несанкціонованих дій або умов\newline Тип: string\newline Значення: \textbf{щорічно}\vspace{1mm}\newline\noindent\rule{\linewidth}{0.4pt}\vspace{1mm}
Параметр: Контролюється вихідний трафік зв'язку частота на предмет незвичних або несанкціонованих дій або умов. вихідного виявлення\newline Тип: string\newline Значення: \textbf{щорічно} \\ \hline
170 & ПОПЕРЕДЖЕННЯ, РЕКОМЕНДАЦІЇ ТА ДИРЕКТИВИ З БЕЗПЕКИ (SI-5) &  & SI-5 & Цілісність даних забезпечується захистом від SQL та XSS ін'єкцій, екрануванням спецсимволів, валідацією JSON Schema та імутабельністю історичних даних (HL7, BPMN).\vspace{1mm}\newline\noindent\rule{\linewidth}{0.4pt}\vspace{1mm}
Параметр: Вибрано одне або декілька з наступних ЗНАЧЕНЬ ПАРАМЕТРІВ: \{персонал або ролі; елементи; зовнішні організації\}\newline Тип: list\newline Значення: \textbf{admin, security\_\allowbreak{}officer}\vspace{1mm}\newline\noindent\rule{\linewidth}{0.4pt}\vspace{1mm}
Параметр: Визначено персонал або ролі, до яких мають бути доведені попередження, поради та директиви з безпеки (якщо визначено)\newline Тип: list\newline Значення: \textbf{admin, security\_\allowbreak{}officer} \\ \hline
171 & ЦІЛІСНІСТЬ ПРОГРАМНОГО ЗАБЕЗПЕЧЕННЯ, ПРОГРАМНОГО ЗАБЕЗПЕЧЕННЯ ТА ІНФОРМАЦІЇ (SI-7) & a. Впровадити механізми захисту від спаму в точках входу та виходу системи, щоб виявляти та протидіяти небажаним повідомленням.\newline b. Оновлювати механізми захисту від спаму, коли доступні нові механізми відповідно до організаційної політики та процедур управління конфігурацією. & SI-7 & Цілісність даних забезпечується захистом від SQL та XSS ін'єкцій, екрануванням спецсимволів, валідацією JSON Schema та імутабельністю історичних даних (HL7, BPMN).\vspace{1mm}\newline\noindent\rule{\linewidth}{0.4pt}\vspace{1mm}
Параметр: Визначені дії, яких слід вжити при виявленні несанкціонованих змін у програмному забезпеченні\newline Тип: list\newline Значення: \textbf{login, logout, failed\_\allowbreak{}attempt}\vspace{1mm}\newline\noindent\rule{\linewidth}{0.4pt}\vspace{1mm}
Параметр: Визначені дії, яких слід вжити несанкціонованих змін у прошивці; при виявленні SI-07\_\allowbreak{}ODP[06] визначені дії, яких слід вжити несанкціонованих змін до інформації; при виявленні SI-07a.[01] використовуються засоби перевірки цілісності для виявлення несанкціонованих змін у програмному забезпеченні; SI-07a.[02] використовуються засоби перевірки цілісності для виявлення несанкціонованих змін у мікропрограмі; SI-07a.[03] використовуються засоби перевірки цілісності для виявлення несанкціонованих змін до інформації; SI-07b.[01] виконуються дії при виявленні несанкціонованих змін у програмному забезпеченні; SI-07b.[02] виконуються дії несанкціонованих змін у прошивці; при виявленні SI-07b.[03] виконуються дії несанкціонованих змін в інформації. при виявленні\newline Тип: list\newline Значення: \textbf{login, logout, failed\_\allowbreak{}attempt} \\ \hline
172 & УПРАВЛІННЯ ТА ЗБЕРЕЖЕННЯ ІНФОРМАЦІЇ (SI-12) & Управляти та зберігати інформацію всередині системи та виводити інформацію із системи відповідно до чинного законодавства, виконавчих наказів, директив, правил, політик, стандартів, керівних принципів та експлуатаційних вимог. & SI-12 & Цілісність даних забезпечується захистом від SQL та XSS ін'єкцій, екрануванням спецсимволів, валідацією JSON Schema та імутабельністю історичних даних (HL7, BPMN).\vspace{1mm}\newline\noindent\rule{\linewidth}{0.4pt}\vspace{1mm}
Параметр: Здійснюється управління інформацією в системі відповідно до чинних законів, наказів, директив, положень, політик, стандартів, інструкцій та експлуатаційних вимог\newline Тип: list\newline Значення: \textbf{default\_\allowbreak{}deny\_\allowbreak{}rule, abac\_\allowbreak{}rule\_\allowbreak{}1}\vspace{1mm}\newline\noindent\rule{\linewidth}{0.4pt}\vspace{1mm}
Параметр: Зберігається інформація в системі відповідно до чинних законів, указів Президента, директив, положень, політик, стандартів, інструкцій та експлуатаційних вимог\newline Тип: list\newline Значення: \textbf{default\_\allowbreak{}deny\_\allowbreak{}rule, abac\_\allowbreak{}rule\_\allowbreak{}1}\vspace{1mm}\newline\noindent\rule{\linewidth}{0.4pt}\vspace{1mm}
Параметр: Управління інформацією, що виводиться з системи, здійснюється відповідно до чинних законів, указів Президента, директив, положень, політик, стандартів, інструкцій та операційних вимог\newline Тип: list\newline Значення: \textbf{default\_\allowbreak{}deny\_\allowbreak{}rule, abac\_\allowbreak{}rule\_\allowbreak{}1}\vspace{1mm}\newline\noindent\rule{\linewidth}{0.4pt}\vspace{1mm}
Параметр: Зберігається інформація, що виводиться з системи, відповідно до чинних законів, наказів, директив, положень, політик, стандартів, інструкцій та експлуатаційних вимог\newline Тип: list\newline Значення: \textbf{default\_\allowbreak{}deny\_\allowbreak{}rule, abac\_\allowbreak{}rule\_\allowbreak{}1} \\ \hline
\multicolumn{5}{|l|}{\textbf{18. УПРАВЛІННЯ РИЗИКАМИ В ЛАНЦЮГУ ПОСТАЧАННЯ (SR)}} \\ \hline
173 & ПОЛІТИКА ТА ПРОЦЕДУРИ УПРАВЛІННЯ РИЗИКАМИ ЛАНЦЮГА ПОСТАЧАННЯ (SR-1) & a. Розробіть, задокументуйте та поширте [Призначення: персонал або ролі, визначені організацією]:\newline 1. [Вибір (один або декілька): Рівень організації; Рівень місії/бізнес-процесу; рівень системи] політика управління ризиками ланцюга постачання, яка: a) Розглядає мету, сферу діяльності, ролі, відповідальність, зобов'язання керівництва, координацію між організаційними підрозділами та відповідність; b) Відповідає чинним законам, виконавчим наказам, директивам, положенням, політикам, стандартам і вказівкам;\newline 2. Процедури для сприяння впровадженню політики управління ризиками ланцюга постачання та відповідних засобів контролю управління ризиками ланцюга постачання;\newline b. Призначити [Призначення: посадова особа, визначена організацією] для управління розробкою, документуванням і розповсюдженням політики та процедур управління ризиками ланцюга постачання;\newline c. Перегляньте та оновіть поточне управління ризиками ланцюга постачання:\newline 1. Політика [Призначення: частота, визначена організацією] та наступне [Призначення: події, визначені організацією];\newline 2. Процедури [Призначення: частота, визначена організацією] та наступні [Призначення: події, визначені організацією]. & SR-1 & Управління ризиками ланцюга постачання базується на перевірці бібліотек через exploit-db та обов'язковому оновленні вразливостей.\vspace{1mm}\newline\noindent\rule{\linewidth}{0.4pt}\vspace{1mm}
Параметр: Визначено персонал або ролі, на які поширюється політика управління ризиками ланцюга постачання\newline Тип: list\newline Значення: \textbf{admin, security\_\allowbreak{}officer}\vspace{1mm}\newline\noindent\rule{\linewidth}{0.4pt}\vspace{1mm}
Параметр: Визначено персонал або ролі, на які поширюються процедури управління ризиками ланцюга постачання\newline Тип: list\newline Значення: \textbf{admin, security\_\allowbreak{}officer}\vspace{1mm}\newline\noindent\rule{\linewidth}{0.4pt}\vspace{1mm}
Параметр: Визначена посадова особа, відповідальна за розробку, документування та розповсюдження політики та процедур управління ризиками ланцюга постачання\newline Тип: list\newline Значення: \textbf{admin, security\_\allowbreak{}officer}\vspace{1mm}\newline\noindent\rule{\linewidth}{0.4pt}\vspace{1mm}
Параметр: Визначена періодичність перегляду та оновлення поточної політики управління ризиками ланцюга постачання\newline Тип: list\newline Значення: \textbf{default\_\allowbreak{}deny\_\allowbreak{}rule, abac\_\allowbreak{}rule\_\allowbreak{}1} \\ 
 &  &  &  & Параметр: Є події, які вимагають перегляду та оновлення поточної політики управління ризиками ланцюга постачання\newline Тип: list\newline Значення: \textbf{default\_\allowbreak{}deny\_\allowbreak{}rule, abac\_\allowbreak{}rule\_\allowbreak{}1}\vspace{1mm}\newline\noindent\rule{\linewidth}{0.4pt}\vspace{1mm}
Параметр: Визначена періодичність перегляду та оновлення поточної процедури управління ризиками ланцюга постачання\newline Тип: string\newline Значення: \textbf{щорічно} \\ \hline
174 & ПЛАН УПРАВЛІННЯ РИЗИКАМИ ЛАНЦЮГА ПОСТАЧАННЯ (SR-2) & a. Розробіть план управління ризиками ланцюга постачання, пов'язаними з дослідженнями та розробкою, проектуванням, виробництвом, придбанням, доставкою, інтеграцією, експлуатацією та обслуговуванням, а також утилізацією таких систем, компонентів системи або послуг для системи: [Призначення: системи, визначені організацією, системні компоненти або системні служби];\newline b. Перегляньте та оновіть план управління ризиками ланцюга постачання [Призначення: частота, визначена організацією] або за потреби для усунення загроз;\newline c. Захистіть план управління ризиками ланцюга постачання від несанкціонованого розголошення та модифікації. & SR-2 & План управління базується на принципі відмови від стороннього коду. Замість інтеграції сотень неперевірених бібліотек, команда розробляє всі необхідні компоненти (веб-сервер N2O, KVS, CA) самостійно в єдиному монорепозиторії.\vspace{1mm}\newline\noindent\rule{\linewidth}{0.4pt}\vspace{1mm}
Параметр: zero\_\allowbreak{}dependency\_\allowbreak{}policy\newline Тип: boolean\newline Значення: \textbf{true} \\ \hline
175 & КОНТРОЛЬ ЛАНЦЮГА ПОСТАЧАННЯ І ПРОЦЕСІВ (SR-3) & a. Встановлення процесу або процесів для виявлення та усунення слабких місць або недоліків в елементах і процесах ланцюга постачання [Призначення: визначена організацією система або компонент системи] у координації з [Завдання: персонал ланцюга постачання, визначений організацією];\newline b. Використовуйте такі заходи захисту, щоб захистити систему, компонент системи або системну службу від ризиків ланцюга постачання та обмежити шкоду чи наслідки від подій, пов'язаних із ланцюгом постачання: [Призначення: заходи захисту ланцюга постачання, визначені організацією];\newline c. Задокументуйте обрані та впроваджені процеси та заходи захисту ланцюгом постачання у [Вибір: плани безпеки та приватності; план управління ризиками ланцюга постачання; [Призначення: документ, визначений організацією]]. & SR-3 & Заборонено використання автоматичних менеджерів пакетів для завантаження коду під час збірки. Весь вихідний код системи та модулів міститься безпосередньо в репозиторії проекту, що гарантує 100\% контроль над кожним рядком коду, який компілюється.\vspace{1mm}\newline\noindent\rule{\linewidth}{0.4pt}\vspace{1mm}
Параметр: strict\_\allowbreak{}code\_\allowbreak{}vendoring\newline Тип: boolean\newline Значення: \textbf{true} \\ \hline
176 & СТРАТЕГІЇ ПРИДБАННЯ, ІНСТРУМЕНТИ І МЕТОДИ (SR-5) & Використовуйте наступні стратегії придбання, контрактні інструменти та методи закупівель, щоб захистити від ризиків ланцюга постачання, визначити та пом'якшити їх: [Призначення: визначені організацією стратегії придбання, контрактні інструменти та методи закупівель]. & SR-5 & Управління ризиками ланцюга постачання базується на перевірці бібліотек через exploit-db та обов'язковому оновленні вразливостей. \\ \hline
177 & ОЦІНКА ПОСТАЧАЛЬНИКІВ (SR-6) & Оцініть і перегляньте ризики ланцюга постачання, пов'язані з постачальниками або підрядниками, системою, системним компонентом або системною послугою, яку вони надають [Призначення: частота, визначена організацією]. & SR-6 & Єдиним зовнішнім «постачальником» є розробник платформи Erlang/OTP (компанія Ericsson та спільнота). Оновлення версій OTP проходять жорстку перевірку перед їх застосуванням на продуктивних серверах.\vspace{1mm}\newline\noindent\rule{\linewidth}{0.4pt}\vspace{1mm}
Параметр: erlang\_\allowbreak{}otp\_\allowbreak{}supplier\_\allowbreak{}review\newline Тип: boolean\newline Значення: \textbf{true} \\ \hline


\end{longtable}

\end{document}
